% ----------------------------------------------------------------------
%              Latex PhD template for the University of Deusto
% ----------------------------------------------------------------------


%: Style file for Latex
% Most style definitions are in the external file PhDthesisPSnPDF.
% In this template package, it can be found in ./Latex/Classes/
\documentclass[twoside,12pt]{Latex/Classes/PhDthesisPSnPDF}

% My own packages
%\usepackage{pdflscape}
%\usepackage{longtable}
\usepackage{algpseudocode}
\usepackage{minted}
\usepackage[printonlyused]{acronym}
\usepackage{enumerate}
\usepackage{footnote} % footnotes inside tables using savenotes
\usepackage[inline]{enumitem} % inline lists
% For inline cites (section 6.2
\usepackage{bibentry} % http://gking.harvard.edu/files/bibentry2.pdf
%\nobibliography* throws an error
\nobibliography{9_backmatter/ref_self} % http://stefaanlippens.net/bibentry


%: Macro file for Latex
% Macros help you summarise frequently repeated Latex commands.
% Here, they are placed in an external file /Latex/Macros/MacroFile1.tex
% An macro that you may use frequently is the figuremacro (see introduction.tex)
%---------------------------------------------------------------
% Macros
% version 3 by Igor Ruiz-Agundez 2011
% version 2 by Jakob Suckale 2007
% version 1 by Harish Bhanderi 2002
%---------------------------------------------------------------

% This file contains macros that can be called up from connected TeX files
% It helps to summarise repeated code, e.g. figure insertion (see below).



%---------------------------------------------------------------
% Figures
%---------------------------------------------------------------


% Makes the \InsertFig macro compatible both with one or two columns
\makeatletter
\newlength \figwidth
\if@twocolumn
  \setlength \figwidth {\columnwidth}
\else
  \setlength \figwidth {\textwidth}
\fi
\makeatother

% \InsertFig allows inserting figures
% Parameters
% 1 --> Filename
% 2 --> Label for referencing
% 3 --> Title describing the figure (caption)
% 4 --> Description of the figure
% 5 --> Figure width, range [0,1]. If parameter is left blank the figure size is not change
% 6 --> Any other option for \includegraphics
% Usage:
% \InsertFig{}{}{}{}{}{}
%

\newcommand{\InsertTab}[5]{
  \ifthenelse{\isempty{#5}}
  { % if empty
    \begin{savenotes}
    \begin{table}
      \caption{#2}{\footnotesize{#3} \medskip}
      
      \centering
      #4
      \label{#1}
    \end{table}
    \end{savenotes}
  }
  { % if not empty
    \begin{savenotes}
    \begin{table}[#5]
      \caption{#2}{\footnotesize{#3} \medskip}
      
      \centering
      #4
      \label{#1}
    \end{table}
    \end{savenotes}
  }
}

\newcommand{\InsertFig}[6]{%
	\ifthenelse{\isempty{#5}}%
	{% if #1 is empty
		\begin{figure}[htbp!]
		\centering
		\includegraphics[#6]{#1}%
		\caption{#3}{\footnotesize{#4}}
		\label{#2}
		\end{figure}    
	}
	{% if #1 is not empty
		\begin{figure}[htbp!]
		\centering
		\includegraphics[width=#5\figwidth,#6]{#1}%
		\caption{#3}{\footnotesize{#4}}
		\label{#2}
		\end{figure}
	}
}

%% Simple version of \InsertFig
%\newcommand{\InsertFig}[5]{
%  \begin{figure}[htbp]
%   	\centering
%    \includegraphics[width=#4\textwidth,#5]{#1}%
%    \caption{#3}
%    \label{#2}
%  \end{figure}
%}



% insert a centered figure with caption
% parameters 1:filename, 2:label, 3:title, 
\newcommand{\figuremacro}[3]{
	\begin{figure}[htbp]
		\centering
		\includegraphics[width=1\textwidth]{#1}
		\caption[#3]{\textbf{#3}}
		\label{#2}
	\end{figure}
}


% insert a centered figure with caption and description
% parameters 1:filename, 2:label, 3:title, 4:description
\newcommand{\figuremacroD}[4]{
	\begin{figure}[htbp]
		\centering
		\includegraphics[width=1\textwidth]{#1}
		\caption[#3]{\textbf{#3} - #4}
		\label{#2}
	\end{figure}
}

% insert a centered figure with caption and description AND WIDTH
% parameters 1:filename, 2:label, 3:title, 4:description, 5: textwidth
% textwidth 1 means as text, 0.5 means half the width of the text
\newcommand{\figuremacroDW}[5]{
	\begin{figure}[htbp]
		\centering
		\includegraphics[width=#5\textwidth]{#1}
		\caption[#3]{\textbf{#3} - #4}
		\label{#2}
	\end{figure}
}

% inserts a figure with wrapped around text; only suitable for NARROW figs
% o is for outside on a double paged document; others: l, r, i(inside)
% text and figure will each be half of the document width
% note: long captions often crash with adjacent content; take care
% in general: above 2 macro produce more reliable layout
\newcommand{\figuremacroN}[3]{
	\begin{wrapfigure}{o}{0.5\textwidth}
		\centering
		\includegraphics[width=0.48\textwidth]{#1}
		\caption[#2]{{\small\textbf{#2} - #3}}
		\label{#1}
	\end{wrapfigure}
}




% Estas definiciones son para el comando \InsertFigBox
\newlength{\anchoFigura}
\newlength{\anchoFloat}
\addtolength{\fboxsep}{2\fboxsep}
%\renewcommand{\capfont}{\normalfont\normalcolor\sffamily\small}
%\renewcommand{\caplabelfont}{\normalfont\normalcolor\sffamily\bfseries\small}

% El comando \InsertFigBox nos permite insertar figuras en un marco
% Los parametros son:
% 1 --> Fichero de la imagen
% 2 --> Etiqueta (label) para referencias
% 3 --> Texto a pie de imagen
% 4 -> Porcentaje del ancho de página que ocupará la figura (de 0 a 1)
% 5 --> Opciones que queramos pasarle al \includegraphics
\newcommand{\InsertFigBox}[5]{%
  \setlength{\anchoFloat}{#4\textwidth}%
  \addtolength{\anchoFloat}{-4\fboxsep}%
  \setlength{\anchoFigura}{\anchoFloat}%
  \begin{figure}%
    \begin{center}%
      \Ovalbox{%
        \begin{minipage}{\anchoFloat}%
          \begin{center}%
            \includegraphics[width=\anchoFigura,#5]{#1}%
            \caption{#3}%
            \label{#2}%
          \end{center}%
        \end{minipage}
      }%
    \end{center}%
  \end{figure}%
}



%---------------------------------------------------------------
% Misc
%---------------------------------------------------------------

% predefined commands by Harish
\newcommand{\PdfPsText}[2]{
  \ifpdf
     #1
  \else
     #2
  \fi
}


%---------------------------------------------------------------
% Locales
%---------------------------------------------------------------


%%
%% Para quitar traducciones raras (Cuadros)
%% A de usarse cada vez que se seleccione el idioma
%%
\newcommand{\MejorarTraducciones}{%
       \renewcommand{\listtablename}{Índice de tablas}
       \renewcommand{\tablename}{Tabla}
       \renewcommand{\lstlistingname}{Lista}
}%



%---------------------------------------------------------------
% Source code
%---------------------------------------------------------------


%%
%% Para escribir extractos de codigo
%%
%% Las tabulaciones se substituyen por dos espacios
%\fvset{tabsize=2}
%% Creamos un nuevo environment de fancyvrb para los ejemplos enmarcados
%\DefineVerbatimEnvironment{VerbEj}{BVerbatim}{fontsize=\small,samepage=true,commandchars=\\\{\}}
%% Colo de fondo
%\definecolor{grisfondo}{gray}{0.9}
%% Environment para extractos de codigo
%\newenvironment{codigo}%
%{\VerbatimEnvironment\begin{Sbox}\begin{VerbEj}}%
%{\end{VerbEj}\end{Sbox}\setlength{\fboxsep}{8pt}\begin{center}\fcolorbox{black}{grisfondo}{\TheSbox}\end{center}}
%
%% Otro formato más bonito para código fuente
%\newcommand{\codigofuente}[3]{%
%  \lstinputlisting[language=#1,caption={#2}]{#3}%
%}






%: ----------------------------------------------------------------------
%:                  TITLE PAGE: name, degree,..
% ----------------------------------------------------------------------
% below is to generate the title page with crest and author name

% if output to PDF then put the following in PDF header
\ifpdf  
    \pdfinfo { /Title  (PhD)
               /Creator (TeX)
               /Producer (pdfTeX)
               /Author (Name surname)
               /CreationDate (D:YYYYMMDDhhmmss)  %format D:YYYYMMDDhhmmss
               /ModDate (D:YYYYMMDDhhmm)
               /Subject (xyz)
               /Keywords (keyword1, keyword2, keyword3) }
    \pdfcatalog { /PageMode (/UseOutlines)
                  /OpenAction (fitbh)  }
\fi





% ----------------------------------------------------------------------
% This section below defines front covert (external and internal)
% Shield logo
\crest{\includegraphics[width=2cm]{Deusto_Shield}}
% Full logo
%\crest{\includegraphics[width=6cm]{UDeusto}}
\makeatletter
% check http://tex.stackexchange.com/questions/27710/how-can-i-use-author-date-and-title-after-maketitle
% Title of the dissertation
\title{Semantic Tuples Spaces for Constrained Devices: A Web-compliant Vision} \let\Title\@title
\university{Universidad de Deusto}
\degree{Tesis doctoral presentada por }
\author{Aitor Gómez Goiri} \let\Author\@author
\collegeordept{dentro del Programa de Doctorado en Your programme name}
\textadvisor{Dirigida por }
\advisor{Dr. Diego López de Ipiña} \let\Advisor\@advisor
\textadvisortwo{y}
\advisortwo{Dr. Íñigo Goiri Presa} \let\Advisortwo\@advisortwo
\textsignaturecandidate{El doctorando}
\textsignatureadvisor{El director}
\cityofbirth{Bilbao}
%\degreedate{\monthname \ \the\year}
\degreedate{Month de \the\year}
\makeatother

\newcommand{\website}{\url{http://gomezgoiri.net}}


% ----------------------------------------------------------------------
       
% turn of those nasty overfull and underfull hboxes
\hbadness=10000
\hfuzz=50pt


%: --------------------------------------------------------------
%:                  FRONT MATTER: dedications, abstract,..
% --------------------------------------------------------------

\begin{document}

\selectlanguage{british}

% sets line spacing
\renewcommand\baselinestretch{1.2}
\baselineskip=18pt plus1pt

% Watermark
%\watermark{DRAFT	DRAFT	DRAFT	DRAFT	DRAFT	DRAFT	DRAFT	DRAFT	DRAFT}


%: ----------------------- generate cover page ------------------------

\maketitle  % command to print the title page with above variables



% Title back
% Thesis Titleback ---------------------------------------------------

\thispagestyle{empty}

\hfill

\vfill

\medskip


\noindent
\textit{
\Title
}




Author: \Author

Advisor: \Advisor

Advisor: \Advisortwo



\vfill

\vfill

\noindent
The following web-page address contains up to date information about this dissertation and related topics: \\
\website


\noindent
Text printed in Bilbao

% TODO final date
\noindent
First edition, 
% Moth and year
\monthname \ \the\year

\vspace{1cm}
\hrule
\bigskip

% \cleardoublepage command ends the current page and causes all figures and tables that have so far appeared in the input to be printed. In a two-sided printing style, it also makes the next page a right-hand (odd-numbered) page, producing a blank page if necessary. 
\cleardoublepage

%%: ----------------------- cover page back side ------------------------
%% Your research institution may require reviewer names, etc.
%% This cover back side is required by Dresden Med Fac; uncomment if needed.
%
%\newpage
%\vspace{10mm}
%1. Reviewer: Name
%
%\vspace{10mm}
%2. Reviewer: 
%
%\vspace{20mm}
%Day of the defense:
%
%\vspace{20mm}
%\hspace{70mm}Signature from head of PhD committee:
%
%
%\cleardoublepage

% ----------------------------------------------------------------------



%: ----------------------- abstract ------------------------

% Your institution may have specific regulations if you need an abstract and where it is to be placed in the document. The default here is just after title.


% The original template provides and abstractseparate environment, if your institution requires them to be separate. I think it's easier to print the abstract from the complete thesis by restricting printing to the relevant page.
% \begin{abstractseparate}
%   
% Thesis Abstract -----------------------------------------------------


%\begin{abstractslong}    %uncommenting this line, gives a different abstract heading


\begin{abstracts}        %this creates the heading for the abstract page
% Ubicomp, mobile e IoT
\acf{ubicomp} envisions environments where devices interact among them to seamlessly work on behalf of the humans.
% TODO TODO TODO In recent years,... in our days.
In recent years, the emergence of the \acf{iot} concept, which opts for connecting everyday objects to the Internet, and the Mobile Computing paradigm have contributed to strengthen UbiComp in our days.
For this reason, UbiComp environments are not necessarily populated by powerful computers.
On the contrary, mainly resource constrained devices (e.g., embedded and mobile devices) are the main actors of these environments.
% Retos de Ubicomp
Thus, it is important for the environment to deal with their heterogeneity, unreliability, and replaceability.


% Tuplespaces y semántica
To face heterogeneity, the Semantic Web has proposed several standards and models to provide a precise meaning to the terms so that they can be reused across applications boundaries.
Regarding unreliability and replaceability, space-based computing (or Tuple Spaces) promotes the uncoupled coordination of the devices.
Solutions based on semantic tuple spaces combine these three beneficial aspects, resulting from bridging the Semantic Web and Tuple Spaces domains, for UbiComp.

% Reto que supone adaptar el paradigma para cumplir los objetivos de la hipótesis
Most of these semantic tuple spaces consider embedded and mobile devices as mere clients of a space managed by more powerful devices.
Such delegation helps to reduce the workload of devices with computing and energy limitations.
However,  this delegation moves the data away from where it is physically generated.
This creates a conflict between providing updated data and generating unnecessary network traffic for unused information.
% TODO TODO TODO each other no!
Besides, this delegation makes constrained devices intrinsically dependent on each other when it might not always be necessary.
This dissertation explores how these constrained devices can act as fully fledged semantic knowledge providers to create a more decentralized space.

% Aspectos que se analizan en la tesis
In conclusion, this dissertation presents a novel adaptation of semantic tuple space which considers the energy and computational impact on the devices.
Specifically, this dissertation proposes the following contributions:
\begin{itemize}
  \item A space model which considers the principles which made the web flourish in the last decades, together with the uncoupling properties of the space-based computing.
  \item An energy-aware search mechanism for autonomous constrained devices.
  \item An alignment of two approaches to act on the physical environment, namely a space-based indirect actuation and a web-based direct actuation.
\end{itemize}

\end{abstracts}



\begin{resumen}        %this creates the heading for the abstract page
La computación ubicua (\acs{ubicomp}) vislumbra entornos donde los dispositivos interactuan entre sí para trabajar de forma imperceptible para los seres humanos.
En los últimos años, la emergencia del Internet de las Cosas (\acs{iot}), que aboga por conectar objetos cotidianos a Internet, y la computación móvil han contribuído a fortalecer \ac{ubicomp}.
Por ello, los entornos ubicuos no se componen necesariamente de computadoras potentes.
Por el contrario, los actores principales de dichos entornos suelen ser en su mayor parte dispositivos con recursos limitados (p.e. dispositivos móviles y embebidos).
Es por ello que es importante para estos entornos enfrentarse a la heterogeneidad, falta de fiabilidad y facilidad de reemplazo de estos dispositivos.


Para enfrentarse a la heterogeneidad, la web semántica propone diversos estándares y modelos para proveer un significado precido a los términos para que puedan reusarse más allá de las fronteras marcadas por las aplicaciones.
Relativo a la falta de fiabilidad y facilidad de reemplazo, la computación basada en espacios (o \acl{ts}) promueve una coordinación desacoplada de los dispositivos.
Las soluciones basadas en espacios de tuplas semánticos combinan estos tres aspectos beneficiosos para \ac{ubicomp} uniendo los dominios de la web semántica y los \acl{ts}.


Muchos de estos espacios de tuplas semánticos consideran a los dispositivos móviles y embebidos como meros clientes de un espacio gestionado por dispositivos más potentes.
Dicha delegación ayuda a reducir la carga de trabajo de los dispositivos con limitaciones computacionales y energéticas.
Sin embargo, esta delegación distancia los datos de donde son generados.
Esto crea un conflicto entre proveer datos actualizados y generar tráfico de red innecesario para información no usada.
Además, esta delegación hace a los dispositivos limitados intrinsicamente dependiente de otros cuando no siempre es necesario.
Esta tesis explora cómo esos dispositivos limitados pueden actuar como autenticos proveedores de conocimiento semántico para crear un espacio más descentralizado.


En conclusión, esta tesis presenta una adaptación novedora de los espacios de tuplas semánticos que considera el impacto computacional y energético en los dispositivos.
Especificamente, esta disertación propone las siguientes contribuciones:
\begin{itemize}
  \item Un modelo de espacio que considera los principios que han hecho florecer a la web en las últimas décadas, junto con las propiedades de desacoplamiento de la computación basada en espacios.
  \item Un mecanismo de búsqueda energéticamente eficiente para dispositivos limitados autónomos.
  \item Un alineamiento entre dos formas de actuar en un entorno físico: actuación indirecta basada en espacios y actuación directa basada en la web.
\end{itemize}

\end{resumen}



\begin{laburpena}        %this creates the heading for the abstract page
% Jarri zure laburpena hemen.


Vivamus ullamcorper erat in sapien dignissim pellentesque.


\end{laburpena}

%\end{abstractlongs}


% ---------------------------------------------------------------------- 

% \end{abstractseparate}


%: ----------------------- tie in front matter ------------------------

% The frontmatter text starts here
\frontmatter

% Thesis Dedictation ---------------------------------------------------

\begin{dedication} %this creates the heading for the dedication page

\textit{
Se lo dedico a todo doctorando que no vea la luz al final del tunel o se sienta desbordado por la frustración. Ánimo.
}

% TODO Dibujo de PhD comics

\emph{Sometimes, there's so much beauty in the world... I feel like I can't take it, like my heart is just going to cave in.}
\qauthor{Ricky Fitts, American Beauty}

\end{dedication}

% ----------------------------------------------------------------------


% Thesis Abstract -----------------------------------------------------


%\begin{abstractslong}    %uncommenting this line, gives a different abstract heading


\begin{abstracts}        %this creates the heading for the abstract page
% Ubicomp, mobile e IoT
\acf{ubicomp} envisions environments where devices interact among them to seamlessly work on behalf of the humans.
% TODO TODO TODO In recent years,... in our days.
In recent years, the emergence of the \acf{iot} concept, which opts for connecting everyday objects to the Internet, and the Mobile Computing paradigm have contributed to strengthen UbiComp in our days.
For this reason, UbiComp environments are not necessarily populated by powerful computers.
On the contrary, mainly resource constrained devices (e.g., embedded and mobile devices) are the main actors of these environments.
% Retos de Ubicomp
Thus, it is important for the environment to deal with their heterogeneity, unreliability, and replaceability.


% Tuplespaces y semántica
To face heterogeneity, the Semantic Web has proposed several standards and models to provide a precise meaning to the terms so that they can be reused across applications boundaries.
Regarding unreliability and replaceability, space-based computing (or Tuple Spaces) promotes the uncoupled coordination of the devices.
Solutions based on semantic tuple spaces combine these three beneficial aspects, resulting from bridging the Semantic Web and Tuple Spaces domains, for UbiComp.

% Reto que supone adaptar el paradigma para cumplir los objetivos de la hipótesis
Most of these semantic tuple spaces consider embedded and mobile devices as mere clients of a space managed by more powerful devices.
Such delegation helps to reduce the workload of devices with computing and energy limitations.
However,  this delegation moves the data away from where it is physically generated.
This creates a conflict between providing updated data and generating unnecessary network traffic for unused information.
% TODO TODO TODO each other no!
Besides, this delegation makes constrained devices intrinsically dependent on each other when it might not always be necessary.
This dissertation explores how these constrained devices can act as fully fledged semantic knowledge providers to create a more decentralized space.

% Aspectos que se analizan en la tesis
In conclusion, this dissertation presents a novel adaptation of semantic tuple space which considers the energy and computational impact on the devices.
Specifically, this dissertation proposes the following contributions:
\begin{itemize}
  \item A space model which considers the principles which made the web flourish in the last decades, together with the uncoupling properties of the space-based computing.
  \item An energy-aware search mechanism for autonomous constrained devices.
  \item An alignment of two approaches to act on the physical environment, namely a space-based indirect actuation and a web-based direct actuation.
\end{itemize}

\end{abstracts}



\begin{resumen}        %this creates the heading for the abstract page
La computación ubicua (\acs{ubicomp}) vislumbra entornos donde los dispositivos interactuan entre sí para trabajar de forma imperceptible para los seres humanos.
En los últimos años, la emergencia del Internet de las Cosas (\acs{iot}), que aboga por conectar objetos cotidianos a Internet, y la computación móvil han contribuído a fortalecer \ac{ubicomp}.
Por ello, los entornos ubicuos no se componen necesariamente de computadoras potentes.
Por el contrario, los actores principales de dichos entornos suelen ser en su mayor parte dispositivos con recursos limitados (p.e. dispositivos móviles y embebidos).
Es por ello que es importante para estos entornos enfrentarse a la heterogeneidad, falta de fiabilidad y facilidad de reemplazo de estos dispositivos.


Para enfrentarse a la heterogeneidad, la web semántica propone diversos estándares y modelos para proveer un significado precido a los términos para que puedan reusarse más allá de las fronteras marcadas por las aplicaciones.
Relativo a la falta de fiabilidad y facilidad de reemplazo, la computación basada en espacios (o \acl{ts}) promueve una coordinación desacoplada de los dispositivos.
Las soluciones basadas en espacios de tuplas semánticos combinan estos tres aspectos beneficiosos para \ac{ubicomp} uniendo los dominios de la web semántica y los \acl{ts}.


Muchos de estos espacios de tuplas semánticos consideran a los dispositivos móviles y embebidos como meros clientes de un espacio gestionado por dispositivos más potentes.
Dicha delegación ayuda a reducir la carga de trabajo de los dispositivos con limitaciones computacionales y energéticas.
Sin embargo, esta delegación distancia los datos de donde son generados.
Esto crea un conflicto entre proveer datos actualizados y generar tráfico de red innecesario para información no usada.
Además, esta delegación hace a los dispositivos limitados intrinsicamente dependiente de otros cuando no siempre es necesario.
Esta tesis explora cómo esos dispositivos limitados pueden actuar como autenticos proveedores de conocimiento semántico para crear un espacio más descentralizado.


En conclusión, esta tesis presenta una adaptación novedora de los espacios de tuplas semánticos que considera el impacto computacional y energético en los dispositivos.
Especificamente, esta disertación propone las siguientes contribuciones:
\begin{itemize}
  \item Un modelo de espacio que considera los principios que han hecho florecer a la web en las últimas décadas, junto con las propiedades de desacoplamiento de la computación basada en espacios.
  \item Un mecanismo de búsqueda energéticamente eficiente para dispositivos limitados autónomos.
  \item Un alineamiento entre dos formas de actuar en un entorno físico: actuación indirecta basada en espacios y actuación directa basada en la web.
\end{itemize}

\end{resumen}



\begin{laburpena}        %this creates the heading for the abstract page
% Jarri zure laburpena hemen.


Vivamus ullamcorper erat in sapien dignissim pellentesque.


\end{laburpena}

%\end{abstractlongs}


% ---------------------------------------------------------------------- 


% Thesis Acknowledgements ------------------------------------------------


% Opening of the acknowledgements

%Sort version
%this creates the heading for the acknowlegments

\begin{acknowledgements}      
%Long version
%uncommenting this line, gives a different acknowledgements heading
%\begin{acknowledgementslong} 

\begin{flushright}
\emph{Sometimes, there's so much beauty in the world - I feel like I can't take it, like my heart is just going to cave in.}
\qauthor{Ricky Fitts (American Beauty, 1999).}
\end{flushright}

\bigskip

Si debo a alguien la posibilidad de haber concluido esta aventura es sin duda a aita y ama.
Esta posibilidad va más allá del mundano y sin embargo tan necesario hecho de haberme proporcionado una cómoda vida en la que siempre pude elegir mi camino.
Por vosotros soy como soy y soy quien soy.
Podría haber salido un hijo un poco más hablador y menos pedante, pero no será porque no os esforzasteis.
Echando la vista atrás no puedo sino estar orgullosísimo de la educación recibida.
Por ello, vaya para vosotros mi primer agradecimiento.


Gracias Diego por haberme dado esta enorme oportunidad.
Gracias por el tiempo y el espacio concedido.
Gracias Íñigo por soportarme en momentos de máxima incertidumbre y pesimismo.
Gracias por servirme de guía.


Gracias a todos mis compañeros y ex-compañeros.
Gracias a mis compañeros de generación por compartir sufrimientos: Edu, Iván y Eguiluz.
Gracias a esa infatigable e incansable máquina con piernas de la que tanto se aprende que es Pablo Orduña.
Gracias a los ancianos del lugar: Aitor y Unai.
Gracias a la gente del Garilab por aceptarme durante mi exilio.
En definitiva, gracias a tanta gente que ha pasado por la vida \emph{esmarlabense} estos últimos años: Don Jaime, Emaldi, Laiseca, el innombrable, Lázaro y un largo etc.
Afortunadamente, ese ``\emph{largo etc}'' se puede consultar en la página web del grupo haciendo uso de la consulta mostrada al final de esta sección o siguiendo el siguiente \emph{QRCode}.
% Poner el SPARQL para seleccionarlos!



\begin{figure}[h]
    \centering
    \includegraphics[width=0.3\textwidth]{qrcode_acks}
\end{figure}


Thanks to all the inspiring people I've known during this journey.
Thanks to Geoff Coulson, Barry Porter and the rest of the people I met at Lancaster University.
Thanks to all the ETH's \emph{Distributed Systems group} members, specially to Simon Mayer, my cicerone in Zurich.
Thanks to rest of you, no matter if I've met you in person or not.


Gracias a todos mis compañeros de aventuras estos años.
A los ocasionales y a los fieles.
En viajes por el mundo, canchas de baloncesto, senderos, \emph{eventos del esmarla},...
Cada segundo en que mantuve la cabeza fuera de la tesis fue un respiro enorme.


Gracias al resto de mi familia, la sanguínea y la \emph{circunstancial}.
Especialmente a mi hermano Mikel y a mi abuela Teo.


\begin{flushright}
\textit{Eskerrik asko,}

Aitor Gómez Goiri

% Month and year
\monthname \ \the\year

% Signature figure

%\begin{figure}[htbp!]
%\end{figure}
%\includegraphics{signature}%

\end{flushright}

\begin{listing}[nolol]
  \scriptsize
  \expandafter\def\csname PY@tok@err\endcsname{}
\begin{Verbatim}[commandchars=\\\{\}]
\PY{k}{PREFIX }\PY{n+nv}{swrcfe:}\PY{n+nn}{ \PYZlt{}http://www.morelab.deusto.es/ontologies/swrcfe.owl\PYZsh{}\PYZgt{}}

\PY{k}{SELECT }\PY{err}{C}\PY{err}{O}\PY{err}{U}\PY{err}{N}\PY{err}{T}\PY{err}{(}\PY{err}{D}\PY{err}{I}\PY{err}{S}\PY{err}{T}\PY{err}{I}\PY{err}{N}\PY{err}{C}\PY{err}{T}\PY{err}{ }\PY{err}{?}\PY{err}{f}\PY{err}{a}\PY{err}{n}\PY{err}{a}\PY{err}{m}\PY{err}{e}\PY{err}{)} \PY{k}{WHERE}\PY{p}{\PYZob{}}
\PY{n+nv}{  ?job}\PY{o}{ swrcfe:doneBy}
\PY{o}{     }\PY{n+na}{\PYZlt{}http://www.morelab.deusto.es/labman/resource/people/aitor\PYZhy{}gomez\PYZhy{}goiri\PYZgt{}} \PY{p}{.}
  \PY{p}{\PYZob{}}
    \PY{n+nv}{?job2}\PY{o}{ swrcfe:doneBy }\PY{n+nv}{?formerworkmate }\PY{p}{.}
    \PY{n+nv}{?formerworkmate}\PY{o}{ foaf:familyName }\PY{n+nv}{?faname }\PY{p}{.}
    \PY{n+nv}{?formerworkmate}\PY{o}{ foaf:firstName }\PY{n+nv}{?finame }\PY{p}{.}
    \PY{n+nv}{?job}\PY{o}{ swrcfe:jobStartDate }\PY{n+nv}{?jobstartdate }\PY{p}{.}

    \PY{k}{OPTIONAL} \PY{p}{\PYZob{}}\PY{n+nv}{ ?job2}\PY{o}{ swrcfe:jobEndDate  }\PY{n+nv}{?jobenddate }\PY{err}{\PYZcb{}}
\PY{n+na}{    FILTER}\PY{p}{(}\PY{n+nv}{?jobstartdate }\PY{err}{\PYZlt{}}\PY{err}{ }\PY{n+nv}{?jobenddate}\PY{p}{)}
  \PY{p}{\PYZcb{}}\PY{err}{ }\PY{err}{U}\PY{err}{N}\PY{err}{I}\PY{err}{O}\PY{err}{N} \PY{p}{\PYZob{}}
\PY{n+nv}{    ?job2}\PY{o}{ swrcfe:doneBy }\PY{n+nv}{?workmate }\PY{p}{.}
    \PY{n+nv}{?workmate}\PY{o}{ foaf:familyName }\PY{n+nv}{?faname }\PY{p}{.}
    \PY{n+nv}{?workmate}\PY{o}{ foaf:firstName }\PY{n+nv}{?finame }\PY{p}{.}

    \PY{k}{OPTIONAL} \PY{p}{\PYZob{}}\PY{n+nv}{ ?job2}\PY{o}{ swrcfe:jobEndDate }\PY{n+nv}{?jobenddate }\PY{err}{\PYZcb{}}
\PY{n+na}{    FILTER}\PY{p}{(}\PY{err}{!}\PY{n+na}{bound}\PY{p}{(}\PY{n+nv}{?jobenddate}\PY{p}{)}\PY{p}{)}
  \PY{p}{\PYZcb{}}
\PY{err}{ }\PY{err}{ }\PY{err}{F}\PY{err}{I}\PY{err}{L}\PY{err}{T}\PY{err}{E}\PY{err}{R}\PY{err}{(}\PY{err}{?}\PY{err}{j}\PY{err}{o}\PY{err}{b}\PY{err}{ }\PY{err}{!}\PY{err}{=}\PY{err}{ }\PY{err}{?}\PY{err}{j}\PY{err}{o}\PY{err}{b}\PY{err}{2}\PY{err}{)}
\PY{p}{\PYZcb{}}\PY{err}{ }
\end{Verbatim}

  %\caption{Consulta SPARQL para obtener la gente relacionada con el grupo MORElab a la que tengo que estar agradecido.}
  \label{lst:workmates_ack}
\end{listing}



%Closing of the acknowledgements
%Sort version
\end{acknowledgements}
% Long version
%\end{acknowledgementslong}

% ------------------------------------------------------------------------




% As abstract contains various languages we set the main language again
\selectlanguage{british}


%: ----------------------- contents ------------------------

\setcounter{secnumdepth}{5} % organisational level that receives a numbers
\setcounter{tocdepth}{2}    % print table of contents for level 3
% "Note that if you have a list of figures/tables, the tocdepth must be at least 2 for them to show up"

\hypersetup{ % requires hyperref package
  bookmarksdepth = 3, % for the PDF's table of contents!
}


%%You can also add extra lines to the ToC or to force extra unnumbered section headings to be included. For example, if you wanted to add an entry called Preface, and you didn't want the Preface to be numbered, you'd use these commands:
%\ subsection*{Preface}
%\addcontentsline{toc}{subsection}{Preface} 

\tableofcontents            % print the table of contents
% levels are: 0 - chapter, 1 - section, 2 - subsection, 3 - subsubsection

%: ----------------------- list of figures/tables ------------------------

\listoffigures	% print list of figures
\listoftables  % print list of tables


%: ----------------------- glossary ------------------------

% Usando ACRONYMS, puedo definir palabras rollo Tuple Spaces para ser uniforme a lo largo del texto.
% Así evito acabar llamandolas tspaces, tuplespaces, o tuple spaces.
% Comprobar sobre texto escrito que uso variantes de \ac siempre
% Faltan por comprobar los capítulos 4, 5 y 6

%% acronyms
%% http://staff.science.uva.nl/~polko/HOWTO/LATEX/acronym.html
%\acrodef{label}[acronym]{written out form}

\acrodef{ts}[TS]{Tuple Space}
\acrodef{tsc}[TSC]{Triple Space Computing}
\acrodef{ami}[AmI]{Ambient Intelligence}
\acrodef{ubicomp}[UbiComp]{ubiquitous computing}
\acrodef{wot}[WoT]{Web of Things}
\acrodef{iot}[IoT]{Internet of Things}
\acrodef{sw}[SW]{Semantic Web}


% si decido gustarme: http://www.latex-community.org/know-how/263-glossaries-nomenclature-lists-of-symbols-and-acronyms
% Tie in external source file for definitions: /0_frontmatter/glossary.tex
% Glossary entries can also be defined in the main text. See glossary.tex
% this file is called up by thesis.tex
% content in this file will be fed into the main document

% Glossary entries are defined with the command \nomenclature{1}{2}
% 1 = Entry name, e.g. abbreviation; 2 = Explanation
% You can place all explanations in this separate file or declare them in the middle of the text. Either way they will be collected in the glossary.

% required to print nomenclature name to page header
\markboth{\MakeUppercase{\nomname}}{\MakeUppercase{\nomname}}

% ----------------------- contents from here ------------------------
%

%
%
%% acronyms

\nomenclature{AmI}{Ambient Intelligence}
\nomenclature{IP}{Internet Protocol}
\nomenclature{WoT}{Web of Things}
\nomenclature{TS}{Tuple Space}
\nomenclature{TSC}{Triple Space Computing}
\nomenclature{UbiComp}{Ubiquitous Computing}




%\begin{multicols}{2} % \begin{multicols}{#columns}[header text][space]
%\begin{footnotesize} % scriptsize(7) < footnotesize(8) < small (9) < normal (10)

%\printnomenclature[1.5cm] % [] = distance between entry and description

% I don't know why it does not work.
% TODO fix or remove (acronyms are already listed)
\printnomenclature % [] = distance between entry and description

\label{sec:glossary} % target name for links to glossary

%\end{footnotesize}
%\end{multicols}




%: --------------------------------------------------------------
%:                  MAIN DOCUMENT SECTION
% --------------------------------------------------------------

% the main text starts here with the introduction, 1st chapter,...
\mainmatter

%\renewcommand{\chaptername}{} % uncomment to print only "1" not "Chapter 1"
\pagestyle{fancy}


%: ----------------------- subdocuments ------------------------

% Parts of the thesis are included below. Rename the files as required.
% But take care that the paths match. You can also change the order of appearance by moving the include commands.

%: ----------------------- introduction ------------------------

% this file is called up by thesis.tex
% content in this file will be fed into the main document

%: ----------------------- introduction file header -----------------------


\begin{savequote}[50mm]
The beginning is the most important part of the work. 
\qauthor{Plato}
\end{savequote}

\chapter{Introduction}
\label{cha:Introduction}

% the code below specifies where the figures are stored
\ifpdf
    \graphicspath{{1_introduction/figures/PNG/}{1_introduction/figures/PDF/}{1_introduction/figures/}}
\else
    \graphicspath{{1_introduction/figures/EPS/}{1_introduction/figures/}}
\fi


% ----------------------------------------------------------------------

\section{Motivation}
\label{sec:Motivation}
% Terminos aclarados: IoT, ubicomp o pervasive y AmI. (otros posibles: Context Awareness?)


% internet ahora lo es todo
% Se podría hasta poner una cita de Vinton Cerf :-P
Originally, the Internet was basically composed of a small number of computers that were physically connected to a wired network.
Over the years, the popularity of the Internet grew and connecting computers became easier and cheaper.
Thanks to wireless technologies, devices can connect to the Internet without having to be physically connected to a network.

% IoT + mobile => UbiComp
These technologies have contributed to the emergence of end-user devices such as smartphones or tablets.
Furthermore, nowadays, everyday objects like cars or washing machines start to be connected to the Internet to exchange information.
This is what is currently known as the \ac{iot} \citep{atzori_internet_2010}.
The objects of the \ac{iot}, together with mobile devices, constitute the clearest sign of the \acf{ubicomp} prominence in our current lives \citep{caceres_ubicomp_2012}.

% UbiComp & AmI
The \ac{ubicomp} concept was envisioned by \citeauthor{weiser1991computer} in the early nineties \citep{weiser1991computer}. % y el term was coined
\citeauthor{weiser1991computer} defended that devices should imperceptibly work on our behalf. % qué es Ubicomp
This idea is particularly stressed by the concept of \ac{ami} \citep{ramos_ambient_2008}.
For instance, the washing machine could plan to finish washing the clothes two minutes before the car gets home.


\bigskip


% Cómo se pueden comunicar los cacharros?
\citeauthor{weiser1991computer} remarked that the real power of \ac{ubicomp} ``\emph{comes not from any one of these devices, it emerges from the interaction of all of them}''.
% WoT
However, integrating devices to accomplish such task is not trivial as they usually communicate using different protocols.
To solve this problem, the \ac{wot} initiative proposes to use well-established web standards to ease their communication \citep{guinard_thesis_2011}.
This, together with the already existing tools and libraries from the web world, simplifies its adoption by developers \citep{guinard_search_2011}.
Besides, the \ac{wot} naturally integrates with the currently widely used \acs{rest}ful web applications.

% Debería hacer el salto de:
% WoT problem coupling => TS ???

% TS
Parallel to this initiative, \aclp{ts} (or \emph{space-based computing}) \citep{gelernter_generative_1985} proposes a different integration style based on the \emph{blackboard model}.
In \ac{ts} participants cooperate with each other by writing and reading information into a shared \emph{space}.
The main benefit of this paradigm is the higher-level of decoupling its indirect communication provides. % y dejo el explicar cuales son para el estado del arte


\bigskip

% Necesidad de semántica
One common problem of both the \ac{wot} and \aclp{ts} is that the format of the data they exchange is multifarious and application domain dependent.
This implies that data will not be meaningful in other domains unless a specialized system converts and reinterprets them.
A way to solve this problem is annotating the data semantically as proposed by the \ac{www} \citep{berners-lee_semantic_2001}.
Specifically, the \ac{tsc} paradigm \citep{fensel_triple-space_2004} employs this semantic knowledge on shared spaces following the \aclp{ts} approach.
Besides, we argue that \ac{tsc} can be designed to fulfil \ac{rest}'s principles.


However, adding semantics may suppose a burden for resource constrained devices.
To reduce this overhead in such devices, part of this computation is usually delegated to intermediaries \citep{honkola_smart-m3_2010}. % todavía no poner esta cita
This approach reduces the overhead of semantically annotated data but brings other problems:
\begin{enumerate}
 \item When devices rely on others to provide information, it is not guaranteed that the information accessed will accurately represent the last information available in the data providers (i.e., the sensors).
 \item Once we rely on those intermediaries, it is required for them to be available at all time.
	Otherwise, the devices will not be able to talk to each other.
\end{enumerate}


\bigskip


% definir más claramente si IoT o AmI
In this thesis, we propose to adapt \ac{tsc} to the \ac{ubicomp} environments, particularly involving limited devices in such task.
To fulfil such an ambitious goal, we need to carefully consider the restrictions imposed by such environments.
Two key aspects to take into account are the limited energy autonomy or computational capacity of many devices.
Similarly, our solution should be flexible enough to support a wide range of scenarios and to reuse third-party applications' data.


This thesis addresses different key aspects of this adaptation:
\begin{itemize}
  \item \emph{\ac{tsc}'s compatibility with the \acs{rest} style and the \ac{wot} vision.}
	Aligning the \ac{tsc} design with \ac{rest} one can inherit many of the benefits provided by this style and interoperate with other \ac{rest}-based solutions.
	At the same time, the use of \ac{tsc} can help developers to focus on higher-level problems.
	% problema: habría que probar o se da por sentado que es lo que hace un middleware?
	
  \item \emph{Decentralized search on the semantic content provided by federated devices.}
	% TODO Se sobreentenderá lo de federated? es la primera vez que lo menciono, pero es un concepto común...
	This search architecture enables the direct communication between participants.
	Furthermore, it reduces the energy consumption of the resource-constrained devices.
	
  \item \emph{Actuation on the physical environment.}
	We present common usage patterns for space-based computing and we analyse how they can be used in \ac{ubicomp}.
	Besides, we propose a space enhancement which seamlessly reuses external \ac{rest} services.
	This way, we build a bridge between our solution and other \ac{wot} and \ac{rest} solutions.
\end{itemize}
%Finally, we also present the incarnation of some of these ideas in a middleware,
%and some scenarios where it has been used.

%: ----------------------- related work ------------------------
% ----------------------------------------------------------------------

\begin{savequote}[50mm]
I was born not knowing and have had only a little time to change that here and there.
\qauthor{Richard P. Feynman}
\end{savequote}


\chapter{State of the Art}
\label{cha:stateoftheart}
\newcommand{\pathchaptwo}{2_state_of_the_art}

% the code below specifies where the figures are stored
\ifpdf
    \graphicspath{{\pathchaptwo/figures/PNG/}{\pathchaptwo/figures/PDF/}{\pathchaptwo/figures/}}
\else
    \graphicspath{{\pathchaptwo/figures/EPS/}{\pathchaptwo/figures/}}
\fi


%------------------------------------------------------------------------- 

% TODO uniformizar labels!
In this chapter we present the related work.
We start introducing, categorizing and describing the related research topics in Section~\ref{sec:soa_intro}. % o llamarle research areas, important concepts?
% nombrarlos? \ac{rest}, \ac{ts} and the \ac{sw}
% "this survey is not intended to be exhaustive (nor exhausting)"
In Section~\ref{sec:soa_ubicomp}, rather than presenting an exhaustive (and exhausting) survey on these topics, we explore their impact in the \ac{ubicomp} field. % their impact vs their combination
To that end, we identify the relevant works remarking their interesting characteristics and how they relate to this dissertation.
Finally, Section~\ref{sec:soa_tsc} analyzes and compares other semantic \ac{ts} middlewares independently of their application domain.


\InsertFig{venn-summary}{fig:venn_ubicomp_tuple_semantic}{
  Background areas of the thesis, which focuses on their intersection.
}{
  This thesis focuses on the intersection of the three areas represented in the Figure.
  Therefore, this chapter will be driven by their analysis.
}{0.6}{}


% TODO valorar meter en otra sección descripción de IoT y UbiComp y sus características?
\section{Introduction}
\label{sec:soa_intro}

% Otra forma de ver la interop:
%  La European Telecommunication Standards Institute (ETSI) define cuatro capas:
%      technical interop: a nivel de comunicación (p.e. de acuerdo en las 7 capas de OSI)
%      syntactic interop
%      semantic interop
%      organizational interop


% no habría que definir de nuevo y de manera formal ubicomp o vale con la intro?
The IEEE \citep{_ieee_1990} defines \emph{interoperability} as ``\emph{the ability of two or more systems or components to exchange information and to use the information that has been exchanged}''.
The heterogeneity of technologies present in \ac{ubicomp} environments makes this a key property to consider.
The definition clearly distinguishes between two requirements: % o goals o incremental requirements
(1) to exchange information; and
(2) to use that information. % to understand others data (i.e. \emph{interoperation}).


% Buscar una referencia mejor: http://en.wikipedia.org/wiki/Interoperability
Exchanging information in distributed systems comprehends the communication between two systems.
% lo de ab-initio practicamente sólo lo he visto en la wikipedia
For the lower communication levels, we opt for interoperability \emph{ab-initio} relying on standard and widely accepted communication protocols. % e.g. HTTP
For a higher-level (i.e. application layer), Section~\ref{sec:soa_integration} categorizes different integration approaches.
This dissertation aims to delve into the \emph{space-based computing} approach. % whose benefits etc. are described
However, \ac{rest} architectures' properties have made them massively accepted mechanisms to integrate applications.
Consequently, we also take into consideration the latter mechanism in our solution design.


% Sintáctica vs semántica
Regarding the second goal, it is usually divided in \emph{syntactic interoperability} and \emph{semantic interoperability}.
On the one hand, syntactic interoperability is associated with the format of the data (i.e. its syntax and encoding) \citep{van_der_veer_achieving_2006}. % e.g. high level: HTML, XML, etc.
On the other hand, semantic interoperability is concerned with ensuring that the exchanged information has a precise meaning.
Its ultimate goal is to make the information ``\emph{understandable by any other application that was not initially developed for this purpose}'' \citep{_european_2004}.
Section~\ref{sec:soa_sw} explains well-accepted mechanisms to achieve the latter goal: the \acl{sw}. % TODO qué es la SW??? un mecanismo, conjunto de estándares, ¿?

\subsection{Application Integration}
\label{sec:soa_integration}

% Definiciones de middleware en el libro de Coulouris:
%
% (pag 17):
% Middleware • The term middleware applies to a software layer that provides a
% programming abstraction as well as masking the heterogeneity of the underlying
% networks, hardware, operating systems and programming languages.
%
% (en otra parte):
% The task of middleware is to provide a higher-level
% programming abstraction for the development of distributed systems and, through
% layering, to abstract over heterogeneity in the underlying infrastructure to promote
% interoperability and portability.


The integration of two applications is driven by how they communicate.
To ease that communication the applications use middlewares.
A middleware is a software layer which provides a higher level of abstraction and masks the underlying heterogeneity.
\citet{coulouris_distributed_2012} define two communication styles on the upper layer of a middleware: % figura 4.1 modificada para incluir elementos
the remote invocation and the indirect communication (see Figure~\ref{fig:middleware_layers}).

\InsertFig{middleware_layers}{fig:middleware_layers}{Middleware layers}{Middleware layers according to \citet{coulouris_distributed_2012} classification.}{1}{}
% According to Coulouris et al. HTTP is An example of a request-reply protocol!


The \emph{remote invocation} involves the most common two-way exchange between senders and receivers in distributed systems.
Within the \emph{remote invocation} style, RESTful style \citep{fielding_architectural_2000} and the well known set of WS-* \citep{alonso_web_2010} standards probably represent the most popular substyles.
In Section~\ref{sec:remote_invocation}, we analyze their application to the \ac{ubicomp} field, particularly focusing on the \ac{iot}.


\medskip


The \emph{indirect communication} style comprehends all the techniques with no direct coupling between the sender and the receiver.
% Otra forma:
% In contrast, the \emph{indirect communication} comprehends decoupled communications between senders and receivers.
The group communication, publish-subscribe systems, message queues or shared memory approaches are examples of indirect communication.
These paradigms are characterized by two key properties \citep{gelernter_generative_1985,coulouris_distributed_2012}
\footnote{
  We use the terminology of the  \aclp{ts}' seminal paper \citep{gelernter_generative_1985}, but please note that:
  % nomenclature, naming or terminology?
  \begin{itemize}
    %\item Gelernter et al. stated that \emph{distributed sharing} was just a consequence of the these properties.
    \item The \emph{space uncoupling} property is referred as \emph{reference autonomy} by some authors \citep{fensel_triple-space_2004}.
	  % El primero fue un tal Angerer en el 2002, pero en un artículo en alemán.
	  % En el 2003 hay un artículo suyo en Internet, pero no sé si mencionarlo.
    \item These authors mention a third property confusingly called \emph{space autonomy} (or \emph{location autonomy}).
	  According to \citet{fensel_triple-space_2004} this autonomy is achieved because:
	  \begin{quote}
	    The processes can run in completely different computational environments as long as both can access the same space.
	  \end{quote}
  \end{itemize}
}
:

\begin{itemize}
 \item \emph{Space uncoupling}, which is achieved when the sender does not need to know the receiver or receivers and vice versa.
 \item \emph{Time uncoupling}, which happens when senders and receivers do not need to exist in the same time\footnote{
	  Although some authors \citep{fensel_triple-space_2004,krummenacher_www_2005} explain this property just in terms of communication asynchrony,
	  % mencionar a otros? o no porque simplemente siguen lo dicho por Fensel?
	  % a Bundler no lo cito, porque era una master thesis "sólo"
	  % en este caso cito a krummenacher porque en esa publicación lo define directamente como el no uso de comunicación sincrona
	  \citet{coulouris_distributed_2012} make a clear distinction between them.
	  In their words, a communication is asynchronous when \emph{a sender sends a message and then continues without blocking},
	  whereas time uncoupling adds an extra dimension: \emph{the sender and the receiver can have independent existences}.
	  }.
 
\end{itemize}

\acl{ts} computing, also called space-based computing, offers an improvement of the shared memory approach.
Whereas the shared memory works at byte-level and accessing to memory addresses,
the \acl{ts} works with semi-structured data which is accessed in an associative manner.
In other words, in \ac{ts} the participants read data specifying patterns of interest.
In Section~\ref{sec:tuplespaces_eoa}, we study the most relevant \acl{ts} solutions proposed for the \acl{ubicomp}.


\input{\pathchaptwo/1_1_1_remote_invocation}
\input{\pathchaptwo/1_1_2_indirect_communication}
\subsection{The \acl{sw}}
\label{sec:soa_sw}

A problem of the initial view of the Web was that it was human-centered.
Regardless of whether the contents were machine processable, a human needed to interpret them to give them a meaning.
% mirar a ver cuando se empezó a usar \ac{sw} de verdad, porque en el artículo de 2001 dan a enterder que se llevaba un tiempo usando
In the early 2000s, \citet{berners-lee_semantic_2001} proposed the use of what they called the \acf{sw} to solve that problem.
The \ac{sw} was conceived as an extension of the ordinary web which would bring structure to its content.
As the ordinary web, the \ac{sw} benefits from the universality the hypertext provides by linking \emph{anything with anything}.
Besides, like the Internet, the \ac{sw} is intended to be as decentralized as possible.

The \acl{sw} and its key features are defined by the \emph{World Wide Web Consortium} \citep{semanticWeb-FAQ} as:
\begin{quote}
The vision of the \acl{sw} is to extend principles of the Web from documents to data.
Data should be accessed using the general Web architecture using, e.g., URI-s;
data should be \emph{related to one another} just as documents (or portions of documents) are already.
This also means creation of a common framework that allows data to be \emph{shared and reused} across application, enterprise, and community boundaries,
to be \emph{processed automatically} by tools as well as manually, including revealing possible \emph{new relationships} among pieces of data.
\end{quote}

Meaning is expressed in the \acl{sw} as sets of triples.
Each triple is composed by a subject, a predicate and an object like in a normal sentence (see Figure~\ref{fig:triples_example}).
Therefore, in  \citeauthor{berners-lee_semantic_2001} words, \emph{a document can assert that things (people, web pages or whatever) have properties (such as "is a sister of," "is the author of") with certain values (another person, another Web page)}.
A key difference with a normal sentence is that each concept is unambiguously defined by an URI.
These URIs form links between different triples as shown by the example of the Figure~\ref{fig:triples_example}.
% This set of triples can be expressed in multiple formats such as RDF, Turtle, Notation3 or NTriples. % TODO citas


\InsertFig{triples_example}{fig:triples_example}{
  Sample triples using different ontologies.
}{
  All the triples are represented graphically and some of them also textually.
  They describe academic and personal details about the author of this thesis.
  The knowledge is expressed using four different ontologies: FOAF, DC, SWRC and CiTE. % citar?
  Note that we use aliases, known as prefixes in most Semantic formats, to shorten some URIs and enhance the clarity of the Figure.
}{1}{}


A problem with the information described so far is that two different databases may use different URIs to express the same concept.
To overcome this the \ac{sw} offers collections of information called ontologies.
An ontology is a document which expresses the relations between terms commonly using a taxonomy and a set of rules to infer content.
The taxonomy defines the classes of a given domain and how they relate with each other.

Of course, different content providers could provide similar data described according to different ontologies.
Fortunately, ontologies can be easily mapped by providing equivalence relations within them.
In the same way, an ontology can be extended to adapt it to different application domains.
In any case, the reuse of the same models is beneficial and is therefore promoted through standardization.

% Remarkably, the use of the \ac{sw} has been promoted in the last years with the Linked Open Data (LOD) initiative.
% The LOD are datasets which follow a series of principles on how to open and publish data.
% The ultimate goal of the LOD is to publish linked terms using full semantics.
\subsection{Conclusions} % or Discussion???


In Section~\ref{sec:remote_invocation} we have analyzed the current trends towards the remote invocation style.
Within them, we can see how \ac{wot} has gained much more attention both for its perceived simplicity and for its seamless integration with the web.
% o decir REST based solutions para no centrar mucho en IoT
Besides, the availability of tools and libraries for many mobile and embedded platforms ease its adoption.


However, as any remote invocation style it introduces coupling between senders and receivers.
% de johanson_extending_2004 he sacado esas características, no la particularización a WoT
These complicates the application changes both in the short-term and in the long-term \citep{johanson_extending_2004}.
In the short-term because nodes constantly join and leave the environment due to mobility or to failures.
In the long-term because the space is used to solve new problems and the obsolete devices are replaced with new technology.
Furthermore, the human interaction is eased by minimizing the configurations.
% quizá lo de johanson_extending_2004 se podría retomar en el capítulo cha:tsc
In Chapter~\ref{cha:tsc}, we deeply compare both \ac{rest} and \ac{ts} to see how they can benefit each other.


\bigskip


% comparación de todas ellas
% de entre los que ofrecen distribución: este no mola por esto, esto otro mola por esto otro, etc.

In Section~\ref{sec:remote_invocation} we have studied the \ac{ts} solutions for \acl{ubicomp}.
Within these solutions, \emph{The event heap} is not considered because of its centralized nature.
The rest present different ways to disseminate the tuples. %  L2imbo, LIME and TOTA

The replication which L$^2$imbo proposes may not be feasible for devices with constrained memory or storage capacity.
Particularly, in spaces populated by many devices like the Internet of Things a large amount of tuples can be generated.
% need to configure the network

TOTA is useful for spatially connected \emph{ad hoc} environments, but it does not seem to suit well on networks where all the nodes can reach each other.
In this case, all the tuples could be replicated on the rest of the nodes showing the same limitations as L$^2$imbo.
% TODO COMPROBAR LEYENDO CON ATENCION LO DE TOTA!!!
Furthermore, \emph{ad hoc} environments are not the focus of our solution since we assume the connectivity between all the devices in a space.

% TODO no darle tan duro que si no a ver como luego me comparo con ellos :-)
Finally, despite of their unrealistic assumptions LIME proposes a model where each tuple is responsible of its own part of the space.
They share spaces whenever they become available.


\begin{table}%[position specifier]
  \centering
  \begin{tabular}{c|cccc}%p{5cm}}
      ~ & The event heap & L2imbo & LIME & TOTA \\
      \hline
      \hline
      Distributed & No & Yes & Yes & Yes \\
      Goal & - & Availability & Federation & Ad-hoc \\
      % Scalability?
      % Coger propiedades del libro de sistemas distribuidos?
  \end{tabular}
  \caption{Comparison of the most prominent \acl{ts} alternatives for \acl{ubicomp}.} % no columpiarme, o poner todas o ninguna?
  \label{tab:ubicomp_ts_comparison}
\end{table}


\medskip

In our model, we argue that in an \ac{ubicomp} environment each node should manage their own semantic information.
By delegating responsibility 
we naturally represent mobility scenarios where users carry their own profiles in their mobiles to new spaces 
and 
we acknowledge the fact that embedded devices directly control their own actuators and sensors.
Therefore, we propose to directly access to the source of the data avoiding the use of intermediaries whenever it is possible.
In this aspect, our solution has more in common with LIME.
Two major differences with LIME are the use of semantics and the natural integration with the web. % i.e. use of REST/HTTP
\section{\acl{ubicomp}}
\label{sec:soa_ubicomp}

The previous section presented the main research topics related to this dissertation.
% The middleware presented in this thesis adopts some relevant characteristics from them.
From their intersection with the \ac{ubicomp} field emerges relevant related work.
Consequently, this section describes these works and presents their differences and similarities with our work.
% hablar de comparacion de ventajas en inconvenientes? hablar de contribuciones a dichas areas?

% presentar secciones
% we will scrutinize how other authors have use the \ac{sw}.


\subsection{The \acl{wot} and other \ac{rest} Solutions for \acs{ubicomp}}
% nueva sección para hablar de WoT (DPWS ya se ha presentado)

\InsertFig{venn-sec1_1}{fig:venn_ubicomp}{
  Non-semantic integration approaches for \ac{ubicomp} apart from \aclp{ts}.
}{
  Scope of this subsection.
}{0.6}{}


The \acl{wot} initiative encourages the use of \acs{rest}-based solutions embedding web servers in daily objects \citep{guinard_internet_2011}.
In this way, the objects can integrate with the \ac{www} as first-class citizens. % decir ``f-c cit of the Web" es reiterativo ya que WWW==web
This integration brings the following benefits:
\begin{itemize}
  \item The smart-things can be linked to enable its discovery by browsing. This involves using the tool most users are familiar with: the browser.
  \item They can be bookmarked or shared through social networks \citep{guinard_sharing_2010}.
  % explicar qué es un mashup?
  \item They can be integrated with other web applications through mash-ups \citep{guinard_towards_2009,ostermaier_webplug:_2010,pintus_anatomy_2011}.
  % TODO citar al canadiense tb? pintus, stirbu y a todo el mundo
  \item Mechanisms such as searching, caching, load-balancing and indexing can be used over the objects to achieve the scalability of the web. % cita al indio
\end{itemize}


% lo he comentado antes, no sé si no me estoy repitiendo más que el ajo...
But besides the functional differences, one of the major advantages of REST compared with the WS-* is that is perceived as simpler.
\citet{guinard_search_2011} proved that developers consider it easier to learn and find it more suitable for programming smart things.
This makes \ac{wot} solutions more popular for resource constrained devices than the WS-* based ones.

% Mencionar un par de trabajos sobre móviles con servidores web embebidos, para que se vea que no es sólo \ac{wot} % WoT (REST + ubicomp)
% Ignorar esta sección? Es realmente necesaria cuando ya he puesto la comparativa de TSC?
\subsection{\acl{ts} for \ac{ubicomp}}

So far, \acl{ts} has been adapted to \acl{ubicomp} by different authors.


\InsertFig{venn-sec1_2}{fig:venn_tuplespaces_ubicomp}{
  \aclp{ts} solutions for Ubicomp.
}{
  Scope of this subsection.
}{0.6}{}


% TODO leer artículos que me faltan y corregir
The \emph{event heap} \citep{johanson_extending_2004} is a system used for a specific \ac{ubicomp} sub-domain: interactive workspaces.
In this scenario there are rooms with different devices deployed and where mobile devices can enter.
Each room has its own space where the devices exchanging tuples to cooperate.
%For example, a video can be presented in a display and through a remote controller, the user can place the tuple for pausing it.
%The display consumes these kind of tuples, so it pauses the video whenever somebody places that tuple on the space.
% This system is simple to implement but is limited since it centralizes the space in a single machine per room.
This work merely identifies the requirements of these environments and the properties which solve them.
Then, it discusses how these properties can be satisfied using \ac{ts} or some extensions.
Finally they compare its implementation both with other \acp{ts} implementations or other coordination infrastructures. % e.g. RMI, MOM, Pub/sub

Although the analysis presented by \citeauthor{johanson_extending_2004} identifies the key requirements for \ac{ts} in \ac{ubicomp},
the solution they propose is purely centralized.
We opted for a distributed solution at the expense of limiting some of the features.
% indicar cuales son o pasar?
% usar sus características a la hora de valorar nuestra solución TSC en futuras secciones?


\emph{L$^2$imbo} \citep{davies_l2imbo:_1998,friday_experiences_1999} replicates the tuples to avoid a single point of failure.
Each node joined to a space uses an IP multicast address to exchange messages with other nodes in that space.
Writing into a space involves sending a multicast message to inform to the rest of the nodes of the tuple written.
Reading operation usually requires local reading.
Destructive reading of the tuple is more complex as it requires a global withdrawal.
In \emph{L$^2$imbo} only the owner of a tuple can remove it from the space.
The ownership of a tuple initially belongs to its creator, but can be transferred. % TODO ver qué añadió Friday


 % meter mas citas de LIME que estas 2?
\emph{LIME} (Linda in a Mobile Environment) \citep{picco_lime:_1999} is a TS solution for mobile systems.
In LIME each mobile device has its own space where it generally writes its tuples.
This space is shared with other devices creating federated spaces, i.e. the aggregation of different shared spaces.
In this way, each mobile can access to tuples in other mobiles whenever they become available.
They also proposed a new writing primitive to insert tuples in remote spaces.
% ejemplo de service discovery?

This is difficult to implement and \citet{coulouris_distributed_2012} complain about the unrealistic assumptions they make to simplify the problem. % cita al libro
These assumptions are the uniform multicast connectivity between devices whose tuple spaces are aggregated and the serialized and ordered connections and disconnections.
In any case, LIME was adapted to several platforms to run both in embedded and mobile devices \citep{murphy_transiently_2006}.


In the \emph{TOTA} (Tuples On The Air) Project \citep{mamei_programming_2009} tuples are disseminated to different devices.
To that end, each tuple has 3 fields:
1) the content of the tuple,
2) a rule which defines how it should be propagated,
3) a rule to define its maintenance.
For instance, they consider a museum where a visitor writes a query tuple describing a piece of art he wants to see.
The propagation rule defines that it should be propagated to all nodes in the vicinity, increasing the distance by one each time.
The tuple is configured to be deleted after a time-to-live period using its maintenance rule.
When the it reaches the room where the piece of art is located, the piece of art will write a response tuple.
This response tuple jumps from a device to another until it reaches the device which queried for it.


% TODO Analizar otras como TOTAM o CRIME % TS para ubicomp
% TODO
%  + reescribirlo como una misma cosa?
%  + añadirle una 

\subsection{\acs{sw} by using Intermediaries}
\label{sec:sw_intermediaries}

% como los smart environments describen contexto usando web semántica

% uso concreto por parte de soluciones significativas: siempre centralizando el uso de semántica en cacharros grandes

% intro a que ahora se va a hablar de soluciones IoT que usen semántica

Adding semantics works well for devices with high computational capacity but may add too much overhead for most of the devices in the \ac{iot}.
To reduce this overhead in such devices, part of this computation is usually delegated to an intermediary.
Some noteworthy example is the one proposed by \citet{broring_semantic_2009}.
% buscar otros ejemplos de enjundia
% The Context Broker Architecture (CoBrA)[Che04] 
% Gu et al. (2007)
% AlarmNet (2008)


\InsertFig{venn-sec2}{fig:venn_semantics_ubicomp}{
  The \acl{sw} for \acl{ubicomp}.
}{
  Scope of this section.
}{0.6}{}


These intermediaries or \emph{Semantic Gateways} are in charge of managing the semantic annotation.
The devices send raw data (which can be compressed) to the intermediaries and the gateways annotate the content semantically.
Thus, the devices do not have to care about any semantic aspect and just collect the data as they did before.

These \emph{Semantic Gateways} reduce the load to embedded devices with limited resources by decreasing the number of requests they have to provide.
In addition, a centralized intermediary can gather all the information and thus, reduce the complexity of managing a distributed environment.

However, using intermediaries to store the semantic data of resource constrained devices also has some drawbacks.
On the one hand, centralization does not faithfully represent mobility situations were individuals carry their own semantic information in their personal devices.
In addition, the data obtained from an intermediary will always be less fresh than the one obtained where it is generated (i.e., sensors).
On the other hand, the servers are critical in centralized systems and therefore, their availability determines the operation of these solutions.
They also impose a burden on the maintenance which may be worthless in some simple scenarios.


% TODO meter alguna referencia a Cabilmonte sobre SPARQL en streaming
% Trata de consultar datos semánticos en streaming a una base de datos relacional.
% Por lo que entiendo, la comunicación sensor-DB se no es semántica.



% diferenciar de una forma menos descarada?

\subsection{\acs{sw} in Providers}
\label{sec:sw_providers}
% enumerar aquellas características de IoT???
% sacar algo del paper que habla de los retos de usar semántica en IoT

% pero ahora los cacharros cada vez son más potentes y no es difícil imaginar un mundo poblado por ellos blablah

% WoT y aquel que hacía cosas de móviles

% TODO generalizar a resource-constrained devices
Lately some solutions have arisen to semantically annotate data where it is generated. % in the provider
In the \acl{wot}, multiple solutions have considered using semantics to enrich the data definition in a machine processable manner.
Generally, these solutions embed the metadata in HTML using microdata, microformats or RDFa \citep{mayer_extensible_2011}.
These contents are returned by the Internet connected objects and are used to enhance the findability of the data by search engines.
% Ontologías no se pueden definir, tiene que haber consenso en su uso por parte de las máquinas de búsqueda.
% en vez de eso, queremos mejorar la busqueda por parte de los cacharros, no de los search engines
This approach does not focus on making the devices able to search and interact with others.
% end-to-end search
Instead, it considers them as mere human-oriented information providers which should be indexed by third party search services.
% aquí se puede enganchar también el discurso de Doulkeridis2007desent: coverage and scalability, freshness y monopoly

A more general way to represent semantic data is using RDF\footnote{\url{http://www.w3.org/RDF/}} based representations (i.e., full semantics).
SPITFIRE European project\footnote{\url{http://spitfire-project.eu}} represents the most remarkable effort on gathering full semantics and the \ac{wot}.
It focuses on fully integrating sensor data with the \ac{lod}. 
The \ac{lod} are datasets which follow a series of principles on how to open and publish data.
The goal of the \ac{lod} is to publish linked terms using full semantics.

SPITFIRE shares with our solution the vision of a world populated by devices acting as semantic data providers no matter how small they are \citep{hasemann_rdf_2012}.
Therefore, many of their efforts are complimentary to this work.
%This is done by defining an ontology for mapping other common ontologies and providing semi-automatic generation of semantic annotations from raw data.
%They also propose an abstraction to represent real-world entities (virtual sensors) using data provided by low-level sensors.
%Finally, and more specific to this work, they propose a searching model which predicts the current state of things by computing their periodic patterns in past states.
To search within these providers, \citet{pfisterer_spitfire:_2011} proposes a model which predicts the current state of things by computing their periodic patterns in past states.
Again, the goal of this method is to adapt search engines to the new fashion of data provided on the \emph{\acl{sw} of Things}.

Instead, we propose a model to enhance the search capability of web-connected things in a distributed fashion.
We aim to promote the interaction and collaboration between any devices in these environments.
Not only sensors, but also regular computers or mobile devices \citep{balandin_access_2011}. % semántica para ubicomp
%\subsection{Conclusions} % or Discussion???

Section~\ref{sec:sw_intermediaries} presents the most common approaches to use \ac{sw} in \ac{ubicomp}: delegate on intermediaries the semantization of the raw data.
However, we do not intend to follow this path.
Instead, we propose to share semantic content directly from the embedded and mobile platforms as the solutions in Section~\ref{sec:sw_providers} do.
This completely distributed approach tries to simplify the maintenance burden and promote the access to the most updated data.


Regarding searching, the solutions on Section~\ref{sec:sw_providers} promote a less distributed search.
Instead, our model tries to enhance searching capability of web-connected things.
This improvement leads to environments where any devices directly interact and collaborate with each other.
\section{Semantic \aclp{ts}}
\label{sec:tsc_soa}

% Añadir  Citron (2005)? (eoa de Almeida)
% Tanasescu
% swarm
% decir que tiene mucho que ver con distributed triple stores
% analizarlos también??? vaya pereza! -> igual comentar diferencia y dejar el análisis para la sección tosca de IGWS

Semantic \aclp{ts} aim to join \aclp{ts} with the \acl{sw} to propose a more uncoupled solution.
Particularly, it benefits from the autonomies introduced by both \ac{ts} and the \ac{sw}:

\begin{itemize}
  \item Space uncoupling.
  \item Time uncoupling.
  \item Data-schema uncoupling. % explicar?
% 	enhances the interoperability of the
% 	applications built on top of TSC. This way, two applications using standard ontologies can interact among them
% 	automatically enriching one each other, as long as they use the same space and standard and linked ontologies.
\end{itemize}


In this section we provide a review of the existing Semantic \aclp{ts} solutions. % citar nixon_tuplespace-based_2008 ???
To summarize, we present a comparison of them at the end of the section, analyzing their limitations.


\InsertFig{venn-sec3}{fig:venn_tuplespaces_semantics}{
  Semantic \aclp{ts}.
}{
  Scope of this section.
}{0.6}{}


\subsection{TSC}
% modelo propuesto
TSC \citep{fensel_tsc_2007} is the middleware which first implemented the ideas of Fensel et al. \citep{fensel_triple-space_2004}.
TSC explores a \acl{tsc} infrastructure based on the \ac{rest} architectural style.
Equivalently to \ac{rest}, TSC proposes a stateless infrastructure where the resources are identified by \acp{uri}.
In the first approximation, each tuple was considered a resource.
These tuples correspond with \ac{rdf} triples.
In a later work,  Krummenacher et al. \citep{krummenacher2006specification} proposed to use identifiers just for a set of triples (i.e. \ac{rdf} graphs).
% bundler cita el de "WWW: or what's wrong with the web" para explicar eso último
They argue that resources usually need to be modeled using more than just a triple.
Besides, the use of graphs reduces the complexity on storage.
To cope with the heterogeneity of data and scalability issues, they propose the use of independent spaces identified by their own \ac{uri}.
Using these disjoint spaces, they can restrict the number of participant in each one.


% cosas relativas al despliegue
TSC combines client/server and \ac{p2p} architectures using a hybrid super-peer one.
The super-peers of this architecture communicate with each other using a middleware which offers a virtual shared memory space of Java objects.
They present 3 different types of nodes:
\begin{itemize}
  \item The \emph{servers} are the responsible for storing and replicating the spaces.
  \item The \emph{heavy clients} are not always connected to the system, but they have the same responsibilities as the \emph{servers}.
	They can work off-line with their replicas.
  \item The \emph{light clients} can access to the spaces using proxies or an \ac{http} \acs{api}.
\end{itemize}


% arquitectura
TSC consist of three main components: a mediation engine, a data access layer and a coordination layer.
The mediation engine manages data heterogeneity by intermediating between nodes using different data schemas.
The data access layer ensures persistence of data, supports template matching and provides some reasoning over the knowledge stored.
The coordination layer replicates named graphs to all involved TS kernels and therefore has to be synchronized with the data access.
% http://tsc.deri.at/deliverables/D21.html
% http://tsc.deri.at/deliverables/D31/D31.html
To do that, they use a \ac{p2p} implementation of a virtual shared data space. % cita a CORSO?
% decir que ofrece transaccionalidad gracias a CORSO? transactionality
% It also offers a transactional context and a simple form to publish and subscribe to certain patterns.



\subsection{Semantic Web Spaces}

% modelo propuesto
Semantic Web Spaces (SWS) \citep{tolksdorf_coordination_2006} defines two different data coordination views: data view and information view.
In both views the tuples are represented with three fields which correspond with \ac{rdf} triples' subject, predicate and object.
The data view stores syntactically valid \ac{rdf} and it is accessed using Linda primitives.
The information view stores consistent and satisfiable data which are managed using new primitives.
This view takes into account the knowledge defined by ontologies to perform semantic matching over inferred triples.

% cosas relativas al despliegue
SWS are virtually partitioned using contexts, which provide a simple form of access control.
Clients can allocate contexts to control their view privately or share it with other clients. % This allow clustering ¿WTF?
The different active spaces and the tuples they contain are represented in a meta-space using a meta-model ontology.
This ontology also allows to express the hierarchical relation between the spaces and security policies.
% por lo que entiendo, está centralizadísimo

% arquitectura
The prototype contains three main components: publication, retrieval and security.
The publication component describes how the platform communicates with the clients or storage managers.
The retrieval component process queries and can reason to enhance results. % no se razona cada vez que se inserta una tupla
The security component restrict the access of client to spaces according to the policies defined.  % hablan de agentes
% dicen que experimentaron problemas de escalabilidad



\subsection{sTuples}

% modelo propuesto
STuples \citep{khushraj_stuples:_2004} was conceived by Nokia Research Center as a pervasive computing work.
It extends the JavaSpace middleware, a object-oriented \acl{ts}, using a semantic match over the semantic tuples.
To do that, a semantic tuple is an object with a field with the semantic description in the DAML+OIL format. % definir DAML+OIL o citarlo?


% cosas relativas al despliegue
The space is centralized and uses a framework called Vigil to create clients and services. % cita Vigil?
Vigil acts as a communication gateway with several protocols and it is directly integrated with sTuples.


% arquitectura
The JavaSpace extension is done in sTuples through two components: the Tuple Manager and the Tuple Matcher.
The manager adds, removes or states changes of tuples.
Every insertion requires the semantic data to be asserted.
The matcher implements the algorithm to match templates.
Besides, specialized agents reside on the space to offer added functionality to the users.
For instance, they implemented recommendations, subscriptions and task-executions in their prototype.



\subsection{Conceptual Spaces}

% modelo propuesto
Conceptual Spaces, or CSpaces, were born to study the applicability of semantic \aclp{ts} to different scenarios including \acl{ubicomp}.
Tuples in CSpaces contain seven different fields.
Some fields enable to express logic theories using different formal languages.
Other remarkable fields are used to define the version, specify the creator or to define to which subspaces the tuple belongs. % explicar los 7 campos?

% cosas relativas al despliegue
The spaces in CSpaces can be individual or shared.
An individual space belongs to a single process.
A shared one allows access to all the processes who have agreed on.

% arquitectura
%Es una IDEA y NO la habían IMPLEMENTADO (posteriormente hicieron un prototipo limitado para un paper de 2006).
CSpaces differentiates between the same three types of nodes as TSC.
They intended to implement individual spaces in heavy-clients and shared spaces in servers.
They first considered a centralized implementation for the shared-spaces, % lo dicen en la comparativa
and later on a P2P network of the nodes managing a shared-space. % paper de después
In any case, to the best of our knowledge, this conceptual exercise never went beyond a rather limited prototype.
% TODO confirmar 100%



\subsection{tsc++}

% modelo propuesto
Tsc++ \citep{krummenacher_open_2009,blunder_distributed_2009} is a new version of the TSC project.
Therefore, although the model is equivalent to TSC, its implementation completely changes.


% cosas relativas al despliegue
% Páginas 54 y 55
Instead of replicating each space in the different kernels as TSC does, tsc++ proposes to divide the spaces among the participants.
Therefore, a space corresponds with a group a nodes on the \emph{Jxta} \ac{p2p} framework \footnote{https://jxta.dev.java.net/}. % jxta de otra forma?
These nodes write their information locally and discover information in other nodes using some special nodes.
This special nodes in Jxta are called \emph{rendezvous} (RDV).
Tsc++ proposes and evaluates two discovery algorithms: flooding and random walk (RDW).
\begin{itemize}
  \item The flooding algorithm sends the request to a RDV, which propagates it to all the known nodes.
	If this nodes are RDV, they may propagate the request until a number of maximum hops is reached.
  \item The RDW algorithm chooses one peer to send the message to after the request has been processed in the current one.
	In Tsc++ the request is sent to a RDV which creates a list of all the nodes and sends it to the first one in the list.
	At each step a new node from that list will be chosen.
	Besides, the \emph{walk} ends when a result is returned or when a maximum number of hops is reached if no response was found.
\end{itemize}
% De dónde leches saqué esto de NB?
%These nodes follow a simple flooding-based strategy to distribute their information called \emph{negative broadcasting}.
%In negative broadcasting each node writes the information locally and read querying to the rest of the nodes of the space.


% arquitectura
Tsc++ has 3 main components: the communication and the data access.
The data access layer relies on a semantic repository to perform queries and store the triples.
% mencionar que son Sesame \citep{broekstra_sesame:_2002} y Owlim \citep{kiryakov_owlimpragmatic_2005} ???
The communication layer maps the reading coordination primitives to group communication.
% No evaluan el mecanismo de subscripción porque dicen que se basa en DHT y que como es escalable no necesita más evaluación
% No me queda claro de qué forma usan DHT sobre los nodos RDV que son los que miran las subscripciones cómo están.
Besides, it offers a subscription mechanism independent to the content stored in the space which uses triples and templates as a basis.


% it does not make inference
% it does not allow expressive querying
% it has not been designed for devices with reduced computing capabilities, because tsc++ middleware focused on architecture and implementation in large scale and we focus in the short scale (local area networks with an intelligent environment).



\subsection{Triple Space Communication (TripCom)}

% modelo propuesto
TripCom\footnote{TripCom (IST-4-027324-STP, www.tripcom.org)} is a Triple Space solution which works with \ac{rdf} triples as resources.
The most remarkable aspects of the model proposed by TripCom are:
\begin{itemize}
  \item The use of subspaces to form form nested multiple spaces.
	Using them it restricts the communication to a part of the whole space.
	This restriction leads to local scalability and completeness.
  \item A mechanism to overlap spaces called \emph{scopes}.
	Using the scopes a client can have a temporary copy of some tuples.
	However, any insertion and deletion would not apply to the whole Triple Space.
  \item A meta-space to describe the relations between spaces or to which graph a tuple belongs.
  \item Different types of templates can be used to query the space, from simple wildcard-based templates to more expressive SPARQL queries.
  \item New primitives: read at triple or a graph level and writing multiple triples at once.
  \item A subscription mechanism.
\end{itemize}


% cosas relativas al despliegue
To access to a space a client needs to contact a TripCom instance, i.e. kernel.
% TODO Duda inglés, está bien hecho eso de "sujeto verbo objeto and verbo" o debería ser "sujeto verbo objeto and sujeto verbo objeto"???
% TODO Uniformizarlo a lo largo del texto
Each kernel stores one or more subspaces and can contact other kernels responsible for different spaces.
To do that, if the space \ac{url} is provided, it simply resolves this \ac{url} using \ac{dns} to contact the other kernel.
Otherwise, the kernel uses 4 different strategies:
\begin{itemize}
  \item Triple Provider. It uses shortcuts to know who answered a query in the past.
  \item Recommender. It uses shortcuts to know which kernel successfully routed a query in the past.
  \item Indexing - \ac{dht}.
	It creates indexes using a hash function over the subject, predicate, object and space \ac{url}.
	Then, it stores these indexes in a distributed database which relies in a structured \ac{p2p} system. % mencionar PGrid???
\end{itemize}


% arquitectura
A TripCom kernel has several components which communicate with each other using a \acl{ts} solution.
The most important components regarding how the triples are distributed are:
\begin{itemize}
  \item Triple Store Adapter.
	It acts as a gateway with the selected repository.
  \item Query Processor. It descomposes query into subqueries.
	Then, to answer each subquery the kernel: a) queries the local data store or b) forwards it to other kernel.
  \item Distribution Manager.
	Forwards queries or write requests to the appropiate kernels.
\end{itemize}
Apart from these, there are also security, mediation, metadata and transaction managers.
On top of these components, TripCom offers three different \acp{api}:
Triple Space \acs{api} to map the primitives,
web service \acs{api}, which offers higher level service consumption \acs{api}
and management \acs{api} to configure the each kernel.



\subsection{Smart-M3}

% modelo propuesto
Smart-M3 \citep{honkola_smart-m3_2010} is an open-source distributed middleware which uses \ac{rdf} graphs.
It supports three types of primitives: \acl{ts}-based primitives, subscriptions and a query over a whole space.
Smart-M3 defines two types of nodes: Semantic Information Brokers (SIB) and Knowledge Processors (KP).
The SIBs store the knowledge and perform reasoning process on behalf of the KPs.
A KP access to a SIB using two different types of interfaces (called KPI):
\begin{itemize}
  \item The low-level API uses triples as a base.
  \item The high-level API allows the developer to use ontological entities such as classes, relations and individuals.
\end{itemize}


% cosas relativas al despliegue
The space can be distributed over more than a SIB using a \emph{distributed deductive closure}.
 % TODO asegurarse 100% de esto que voy a decir!
However, to the best of our knowledge no evaluation with more than a SIB has been presented yet, making it de facto an centralized space.
For the communication between SIBs and KPs, Smart-M3 defines a protocol called \emph{Smart Access Protocol (SSAP)}. % decir que está basado en XML?
SSAP can be implemented on top of different communication mechanisms (e.g. SOA, XMPP, Bluetooth or TCP/IP) making the solution communication agnostic. % TODO acrónimos?


% arquitectura
% TODO repasar cómo está escrito, porque con esos nombres queda un poco confuso
A SIB is divided in 5 layers: transport, request handling, graph operations, triple operation and persistent store.
The transport layer has different processes for each communication mechanism supported.
Once the SSAP operations have been extracted, the request handling creates a thread for each primitive.
The \emph{graph operations} layer works with primitives which use \acs{rdf} graphs and it uses the underlying \emph{triple operations}. % ¿?
In the \emph{triple operation}, expressive queries may need to be translated to query the store.
% decir que también soporta subscripción?

% dibujo?

% De: http://www.diem.fi/files/KP_reference%20implementation.pdf
%  The Semantic Information Broker is the information repository of the Smart 
%  Environment. In theory, the Smart-M3 Smart Space can consist of one or more SIB 
%  entities. However, the SIB reference implementation does not currently support any 
%  kind of interaction between different SIBs, which would enable distribution of the 
%  Smart Space. On the IOP, the SIB is implemented as a NoTA SN.



\subsection{Nardini et al.}

% modelo propuesto
\citet{nardini_semantic_2013} explores the use of Description Logics on top of an existing \emph{tuple centre} implementation called TuCSoN. % TODO CITA TuCSoN???
A tuple centre is a tuple space where the behavior can be programmed as a result of the basic coordination primitives.
TuCSoN aims to provide \emph{tuple centres} that agents can use through Linda-type primitives.
This interaction can cause reactions which trigger behavior previously implemented in the \emph{tuple centre}.

% cosas relativas al despliegue
In TuCSoN each \emph{tuple centre} is deployed in device's port and it is identified by a name.
The communication with each centre is direct, not based in any underlying network. % no les mola P2P decian en algun lado
The agents can access to remote \emph{tuple centres} and migrate to them, but these interactions need to be programmed.

% arquitectura
Nardini et al. defines a domain ontology attached to each centre to provide semantic interpretation to the tuples stored.
They extend TuCSoN providing semantic tuples, templates, primitives, reactions and matching mechanism.
\begin{itemize}
  \item The semantic tuples express the logic terms.
  \item The semantic template defines a language based on logic.
  \item To discriminate common syntactic primitives from semantic ones they use a keyword.
  \item The semantic reactions redefine events, guards and reactions enabling the use of first-order-logic.
	The events define when a reaction should be triggered.
	The reaction guard represent a set of conditions which need to be satisfied to execute a reaction.
  \item For the semantic matching mechanism they integrate TuCSoN with a reasoner (Pellet). % citar pellet?
\end{itemize}



\subsection{Comparison}
% TODO CSpaces no hace más que liar el analisis, me planteo no incluirlo
%		Porque es muy bonito poner ideas felices, pero no ha sido implementado nunca!


\subsubsection{\acl{ts} model}

A \acl{ts} model is mainly characterized by the information hold (i.e. \emph{tuple}),
how it is hold (i.e. the \emph{space} model)
and how a node can query the space (i.e. the type of \emph{query}).
The Table~\ref{tab:comparisonTS} summarizes the model used in each of the semantic \aclp{ts}.

The \emph{flat} model is the simplest manner of disjoining spaces.
Identifying this spaces with URIs, one can link them in the knowledge similarly to what other tuple models do (e.g. Nardini et al or CSpaces). % comprobar
More sophisticated models propose virtual view of spaces (i.e. scopes) or hierarchies between spaces.
Nevertheless, we argue that this models are not easy to implement in a distributed way. % extenderlo despues?

The most common tuple model is the three-field tuple with \emph{subject}, \emph{predicate} and \emph{object}.
Other solutions propose their own tuple format.
We believe that for a developer used to work with the Semantic Web the use of \ac{rdf} triples is straightforward.
Grouping these triples in \ac{rdf} graphs also provides a useful abstraction for resources.
% TODO aclarar bien porqu TSC se centra en grafos y tsc++ no!

Regarding the queries, all the solutions offer graph patterns-based queries.
These queries can be easily processed by any device, whereas complex queries need parsers.
These parsers are not always available for resource constrained platforms.
Therefore, in this thesis we focus on queries based on graph patterns.
The adoption of more complex querying languages such as SPARQL can be done as a future work. % cita SPARQL?




\begin{table}[htbp]
\caption{Tuple Space model comparison for the analyzed solutions.}

\begin{tabular}{ l p{3cm} p{3cm} p{3cm} }
\hline 
  & Space  & Tuples  & Queries \tabularnewline
\hline 
 TSC & Flat  & Named graphs  & Graph patterns \tabularnewline
 SWS & Disjoint nested  & Triples  & G. patterns and complex \tabularnewline
 sTuples & Centralized nested  & Tuples with DAML+OIL field  & DL-based \tabularnewline
 CSpaces & Interconnected  & 7-field tuple  & Formal lenguage \tabularnewline
 tsc++ & Flat  & Triples  & Queries \tabularnewline
 TripCom & Nested  & Triples  & G. patterns and complex \tabularnewline
 Smart-M3 & Flat  & Triples  & G. patterns and complex \tabularnewline
 Nardini et al. & Tuple centres managed in different nodes & Tuples based on formal lenguage  & Formal lenguage \tabularnewline
\hline 
\end{tabular}
\label{tab:comparisonTS}
\end{table}



\subsubsection{Additional features}

The Table~\ref{tab:compAdds} summarizes the additional mechanisms provided by the solutions analyzed.
Most of them offer subscription mechanisms which are particularly useful for distributed environments.
% TODO modificar esto que ahora no hay implementacion!
Although this dissertation focuses more in the querying mechanism, in the Chapter~\ref{cha:implementation} we will propose a subscription mechanism.
Other solutions offer , but we have discarded its usage due to the hard implementation for distributed systems running on dynamic environments.


% transacciones
% subscripciones
% etc.

\begin{table}[htbp]
\caption{Features offered by the analyzed solutions.}

\begin{tabular}{ l p{3cm} l p{4cm} }
\hline 
  & Subscriptions  & Transactions  & Other features \tabularnewline
\hline 
 TSC & Yes  & Yes  &  \tabularnewline
 SWS & No  & No  &  \tabularnewline
 sTuples & Yes  & No  &  \tabularnewline
 CSpaces & Yes  & No  & Multiple read and writes \tabularnewline % bussler dice que tiene algún tipo de transaccionalidad, yo no lo creo
 tsc++ & Yes  & No  &  \tabularnewline
 TripCom & Yes  & Yes  &  \tabularnewline
 Smart-M3 & Yes  & No  &  \tabularnewline
 Nardini et al. & No \footnote{But it has been implemented in TuCSoN before \cite{ricci_extending_2002}.} & No  &  \tabularnewline
\hline 
\end{tabular}
\label{tab:compAdds} 
\end{table}





\subsubsection{Use of semantics}
% TODO evitar verengenales que no comprendo y que sólo me complican la vida

The Table~\ref{tab:comparisonSemantics} shows different features regarding the semantic model proposed by each solution.
First we can see how most of the solutions use the standard Semantic Web language \ac{rdf} to express knowledge.
The use of \ac{rdf} does not impose or limits the use of more advanced languages such as \ac{rdfs} or \ac{owl}.

CSpaces and Nardini et al. propose the use of a language-independent \ac{fol}.
However, since most advanced reasoners only implement a subset of \ac{fol}, any implementation is limited by them. % también estas
% theoretical counterpart of OWL DL [7]. Since OWL is the W3C standard ontology description language for the Semantic Web,
% and the standard de-facto for semantic applications in general, we adopt OWL as the ontology language for the domain
% ontologies associated to ReSpecT semantic tuple centres in TuCSoN. While for the details of OWL we forward interested

Reasoning is used both to validate the data inserted into the space or to query over inferred triples.
We do not validate the data inserted on the space to reduce the insertion time.
The possibility of querying over the inferred knowledge is limited by the availability of reasoners.
As will be shown in the next sections, nowadays there is no reasoner light enough to run in a reasonable time span in small devices.
Therefore, reasoning is not mandatory in our solution.
In any case, the use of reasoners over local content is tangential to the model proposed in this thesis.

Semantic matching is usually derived from the reasoning process.
Although it is a desirable characteristic, it is not imposed by any mean.
For our solution, describing the data semantically is a first step towards semantically interpreting it.




\begin{table}[htbp]
\caption{Use of semantics in the analyzed solutions.}

% TODO En Languages decir si se basa en RDF o FOL, el hecho de que luego se use OWL para definir conocimiento, se sobreentiende
% TODO añadir reglas?
% TODO añadir Validity and consistency?
% TODO añadir semantic clustering

\begin{tabular}{ l p{3cm} p{3cm} p{3cm} }
\hline 
  & Languages  & Reasoning  & Semantic matching \tabularnewline
\hline 
 TSC & RDF  & -  & Graph pattern templates and N3QL resolution \tabularnewline
 SWS & RDF(s) extendable to OWL, SWRL  & Yes (to match and validate)  & Subsumption-based \tabularnewline
 sTuples & DAML+OIL  & Tuples with DAML+OIL field  & Subsumption-based \tabularnewline
 CSpaces & Language-independent, up to FOL  & Yes (used for query answering, rewriting and consistency checking)  & Based on query engines \tabularnewline
 tsc++ & RDF & No  & No \tabularnewline % comprobar %  lo de para query answering lo hice yo en su día
 TripCom & RDF  & Yes (query rewriting) & Yes \tabularnewline % TODO rehacer esta columna comprobandolo
 Smart-M3 & RDF  & ?  & G. patterns and complex \tabularnewline
 Nardini et al. & Language-independent, up to FOL & ?  & ? \tabularnewline
\hline 
\end{tabular}
\label{tab:comparisonSemantics}
\end{table}

% TODO añadir tabla con análisis de las autonomías?



\subsubsection{Distribution}
%	2. centralizado / distribuído
%	4. basado en clientes tontos

Centralized \aclp{ts} are much simpler and easier to implement.
Therefore, they usually offer more features than the distributed ones.
However, they also impose a single-point-of-failure.

This thesis aims to solve that by proposing a distributed model.
Within the distributed \aclp{ts}, a wide range of distribution strategies have been detected.
The Table~\ref{tab:comparisonDistribution} summarizes the space distribution strategy adopted by each solution.

TSC uses flooding to replicate the knowledge in all the nodes.
Tsc++ propagates the queries in flooding-based approaches.
TripCom is built on top of a structured-network which imposes where the content is stored.
The suitability of these models is subject of analysis in the Chapter~XX.

% TODO reescribir desde aquí

%In this thesis, we aim to implement a completely distributed Semantic \acl{ts}.
%In that space all the nodes should be able to participate trying to avoid any distinction between clients and servers.
%Ubicomp is be characterized for being populated by very dynamic nodes.
%This nodes usually face severe computation and energy autonomy restrictions.



% TODO y así te evitas engorrosas explicaciones aquí
%We argue that the replication is not suitable for devices with severe storage, memory and computation restrictions.
%Flooding-based queries does not scalable for spaces with many nodes, .
%Finally structured networks does not comply on

\begin{savenotes}
  \begin{table}[htbp]
    \caption{Distribution of the spaces.}
    \centering
    \begin{tabular}{ l c c p{4cm} }
    \hline 
    & C/S & Distributed & Distribution \\
    & access & space & strategy \\
    \hline 
    \midtsc{} & $\checkmark$ & $\checkmark$ & Replication \\ % Possitive broadcasting
    \midsws{} & $\checkmark$ & × & - \\
    \midstuples{} & $\checkmark$ & × & - \\
    \midcspaces{} & $\checkmark$ & $\checkmark$ & Not detailed \\ % además no implementado :-S
    \midtscpp{} & × & $\checkmark$  & Local writing, different query strategies \\ % Flooding, RW, etc.
    \midtripcom{} & $\checkmark$ & $\checkmark$ & Structured network, different strategies \\
    \midsmartmt{} & $\checkmark$ &  $\checkmark$ & Theoretical \\ % decidir si me convence que no haya sido implementado, TODO citar al tipo que dijo como hacer SIBs distribuidos
    \midnardini{} & $\checkmark$ & × & - \\
    \hline 
    \end{tabular}
    \label{tab:distribution_comparison}
  \end{table}
\end{savenotes}



\subsubsection{Implementations}
%	5. Implementado?
%	6. Implementación disponible
%	Interoperability
%	Multiplatform
%	8. dependencias: librerias, etc. (para decir que lo nuestro se puede implementar en casi cualquier plataforma de forma fácil)



\subsection{Conclusion}

% extraer algunas conclusiones de lo que falta
% tono general: están muy bien de características, pero a la hora de hacerlo distribuido falla
% hacen aguas en proponer un modelo que pueda ser usado más allá del client-servidor
% Suitability for Ubicomp

So far, semantic \aclp{ts} solutions have offered a sophisticated models with interesting additional features.
However, in most of the cases this is possible by delegating such tasks on a subset of more powerful and static nodes.
Therefore, the distribution of the space is restricted to some nodes when in the best case, or centralized in a unique node in the worst case.
In these cases, the rest of the nodes are head toward the use of the space as simple clients.


In this thesis we will explore the most decentralized case: the one in which each node manages its own information.
Doing so, a user can carry out his profile in the smart-phone or a sensor can provide the last measures on-demand.
% sólo tsc++ seguía esa premisa

% qué es necesario para cumplir esa visión y porque las alternativas vistas no se ajustan:
% dependencias mínimas, protocolos estándares, etc.


% TODO CONCLUSION
Although both semantic space-based and Triple Spaces implementations exist, none of them has been specifically designed to be run in devices with constrained capabilities apart from our solution.




% ----------------------------------------------------------------------



%: ----------------------- TSC for resource constrained devices ------------------------
% ----------------------------------------------------------------------

\begin{savequote}[50mm]
The original idea of the web was that it should be a collaborative space where you can communicate through sharing information.
\qauthor{Tim Berners-Lee}
\end{savequote}


\newcommand{\codigo}[1]{``\texttt{#1}''}
\newcommand{\primquery}{\emph{query}}
\newcommand{\primread}{\emph{read}}
\newcommand{\primtake}{\emph{take}}
\newcommand{\primwrite}{\emph{write}}


\chapter{Triple Spaces for Limited Devices}
\label{cha:tsc}
\newcommand{\pathchapthree}{3_tsc}

% the code below specifies where the figures are stored
\ifpdf
    \graphicspath{{\pathchapthree/figures/PNG/}{\pathchapthree/figures/PDF/}{\pathchapthree/figures/}}
\else
    \graphicspath{{\pathchapthree/figures/EPS/}{\pathchapthree/figures/}}
\fi



% ----------------------------------------------------------------------



% TODO TODO TODO me he dado cuenta que si quiero interoperar, la clave sería meter en algún lado algo sobre LOD y SPARQL Update 1.1!!!
% Igual que quiero interoperar con Semántic WoT, debería interoperar con LOD !!!



In recent years, the \acf{iot} has become a reality due to the increasing number of everyday objects equipped with computing and networking capabilities.
The use of these objects together with the rising of mobile computing greatly contributes to the \ac{ubicomp} vision.
%In the beginning the community put more effort on those devices' connectivity issues, materialized in the spread of technologies such as Zigbee\footnote{http://www.zigbee.org}, 6LoWPAN\footnote{http://datatracker.ietf.org/wg/6lowpan/charter/} or Bluetooth\footnote{http://www.bluetooth.com}.
%However, nowadays it is focusing on the architectural solutions used by the applications built around these objects.
%A classification of these solutions has been already presented in the previous chapter. % Coulouris et al.
In this thesis, we defend that \acf{tsc} provides some valuable properties to \ac{ubicomp} environments. % mencionarlas?
However, some other properties are derived from how \ac{tsc} is implemented. % TODO implemented suena muy umpa-lumpa? mejor designed o así? o hablar de model mejor que de middleware?


% TODO problema: y por qué no otros? Habría que redactarlo de tal forma que no haga que un lector tiquismiquis requiera un puto survey de network-based architectural styles!
In this chapter, we focus on the design of the space model for our solution.
%selecting the underlying network-based architectural style.
Its application to resource constrained devices brought the successful \ac{wot} initiative.
\ac{wot} is probably the most prominent architectural solution used for smart-objects. % probably porque no tengo datos que lo defiendan
%The benefits that \ac{rest} brings to resource constrained devices have been deeply described in the previous chapter.


To that end, Section~\ref{sec:align_tsc_http} analyzes the characteristics of the \ac{rest} architectural style and the similarities it shares with \ac{tsc}.
Later, it describes a \ac{rest}-like \ac{api} to access a triplespace. % TODO ver cómo escribir triplespace
The goal pursued is two-fold: (1) inherit valuable the properties from \ac{rest} and (2) maintain a high degree of compatibility with the \ac{wot}.
Section~\ref{sec:distributed_tsc} presents the need and requirements of a distributed architecture for \ac{ubicomp}.
These requirements lead us to adapt the \ac{tsc} model.
Finally, Section~\ref{sec:middleware_qualitative_evaluation} evaluates the properties achieved with this model.


%    propiedades aportadas por REST a IoT

\section{On the \acs{rest}fulness of \acs{tsc}}
\label{sec:tsc_restfulness}

\subsection{\acl{rest}}
\label{sec:rest}

\acf{rest} is a network-based architectural style proposed by \citet{fielding_architectural_2000}.
% descripción de propiedades de REST: hipermedia distribuído
% que es hipertexto? http://roy.gbiv.com/untangled/2008/rest-apis-must-be-hypertext-driven#comment-718
It aims to cover certain properties, explained in Section~\ref{sec:network_properties}.
% Estilos de los que se deriva: ¿?
To achieve these properties, \ac{rest} establishes the following constraints from other network-based architectural styles:
% explicarlo o es demasiado obvio?
\begin{description}
 \item[\acf{rest_cs}.] Providing an \ac{api} to the clients, they are isolated from back-end implementation details.
		       % esto ayuda a: , scalability, evolvability (apps independientes pueden evolucionar mejor)
 \item[\acf{rest_s}.] The state is fully stored in the client and therefore each request has all the information needed to process it.
 \item[\acf{rest_cache}.] When added to the \ac{rest_cs} constraint, this style replicates content obtained from a server in the client.
 \item[\acf{rest_u}.] It is the key constraint which distinguishes \ac{rest} from other architectural styles.
                      This constraint is composed by the following ones:
    \begin{description}
	% explicados en sección 5.2 de Fielding, resumir?
	\item[\acf{rest_id}.] Resources are the conceptual targets of hypertext references.
			      Their identification offers a generic interface to access and change the values of a resource.
	\item[\acf{rest_rep}.] Representations are composed by a sequence of bytes and the metadata to describe those bytes.
	\item[\acf{rest_desc}.] The client and server have to agree on standard methods and media types. % linkar al tipo de la Web esa
				% Explicación Fielding:
				% interaction is stateless between requests
				% standard methods and media types are used to indicate semantics and exchange information
				% responses explicitly indicate cacheability.
		                %\ac{http}'s content negotiation for instance, allows to reach an agreement on the media types.
	                        Beyond that point, each request or response should contain all the needed data to process it.
	                        Therefore, in \citeauthor{fielding_architectural_2000}'s words,
	                        the type should be registered,
	                        the registry should point to a specification and 
	                        the specification should explain how to process data according to its intent \citep{wahbe_self-descriptive_2010}. % no necesariamente tiene que ser estándar
	\item[\acf{rest_hateoas}.] This is a controversial constraint because most of the self-proclaimed ``\ac{rest}'' \acp{api} fail to follow it \citep{house_how_2012}.
	                           It states that no out-of-the-band information should guide the interaction with an \ac{api}.
	                           Instead, the hypertext should guide it.
	                           In other words, the client must know just an initial URL and the application's media types.
	                           From that point, it should select the alternatives proposed by the server to change to the next application state \citep{fielding_rest_2008}.
	                           % propuestas: respresentaciones (links) y la manipulación implicita de las mismas (CRUD)
    \end{description}
 \item[\acf{rest_l}.] Each layer provides services to the top layer. % e.g. TCP/IP
 \item[\acf{rest_cod}.] The client has a set of resources which does not know how to process.
                        This \emph{know-how} is downloaded from the server.
\end{description}


\subsection{Suitable Protocols for \ac{rest}}
\label{sec:protocols}

% sólo por el correcto uso de ciertos protocolos se puede conseguir todo menos HATEOAS => citar el vídeo aquel

% hablar un poco de HTTP y CoAP y decir por qué no hablamos más a menudo de CoAP

Historically, \acf{http} has been considered a suitable protocol to follow the \ac{rest} style. % TODO cite
\ac{http} is a simple protocol whose adoption by computing platforms is massive.


However, in the last few years \acf{coap} has emerged as a specialized web transfer protocol for resource constrained devices. % has emerged, has arisen?
Some noteworthy features of \ac{coap} are
(1) the reduced message size,
(2) the use of UDP as a transport layer (with the possibility of using \emph{multicast} communication),
(3) similarity with \ac{http} (both to reuse its properties and to ease cross-protocol proxying),
(4) a resource discovery mechanism. % TODO cita a CORE
% mencionar seguridad?


One could defend that to design a lightweight \ac{tsc} solution \ac{coap} should be used as a baseline.
However, we have chosen to work with \ac{http} for the following reasons:
\begin{itemize}
  \item RDF-based media-formats can be rather verbosed.
	This contrasts with \ac{coap}'s message size limitations.
	Dealing with this limitations was not one main goals of the thesis.
	However, we have considered them at some points of the dissertation to avoid unrealistic assumptions.
  \item \ac{coap} is an ongoing standard.
        Therefore, its definition slightly changes in each draft version.
        A practical limitation of this is that at the moment there are few libraries
        and tools to work with \ac{coap}.
        This limits the range of platforms which could adopt any solution proposed.
  \item Due to its similarities with \ac{coap}, the future adoption of the latter should be relatively straightforward.
        % sin contar con los aspectos destacados anteriormente
\end{itemize}

% Ideas no usadas:
% REST no dice nada, pero de facto el protocolo es HTTP
% durante el desarrollo de esta tesis en resource constrained ha surgido CoAP
%   muchas de las cosas aquí comentadas son iguales o mejor sobre CoAP
%   de todas formas, ha habido algo que ha guiado nuestro diseño: librerias existentes => afecta a plataformas
%         si, pero para eso está lo de layered approach! => motivos prácticos



\subsection{\acl{tsc}}
\label{sec:tsc_vs_rest}

\ac{tsc} consist of a shared space which is accessed using some known primitives.
\emph{Per se}, \ac{tsc} does not contradict in any sense the \ac{rest} principles:

\begin{description}
 \item[\ac{rest_cs}.] Providing the access to each space is centralized in a machine, clients can access to it using \emph{CS}.
           This does not prevent to use a distributed semantic repository in the backend.
 \item[\ac{rest_s}.] The primitives to access the space imply simple read and writes which do not store any state in the server.
 \item[\ac{rest_cache}. ] The semantic content stored in the space can be cached.
           However, the dynamism of the underlying network makes caching challenging to configure. % quería poner "tune up" pero "configure" me ha sonado más fino
 \item[\ac{rest_u}:]
    \begin{description}
	\item[\ac{rest_id}.] There are types of identified resources in \ac{tsc}: spaces, RDF Graphs and certain elements of the RDF Triples.
	                 All of them are identified by URIs.
	                 The space can be seen as a coarse-grained view of the underlying graphs. % TODO coarse grained es lo contrario de fine-grained?
	                 The graphs have sets of triples which are usually related and describe a unit of knowledge. % A RDF triple by itself cannot transmit too much information
	                 Self-identified triple's subjects, predicates or objects are the source of concept linking.
	\item[\ac{rest_rep}.] The RDF graphs and triples mentioned above can be represented using different standard serializations.
	\item[\ac{rest_desc}.] The messages derived from the primitives can be self-describing. % including meta-data
	\item[\ac{rest_hateoas}.] This is the most challenging property to achieve with RDF-based representations.
	                      However, some relevant solutions propose means to fulfill this property \cite{verborgh_functional_2012}. % TODO citar más de este y de otros que han hablado LOD, SPARQL y RESTful
    \end{description}
 \item[\ac{rest_l}.] encapsulation of functionalities can be achieved through a layered system.
 \item[\ac{rest_cod}.] Although it may not be as common as scripting in ordinary human-browsing, % TODO reexplicar esto mejor?
            the \emph{know how} can be downloaded from the server:
            \begin{enumerate}
	      \item Modeling it through ontologies.
	      \item Through semantically expressed rules which are processed with an inference engine (e.g. N3 rules) \citep{verborgh_functional_2012}.
            \end{enumerate}
\end{description}

% CITAR:
% Algunas ideas rescatables de la anterior intro:
% Both \ac{wot} and \ac{tsc} can be considered resource oriented solutions since they put emphasis in giving access to data resources.
% While \ac{tsc} enables expressive queries, dynamic discovery and non human-mediated cooperation among objects, the \ac{wot} adopts the scalable properties of the World Wide Web and it is entirely based on web standards.



\subsection{Aligning \acs{tsc} with \acs{http}}

% TODO Revisar cosas del texto original que no encajaban aquí (están comentadas)
% TODO citar trabajos de:
%      TSC es como Web para semántica
%      RESTful semántico
%      So far, the compatibility of \ac{tsc} with the \acs{http} RESTful style has been proved from a formal \cite{hernandez_formal_2010} point of view.
%      Papers sobre TSC y su discurso sobre la WWW

In this section, we present a \ac{tsc} \acs{api} over \ac{http} to give a practical overview of \ac{tsc}'s \acs{rest}fulness.

\subsubsection{\acs{tsc} resources}

Three key concepts are important to understand the resources: agents share information in a common \textbf{space}.
A space is identified by an \acs{uri}.
Therefore, all the operations in \ac{tsc} are performed against a particular space.
%By default, all applications connect to a common standard space, but they can optionally choose to connect to a particular private space.
Within a space, the information is stored in sets of \textbf{triples} called \textbf{graphs}.
Each graph can also be identified by an \acs{uri}.
The \acs{rdf} \textbf{triples} are the underlying concept of all the \ac{sw} languages.
Each triple is composed by a subject (which is a \acs{uri}), a predicate (also a \acs{uri}) and a value (which can be a \acs{uri} or a literal), as shown in the Figure~\ref{fig:triples_example}.

As detailed later, most common operations in \ac{tsc} attempt to add or remove graphs.
% este es un tema "delicado" desde el punto de vista de TSC, así que mejor si lo tratamos a parte:
% as well as to query for graphs or for sets of triples retrieved from different graphs.
In order to perform the queries, which enable the selection of a subset of the semantic content hold in a given space, a \textbf{template} is required.
The wildcard templates used by default
\footnote{
  % Por qué no SPARQL?
  More sophisticated query languages like SPARQL (\url{http://www.w3.org/TR/rdf-sparql-query/}) could be used. % sophisticated o expressive?
  However, not many embedded platforms have parsers for these languages available.
  This could be an obstacle for the adoption of the middleware.
  Therefore, we have opted for using wildcard templates, which are in turn much more simple to process.
  In any case, the nodes able to parse these query languages can easily decompose a query on wildcard templates.
}
are special triples with optional wildcard subject, predicate and/or object.
For example, the template \texttt{?s foaf:knows gomezgoiri:aitor} could be employed to select instances which represent people who know Aitor (see Figure~\ref{fig:tsc_resources}).

% TODO Discussión profunda de SPARQL y REST???

\InsertFig{tsc_resources}{fig:tsc_resources}{
  Schematic view of a space with four graphs and sample triples for one of these graphs.
}{
  Aliases for the beginning of most URIs, known as prefixes in most Semantic languages, are used for the lack of clarity.
}{1}{}


\subsubsection{Adopted \acs{tsc} primitives}
\label{sec:primitives}

\acs{tsc} derives some primitives originally defined in the Linda language \citep{gelernter_generative_1985} to access to the semantic information hold in each graph.
In this section, we explain these primitives.

\begin{itemize}
 \item The \textbf{write} primitive allows writing a graph into a given space (identified by its \acs{uri}).
	The set of triples received by this primitive will be stored together in the same graph, returning the \acs{uri} which identifies that graph.
	The graph \acs{uri} can be used to access directly to that graph later on, or to create new triples and relate contents.

  \begin{lstlisting}
    write(space_URI, triples): URI              [1]
  \end{lstlisting}


  \item The \textbf{read} returns a graph belonging to a given space which contains at least a triple matching the given template or has the given \acs{uri} as its identifier.
	If more than one graph fulfill one of these conditions, just one of them is returned (nondeterministically).
	It should be remarked that it has been designed as a non blocking operation.
% why?
% TODO READ al ser no-deterministica puede afectar a caching???
  \item The \textbf{take} primitive behaves like a destructive \textbf{read}, deleting the graph returned from the space.

  \begin{lstlisting}
    read(space_URI, graph_URI): triples         [2]
    read(space_URI, template): triples          [3]
    take(space_URI, graph_URI): triples         [4]
    take(space_URI, template): triples          [5]
  \end{lstlisting}

%  \item Space management primitives.
%	A node can join or leave a space using \linebreak \texttt{joinSpace(space\_URI)} or \texttt{leaveSpace(space\_URI)}.
\end{itemize}



\subsubsection{\acs{http} \acs{api} for \acs{tsc}}
\label{sec:httpapi}

As both \ac{tsc} and \ac{http} are \ac{rest} compliant, their similarities are evident.
In the same way \ac{tsc} has the already explained primitives, \ac{http} has verbs to get, create, update or remove resources (GET, PUT, POST, DELETE).
%The main purpose of \acs{tsc} primitives is not to access data but to coordinate the nodes accessing to that data.
Consequently, the translation between these two worlds is straightforward.

\begin{sloppypar}
According to the \ac{rest} principles the interaction with an API must be hypertext-driven \citep{fielding_rest_2008}.
Therefore, the human-oriented representation of the \ac{api} should be completely \emph{browseable} (e.g. using \acs{html}).
Regarding the machine-oriented (i.e. semantic) representation of the \ac{api}, several solutions to enable hypermedia driving have been proposed by \citet{verborgh_functional_2012} and \citet{kjernsmo_necessity_2012}. % mencionar que son experimentales?
Independently of the mechanism adopted, in this section we propose a optional \ac{api} which tries to stress \ac{tsc}'s compliance with \ac{http}.
\end{sloppypar}

The list of spaces a node is joined to are available under \textit{/spaces}.
Each space is identified by an \acs{uri} (e.g. \url{http://space1}).
All the resources of that space, both real (i.e. graphs) or virtual (i.e. query) are listed under $/spaces/\{space\_uri\}$ (summarized by \emph{sp} from now on).
Each graph is available on ${\{sp\}/graph/\{graph\_uri\}}$.
If we make an \acs{http} DELETE to that resource, under \ac{tsc}'s perspective, we take that graph from the space.
The rest of the mappings are shown in the Table~\ref{tab:tscAPI}.

\begin{table} %http://en.wikibooks.org/wiki/LaTeX/Floats,_Figures_and_Captions#Wide_figures_in_two_column_documents
  \centering
  \caption {
    \acs{http} mapping for the primitives detailed in the Section~\ref{sec:primitives}. \textit{sp} is a space \acs{uri},
    \textit{g} is a graph \acs{uri}, \textit{s}, \textit{p} and \textit{o-uri} are subject, predicate and object \acsp{uri} or wildcards (represented with an as \textit{*}).
    When the template's object is a literal, it can be expressed specifying its value (\textit{o-val}) and its type (\textit{o-type}).
    \medskip
  }
  \begin{tabular}{c|l|c}
      \acs{http} request & \acs{url} & Returns \\
      \hline
      POST & \{sp\}/graphs/ & [1] \\
      GET & \{sp\}/graphs/\{g\} & [2] \\
      GET & \{sp\}/graphs/wildcards/\{s\}/\{p\}/\{o-uri\} & [3] \\
      & \{sp\}/graphs/wildcards/\{s\}/\{p\}/\{o-type\}/\{o-val\} & \\
      DELETE & \{sp\}/graphs/\{g\} & [4] \\
      DELETE & \{sp\}/graphs/wildcards/\{s\}/\{p\}/\{o-uri\} & [5] \\
      & \{sp\}/graphs/wildcards/\{s\}/\{p\}/\{o-type\}/\{o-val\} & \\
  \end{tabular}
  \label{tab:tscAPI}
\end{table}


\paragraph{Status Codes}
\label{sec:status_codes}
The \acs{api} should be compliant with the standardized \acs{http} status codes.
These codes are sent back in the response as part of the header.
For instance, the \ac{tsc} middleware returns the 404 error when no significant result can be found for a primitive.
This adoption, apart from enhancing the compatibility with other web applications, could enable the modular adoption of the API.
For example, if a node does not offer a wildcard based \textit{query}, it will not affect the behavior of the rest of the nodes of an space.
Instead, they will interpret these cases as empty responses.
This modularity becomes crucial to ease the partial adoption on new platforms.

\paragraph{Content Negotiation}
Another key aspect of the \acs{http} protocol the \acs{api} should take advantage of is the \textit{content negotiation}.
This mechanism allows to specify the desired representation for a content on the client side and to express what representation is sent as a response from the data provider side.
For that purpose, the client adds an \textit{Accept} field to the \acs{http} header with a weighted list of media types it understands.
Then, the server will answer with the best possible format it knows about, specifying the \textit{Content-type} in the response.

The benefits of using this mechanism in a \ac{tsc} middleware are two-fold.
Firstly, it enhances the browsability of the primitives with human understandable \acs{http} responses.
Secondly, it allows different semantic representations (e.g. RDF/XML\footnote{\url{http://www.w3.org/TR/REC-rdf-syntax/}}, N-Triples\footnote{\url{http://www.w3.org/2001/sw/RDFCore/ntriples/}} or N3\footnote{\url{http://www.w3.org/TeamSubmission/n3/}}).
The latter characteristic becomes crucial since not all the nodes may understand all the formats (e.g. a mobile phone may not have a RDF/XML parser).
%This is true even if all the languages use the same basic concepts (i.e. \acs{rdf} Triples).
In these cases, the compatibility of both sides can be ensured through a conversion carried out in the server side.
Furthermore, expressing the preference for a semantic format can be useful too in other cases.
For example, to obtain the less verbose answer.
\section{Distributed \ac{tsc}}

% cómo y porqué hacer un TSC distribuído
In the previous subsection, we showed a basic alignment between \ac{tsc} and \ac{http} to stress the closeness between \ac{tsc} and \ac{rest}.
An important point of this alignment was the access to the data in client-server basis.
This implies that all the data of a space goes through a server. % para posiblemente almacenarlo en otras máquinas

% Pero un entorno Ubicomp es eminentemente distribuído
However, the data in a \ac{ubicomp} environment is not generated in a single point.
%   La información se genera en fuentes separadas, ¿cómo accedo a los últimos datos?
In fact, each sensor in an intelligent environment is potentially generating data constantly.
Besides, the magnitude of this sensors can be enormous.
This creates a trade off between efficiency and update difficult to overcome:
\begin{itemize}
  \item The more frequently each sensor sends contents to the server, the more inefficient is the solution in terms of network usage.
  \item The less frequently the communication is, the less updated is the data in the server.
        This leads to spaces which misrepresent the environments.
\end{itemize}
A sensitive solution is to let these sensor manage their own information.

%   ademas no practico:
%        las fuentes del conocimiento son las que mejor saben como gestionar una informacion (crearla, acualizarla o eliminarla)
%        la información se transporta en dispositivos
%        reusar datos en distintas apps
Furthermore, this strategy can not only be useful for devices which constantly generate new data.
Delegating responsibilities between the devices presents several other benefi y a for a number of reasons:
\begin{itemize}
  \item They will know how to represent these contents better than others. % ejemplo en grados centigrados o celsius?
	For example, the unit of a temperature measure.
  \item They know when to create, update or delete data. % ejemplo if a server receives 2 contradictory measures... algo
        In fact, they can opt for creating data on demand.
  % TODO esto es más un beneficio de REST que de otra cosa
  \item The data can be reused by other applications or even other spaces. % Interop!
        These applications would not be required to use a space-based approach.
  \item Carrying that information can be useful in certain mobility situations.
        For example, a person which stores her user profile in the smartphone.
        She can shares it with different spaces or applications as it moves along the city.
\end{itemize}


% Aún así, que los dispositivos se coordinen a través de una entidad central es conveniente:
%       Sencillo de implementar: notificaciones, razonamiento, etc.
%       Incluso si ese espacio luego está distribuído o cambia de máquina, acceso a través de dicha interfaz (HTTP o CoAP el día de mañana)
%       Modo: caching (para reliability) o modo ver qué hay en el espacio ahora mismo (visión en tiempo real)
%       -->> INDEPENDENCIA DE LOCALIZACIÓN!!!
Still, most of the \ac{iot} devices or mobile phones are unreliable: they can join and leave at any moment.
Therefore, distributing the shared space among unreliable nodes comes with a number of drawbacks:
\begin{itemize}
  \item This unreliability makes difficult to maintain the consistency of the space. % TODO CUIDADO, puedo estar metiendome en un fregado, la parte de los dispositivos no será consistente!
        Since this information drives the coordination between the devices, its consistency its important.
  \item Blocking mechanism is important in space-based computing.
        For example, a worker node may block until some task is become available in the space.
        A way to implement it in a distributed space is by means of a notification system.
        Again, implementing it using unreliable devices is hard to implement. % no sé si esto último no sonará demasiado a mofa...
  \item And last but not least, if the access to some data in the space cannot be guaranteed no matter if the device which wrote that information is available, there will not be \emph{time uncoupling}.
\end{itemize}
Overcoming the previous difficulties, usually implies a high network traffic.
However, these negatively influences the energy consumption of some nodes that whose energetic autonomy is limited.


% solución enfoque híbrido:
%      grafos autogestionados (no takeables por otros) => write en local
%      cambios en actuadores => directamente a través de HTTP o indirectamente a través de tasks escritas en espacio
%                               o mejor: podrían esas tasks ser directamente esos servicios???
%      grafos que sí => write al "servidor HTTP"
%      query => en todos los dispositivos (capítulo 4) - Porque a veces es necesario a través de todos los nodos
%      take + read => sobre los grafos takeables (o inferencia con toda la info del espacio, cómo prefieras)




% Ventajas de meclar 2 cosas aparentemente separadas?
%   Para WoT:
%      comunicación indirecta para cacharros
%          local+ref
%              de contenidos relativos a coordinación
%              podríamos cachear contenidos de otros en el espacio central (reliability+energy)
%          ref
%              sacar a colación los problemas aquellos de query sobre muchos recursos
%
%   Para TSC (respecto a otras como Smart-M3):
%      limitamos qué gestiona la entidad central => sólo coordinación, el resto lo coge de los cacharros
%      se reconoce su autonomía permitiendo gestionarse su propio contenido => REST, escalabilidad
%      interop => SW y posibilidad de extender a otros servicios REST
%      multiplataforma => facilidad para portarlo
%
%   Para desarrollador:
%      API integrada y simple

By doing so we retain the desirable properties of both approaches:



% presentar las 2 partes
% condiciones relajadas: consistencia, razonamiento distribuído

















The \acs{http} mapping is done among a \textit{client} and a \textit{server}.
However, as previously detailed, \ac{tsc} provides reference autonomy, so a consumer will query the space without knowing the particular addresses of the nodes composing that space.
This autonomy is managed at other upper and optional layers explained in the following section.
However, this \acs{http} mapping is all a data provider needs to serve information in the space.
% Quitado porque ni yo mismo lo entendía ahora...
%, and it is also what a data consumer needs to interrogate a multicast provider (a component provided by Otsopack to contact multiple nodes in a particular space).


\footnote{
  Note that the write primitive is explicitly excluded from the table.
  This primitive could be directly mapped to a \acs{http} POST,
  but we defend that each node should manage its own information
  (see Section~\ref{sec:knowledge_distribution}).
  Therefore, the \acs{http} \acs{api} does not allow remote writing.
}
\section{Conclusions} % TODO llamarlo Evaluation?
\label{sec:middleware_qualitative_evaluation}

With the solution presented in the previous section, we aimed to retain the desirable properties of both \ac{rest} and \ac{tsc}.
% analísis más largo en la siguiente sección
% 1. cómo sé que esto es mejor que usar REST y TSC por separado? Sinergia?
% 2. que aporta unirlo bajo un mismo middleware?
% 3. tiene beneficios?
However, some questions arise from that integration:
(1) does in fact retain all the properties of \ac{rest} and \ac{tsc} separately? Otherwise, which properties are affected?
(2) what other benefits carries comparing with using already existing \ac{tsc} and \ac{rest} middleware?

% Fielding distingue entre:
%      + "distributed-system": hace que el usuario vea el sistema como un todo, sin saber cuando se realizan llamadas en remoto y cuando no
%      + "networking-system": no necesariamente oculta las particularidades de la red <- el se centra en estos


To answer these questions, we analyze the presented solution from different point-of-views:
(1) its coordination properties (Section~\ref{sec:XXX_properties}),
(2) its network-level properties (Section~\ref{sec:network_properties}),
(3) how the latter work together to contribute to the challenges of an \ac{ubicomp} environment (Section~\ref{sec:middleware_properties}).
Section~\ref{sec:middleware_eval_summary} summarizes these analysis. % cómo dan respuesta a las preguntas del comienzo


% TODO cosas rescatables del paper original de WoT 2011:
%     Uniform API vs. specific usage of each API !!!
%     Coupling: lo comido (autonomies), por lo servidor (client-server coupling por tener API común)
%          In spite of the outlined decoupling nature of both approaches, the data definition can be considered a coupling mode indeed.
%          While in \ac{wot} each resource defines its own data formats and contents themselves, in \ac{tsc} the ontology in which the semantic
%          concepts are described must be known by each part of a distributed application to effectively cooperate among them.
%     Discovery: HATEOAS puede ser complejo para cacharros peques (parseo),
%                proponemos discovery basada en API uniforme y a nivel de cacharros con algún mecanismo en el que no entramos.
%          One of the main drawbacks in \ac{wot} is the lack of a discovery mechanism for new objects and the data they provide.
%          Even when this data can be linked in each object response (using HATEOAS) and microformats are sometimes included to ease the search-ability of these objects by search engines, it is difficult for an object which may change of location and context to be referred.
%          Thus, \ac{wot} may have a tendency to create isolated islands of data.
%          Several workarounds have been proposed to overcome this limitation, such as using a central repository\footnote{http://www.pachube.com/}, a framework which uses federated repositories responsible for different administrative domains \citep{stirbu_towards_2008} or making each connected sensor announce itself to let an intermediary know its presence \citep{kamilaris_smart_2010}.

%     Scalability: se ha hablado y ahora no es tan mala como antes ;-) Entre cacharros seguiría dando asco
%     Semantics: Microformats vs Full (hablar de esto)
%     Comprensión por parte de usuario: sencillo misma interfaz, no tienes que andar descubriendo nuevas APIs y viendo cómo funcionan, qué formato devuelven, etc.
%     Comprensión por parte de desarrollador: API sencillica siempre


\subsection{Coordination properties}
\label{sec:coordination_properties}

% TODO añadir la autonomía que se refería a la semántica?
As we have already discussed in Section~\ref{sec:integration}, indirect communication middleware can have two key properties:
\begin{itemize}
  \item The sender does not need to know the receiver or receivers and vice versa (i.e. \emph{space uncoupling}).
  \item Senders and receivers do not need to exists in the same time to communicate with each other (i.e. \emph{time uncoupling}).
\end{itemize}


The primitives used in our space-based computing implementation force it to be \emph{space uncoupled}.
On the contrary, it is not always \emph{time uncoupled}.
If two nodes write and read from the coordination space, they are \emph{time uncoupled};
but if a node accesses to others' content (i.e. \emph{self-managed graphs}), they are not.
In other words, \emph{time uncoupling} is not achieved by the extension of the model presented in Section~\ref{sec:halfway_solution}.
As a future work, this could be alleviated by a caching mechanism implemented in the space holder or holders. % TODO hablar con propiedad de esos elementos => nombrarlos antes
% see \ref{tab:middleware_coordinationprop}

\InsertTab{tab:middleware_coordinationprop}{Uncoupling levels achieved by the different parts of the middleware presented}{}{
  % Tablita de cómo hereda propiedades de los estilos anteriores?
  % (y si quieren más información, que miren en la tesis de Fielding)
  \begin{tabular}{lcc}
    \hline
    ~ & Space &	Time \\
    ~ & uncoupling & uncoupling \\
    \hline
    % O mejor ponerlo con Yes y No?
    Coordination space & $\checkmark$ & $\checkmark$ \\
    Outer space & $\checkmark$ & × \\
    \hline
  \end{tabular}
}{htbp}


\subsection{Networking properties} % Architectural Properties for Network-based Styles
\label{sec:network_properties}
% Cómo afecta la app propuesta => ir actualizando al final de otras secciones???

In Section~\ref{sec:tsc_vs_rest} we have described how \ac{tsc} does not contradict any of the \ac{rest} principles.
However, the adaptation presented in this thesis does not completely adhere to the \ac{rest} style.
In this section we present how this divergences affect the \ac{rest}'s properties.
% decir que el objetivo final es que pudiera ser REST

% TODO antes de entrar al lio, explicar cómo funciona esto de REST, en qué parte de nuestra solución lo necesitaremos
% decir que API para TSC mola
% decir que API para los cacharros deberia de ser independiente a lo que nosotros digamos y por eso estaria bien REST
%     de todas formas, como es dificil, nos quedamos en un nivel de madurez menor


In \citeauthor{fielding_architectural_2000}'s words, the relevant properties which describe a network-based system are the following ones \footnote{We refer to the reader to \citet{fielding_architectural_2000} for the complete thorough analysis.}:
% verificar que no se parece demasiado a sus definiciones o sino poner cursiva.
\begin{description}
  \item[Performance] is divided into network performance, user-perceived performance and efficiency.
    \begin{description}
      \item[Network performance] is affected by the styles in the number of interactions and the granularity of data elements.
      \item[User-perceived performance] refers to the impact a user in front of an application perceives. % UP Performance es una feature que no aporta mucho a nuestro caso
      \item[Efficiency] is achieved by minimizing the use of the network.
    \end{description}
  \item[Scalability] measures how an architecture supports a big amount of components and interactions between them.
  \item[Simplicity] is achieved through the separation of concerns for the components and the generality of architectural elements.
  \item[Modifiability] indicates how easily gradual changes can be introduced in the system.
                       These changes form different implementations which should coexist.
                       % Dynamic modifiability avoids the whole system to restart.
    \begin{description}
      \item[Evolvability] refers to the degree in which a component can be implemented without negatively impacting on others.
      \item[Extensibility] measures the ability to add functionality to the system.
      \item[Customization] is the ability to customize the behavior of an architectural element temporarily.
      \item[Configurability] is related with the extensibility and the reusability.
      \item[Reusability] is the ability to reuse components, connectors or data elements without modifying other apps.
    \end{description}
  \item[Visibility] is the ability of a component of monitoring or mediating in the interaction between other two components.
  \item[Portability] is the ability of working in different environments.
  \item[Reliability] is the degree in which an architecture depends on the failures of the system or components, connectors or partially incorrect data.
\end{description}

% Requisitos de ealy web, propiedades deseables
%   Low entry-barrier
%   Extensibilidad
%   Distributed Hypermedia
%   Internet-scale

The Table~\ref{tab:network_properties} summarizes how different architectural styles achieve these properties.
Particularly, it shows the styles from which \ac{rest} derives (see Section~\ref{sec:rest}).
The ultimate goal of the solution explained in the previous section is twofold:
\begin{enumerate}
  \item Provide a \ac{rest} access to a semantic space.
       This would ease its integration with the rest of the web.
  \item Enrich that space with the knowledge provided by other \ac{rest} \acp{api}.
\end{enumerate}
Therefore, ideally, its networking properties would be the sames as the \ac{rest} style.


However, defining a 100\% \ac{rest} compliant \ac{api} is not easy.
Indeed, most of the self-proclaimed \ac{rest}ful \acp{api} are not \citep{house_how_2012}. % la nueva buzz word para eso: hypermedia API
The main cause is the \emph{HATEOAS} constraint seen in the Section~\ref{sec:rest} \citep{fielding_rest_2008}.
Using \ac{http} as an application-level protocol forces a developer to comply the rest of the constraints \citet{moore_hypermedia_2010}, but not \emph{HATEOAS}.
% TODO hablar del Richardson Maturity model!


% semantic-based \ac{api} is not easy. => referencias a los que lo han intentado, a que no hay hipertexto, a que te puede interesar inferir con bastantes
Regarding the \ac{sw}, some recent efforts have tried to move closer to the \ac{hypermedia} constraint \citep{steiner_fulfilling_2011,kjernsmo_necessity_2012}.
However, these worlds remain quite isolated.
In fact, \ac{sw} API  usually present another main divergence with the \ac{rest} style: the use of query endpoints.
% i.e. no devuelven listas de enlaces, sino todo a lo bruto
% esto quiere decir que no se devuelve todo y ala, tu procesa esa burrada de datos, más bien se filtran
These endpoints intend to solve some inefficiency issues which \ac{rest} shows when working with a big amount of data. % TODO mencionar \ac{lod} en algún punto de aquí?
To that end, they allow expressive query languages and offer results where the boundaries of the different resources often blurs \citep{wilde_restful_2009}.
The most paradigmatic case may be to decide to which graph does a inferred content belongs. % ejemplo de lo anterior


% Aterrizar esto a IoT y a nuestro caso:
%        La eficiencia es importante => para evitar hacer 500 llamadas a un cliente
%        La reusabilidad es importante también => para potencialmente adaptarse a cualquier WoT semántico
%           en el futúro => más trabajo en este área
%       Solución de compromiso:
%           nivel de madurez X de Richardson. Suponemos que alguien debe serguir nuestro API mínimo en los objetos WoT

In \ac{iot} both the efficiency on the communications and reusability of the \ac{api} are important:
\begin{description}
  \item[Efficiency:] mobile and embedded devices' have restricted energy autonomy.
                    This autonomy is severely affected by network communications. % TODO cite
                    Particularly, the access to the data they provide by means of hyperlinks may result in many HTTP requests. % aunque eso se puede adaptar por lo visto
                    %  Por que sería muy costoso hacer rollo araña?!
		    %     lo que ganas de reusabilidad, lo aumentas en complejidad en el cliente
                    %     interpretar código del servidor no es tampoco super-sencillo (xHTML!)
  \item[Reusability:] due to the heterogeneity of the devices, assuming that they all share a common and unevolvable \ac{api} may not be realistic. % TODO paliar para que no parezca que no tiene sentido
\end{description}
The first solution together with the current difficulties on achieving a fully and standard \emph{HATEOAS} for machines in the \ac{sw}, lead us to opt for the second level on the Richardson Maturity level.
Besides, our \ac{tsc} \ac{api} will be \emph{HATEOAS} \ac{api} in its humans representation to ease their learning level.% pero aún así el API de TSC seguirá siendo HATEOAS para los usuarios, no para las máquinas
However, we do not discard to work towards a fully \ac{rest}ful \ac{api} as future work.


In conclusion, the properties of the properties for network communication style we adopted corresponds with \emph{LCODC\$SS} plus simplicity and visibility. % hacer una nueva fila con esto
% TODO comprobar que se explica eso tal cual en la tesis de Fielding: HATEOAS => sólo afecta a reusabilidad

% TODO Comprobar que no se salga por el borde derecho!
\InsertTab{tab:network_properties}{Properties of different architectural styles for network-based applications as defined by \citet{fielding_architectural_2000}}{
  Note that the original table has been slightly adapted.
  These adaptations are remarked inside the table. % using several footnotes.
}{
  % Tablita de cómo hereda propiedades de los estilos anteriores?
  % (y si quieren más información, que miren en la tesis de Fielding)
  \footnotesize
  \begin{tabular}{lccccccccccccc}
    Style &
    \rotatebox{90}{Net Perform} &
    \rotatebox{90}{UP Perform} &
    \rotatebox{90}{Efficiency} &
    \rotatebox{90}{Scalability} &
    \rotatebox{90}{Simplicity} &
    \rotatebox{90}{Evolvability} &
    \rotatebox{90}{Extensibility} &
    \rotatebox{90}{Customiz.} &
    \rotatebox{90}{Configur.} &
    \rotatebox{90}{Reusability} &
    \rotatebox{90}{Visibility} &
    \rotatebox{90}{Portability} &
    \rotatebox{90}{Reliability} \\
    \hline
    CS & ~ & ~ & ~ & $+$ & $+$ & $+$ & ~ & ~ & ~ & ~ & ~ & ~ & ~ \\
    S\footnote{\emph{S} represents the difference between \emph{CSS} and \emph{CS} in \citep{fielding_architectural_2000}.}
      & $-$ & ~ & ~ & $+$ & ~ & ~ & ~ & ~ & ~ & ~ & $+$ & ~ & $+$ \\ % = CSS - CS
    \$ & ~ & $+$ & $+$ & $+$ & $+$ & ~ & ~ & ~ & ~ & ~ & ~ & ~ & ~ \\
    \hline
    Early web\footnote{Corresponds to the \emph{C\$SS} style in \citep{fielding_architectural_2000}.}
      & $-$ & $+$ & $+$ & $++$ & $+$ & $+$ & ~ & ~ & ~ & ~ & $+$ & ~ & $+$ \\ % = C$SS
    L & ~ & $-$ & ~ & $+$ & ~ & $+$ & ~ & ~ & ~ & $+$ & ~ & $+$ & ~ \\ % = LS
    COD & ~ & $+$ & $+$ & $+$ & $\pm$ & ~ & $+$ & ~ & $+$ & ~ & $-$ & ~ & ~ \\
    \hline
    LCODC\$SS & $-$ & $++$ & $++$ & $+4+$ & $+\pm+$ & $++$ & $+$ & ~ & $+$ & $+$ & $\pm$ & $+$ & $+$ \\
    U\footnote{Although it is not explicitly included in the original table, \emph{U} has been derived from \citeauthor{fielding_architectural_2000}'s description.}
      & ~ & ~ & ~ & ~ & $+$ & ~ & ~ & ~ & ~ & $+$ & $+$ & ~ & ~ \\ % Uniform Interface (simple, visible, reusable) && en el texto dice que degrada efficiency
    \hline
    REST\footnote{Derived from the addition of \emph{U} to \emph{LCODC\$SS}.} % esto se podría intuír por la línea, pero quién sabe..
      & $-$ & $++$ & $++$ & $+4+$ & $+\pm++$ & $++$ & $+$ & ~ & $+$ & $++$ & $+\pm$ & $+$ & $+$ \\ % = LCODC$SS + U
%	TSC (own) & × & × & × & × & × & × & × & × & × & × & × & × & ×\\ % mejor describirlo sólo con palabras que meterme en berenjenales
    \hline
  \end{tabular}
}{htbp}



\subsection{Properties for \acs{ubicomp}}
\label{sec:middleware_properties}

A middleware is a software layer which provides a higher level of abstraction and masks the underlying heterogeneity.
The middleware presented in this thesis is oriented for \ac{ubicomp} environments and the devices which populate them (i.e. mobile and embedded devices).
Whereas the devices part of the \ac{iot} are a subset of the ones present in \ac{ubicomp}, we consider that the challenges they have to face are the same ones.
For instance, smartphones are not part of the \ac{iot}, but face similar energy and computational limitations as the embedded devices. % posiblemente no todas las propiedades!

% ¿Qué propiedades deseables para IoT añadiríamos?
Therefore, we will adhere to the challenges presented by the \emph{Internet-of-Things Architecture} European project \citep{walewski_project_2011} to analyze our solution.
\citeauthor{walewski_project_2011} state that the \ac{iot} must overcome the following challenges:
interoperability, scalability, manageability, mobility, security and privacy and reliability.
Furthermore, we also consider energetic and computational constraints which devices in \ac{ubicomp} have to face. % ponerlo de forma que sea menos pegote?


% TODO cómo las distintas propiedades vistas antes se combinan para aportar a estas
Table~\ref{tab:middleware_netprop}, shows how the properties explained in the previous sections contribute to achieve these challenges.


% TODO Comprobar que no se salga por el borde derecho!

% Tablita de cómo hereda propiedades de los estilos anteriores?
% (y si quieren más información, que miren en la tesis de Fielding)
\begin{savenotes} % to use footnotes inside
  \begin{table}[htbp]
    \caption{Direct relations between network and coordination properties and challenges for lightweight middleware.}
    \begin{center}
      \begin{tabular}{lcccccccp{0.3cm}ccp{0.3cm}c}
        ~ & \multicolumn{7}{c}{\multirow{2}{*}{Network prop.}}& ~ & \multicolumn{2}{c}{Coord.} & ~ & ~\\
        ~ & \multicolumn{7}{c}{~} & ~ & \multicolumn{2}{c}{prop.} & ~ & ~\\[0.3cm]
        \cline{2-8}\cline{10-11} % column 8 only exists to create that separation
	~ &
	\rotatebox{90}{Performance} &
	\rotatebox{90}{Scalability} &
	\rotatebox{90}{Simplicity} &
	\rotatebox{90}{Modifiability} &
	\rotatebox{90}{Visibility} &
	\rotatebox{90}{Portability} &
	\rotatebox{90}{Reliability} &
	~ &
	\rotatebox{90}{Space uncoupl.} &
	\rotatebox{90}{Time uncoupl.} &
	~ &
	\rotatebox{90}{Semantic Web} \\ % TODO cómo hacer para que esto tenga sentido y esté a la misma altura
	\hline
	Interoperability & ~ & ~ & ~ & × & ~ & × & ~ & ~ & ~ & ~ & ~ & × \\ % ×
	Scalability & ~ & × & ~ & ~ & ~ & ~ & ~ & ~ & ~ & ~ & ~ & ~\\
	Manageability & ~ & ~ & ~ & ~ & ~ & ~ & ~ & ~ & ~ & ~ & ~ & ~\\
	Mobility & ~ & ~ & ~ & ~ & ~ & ~ & ~ & ~ & × & × & ~ & ~\\
	Security &  ~ & ~ & ~ & ~ & ~ & ~ & ~ & ~ & ~ & ~ & ~ & ~\\
	Reliability & ~ & ~ & ~ & ~ & ~ & ~ & × & ~ & ~ & × & ~ & ~\\
	Restricted energy & \multirow{2}{*}{×} & ~ & ~ & ~ & ~ & ~ & ~ & ~ & ~ & ~ & ~ & \multirow{2}{*}{×}\\ % también simplicity
	~~\& comp. & ~ & ~ & ~ & ~ & ~ & ~ & ~ & ~ & ~ & ~ & ~ & ~ \\
	\hline
      \end{tabular}
    \end{center}
    \label{tab:middleware_netprop}
  \end{table}
\end{savenotes}


% TODO ver si tiene sentido hablar de estas cosas
% explicar cómo HTTP trataba de evitar complejidades
% Propiedades que jodes de HTTP de cara al usuario: simplicidad...
% explicar como al crear una capa por encima, complicas todo un poco
%       sencillez
%       user-based library vs a lo otro
% drawbacks frente a HTTP: sencillez, etc.
%A middleware is a software layer which provides a higher level of abstraction and masks the underlying heterogeneity.
%\ac{http} is a middleware indeed, but sacrifices masking exceptions on behalf of the simplicity. % TODO checkear con la tesis de Fielding
% hablar de user-based library vs a lo otro


\subsubsection{Interoperability}

% 2 cosas contribuyen a esta propiedad:
%     Uso de HTTP
%     Semántica => ponerlo como una autonomía (encaja bien en la tablita) o no (en realidad se podría usar tanto sobre TSC como sobre HTTP)?

Interoperability is the key to deal with a wide range of heterogeneous technologies.
The use of both the \textbf{network style} and the \textbf{semantic data} contribute to this challenge.


\paragraph{Semantics}

Using semantics we promote the reuse of data describing them in a more rich and abstract way.
The semantic web is composed by a series of standardized technologies which allow to formally describe the models (i.e. the concepts and how they relate to each other).
Two key mechanisms of the Semantic Web in this aspect are the inference of new data and the mapping of equivalent models.
% hablar de cómo contribuye o no a su portabilidad??
% Besides, the Semantic Web is built on top of well-defined and vastly supported tools.

% Drawbacks
% TODO referencia al paper de qué impide conseguir challenges semánticos


\paragraph{Networking}

For the communication between nodes, we rely on the widely supported \ac{http} protocol. % usar la palabra ubiquitous o pervasive???
Communication with devices using other protocols must be done by means of specialized gateways. % out-of-the-scope
Therefore, adhering to a unique protocol does not help dealing with various communication technologies.

% opción muy común hoy en día para definir APIs
However, a number of current applications use \ac{http} to define their network accessible \ac{api}'s. % through so-called REST APIs
% posiblemente indirectly due to its simplicity => pervasiveness and massive adoption
This makes it a \emph{de facto} requirement to interoperate with these applications. % te da acceso a más apps
% TODO ejemplo de ello: mashups
% también en resource constrained devices: el éxito de WoT es sólo una prueba de ello
In fact, this success has also its counterpart on resource constrained platforms.
The clearer sign for that is the upcoming \ac{wot} initiative.


At a finer grained level, there are several properties which help different \ac{http} implementations to interoperate: % (of interoperability)
\begin{description}
  % \item[Simplicity] => a nivel de red no, pero un middleware simple contribuye a una mayor portabilidad % PROPIEDAD INDIRECTA
  \item[Modifiability] is the ease of gradually introducing changes in the system.
                       Particularly dynamic modifiability avoids restarting the entire system when introducing a change.
                       This property allows different implementations to coexist,
                       reducing issues on the communication between two nodes through \ac{http}.
                       
		      % mejor no meterme en fregados de justificar cada subpunto:
		      % is important due to the \emph{evolvability} and the \emph{reusability}.
		      % Evolvability allows two nodes to have components implemented at different degrees without negatively impacting each other.
		      % contribuyen a eso (sacado de Fielding):
		      %   + CS: The separation of concerns also allows the two types of components to evolve independently, provided that the interface doesn’t change.
		      %   + Layered systems reduce coupling across multiple layers by hiding the inner layers from all except the adjacent outer layer, thus improving evolvability and reusability.
		      % extensibility no: lo introduce COD, así que echa calculos de si me importa esa característica de HTTP
		      % visibility => qué tiene que ver lo fácil que es monitorizarlo con esto?
  
  \item[Portability:] \ac{http}'s portability is backed by the plethora of \ac{http} libraries and frameworks available for most of the computing platforms.
                    % Cachis, no he podido usar una palabra tan molona como: pervasiveness
                    % Justificación: permite interoperar entre distintas plataformas?
                    Consequently, a huge number of platforms can interoperate through \ac{http}.
\end{description}


% No sé dónde meter esta ventaja independiente
% TODO Besides, we enable its \emph{evolvability} in different nodes by allowing them to implement at different degrees (see Section~\ref{sec:status_codes}). % TODO esto es un poco jugón, porque no funcionaría bien sin todo!

Besides, in our \ac{api} \ac{http} the following points remain unsolved:
% estos 2 últimos son como redefinir interop
\begin{itemize} % posiblemente incluir lo que se explica aquí en otro punto
  \item \emph{Using third applications' data}:
  \begin{enumerate}
    \item In Section~\ref{sec:network_properties} we detailed how a pure \ac{rest} \ac{api} contributes to a higher decoupling between clients and servers.
	  This could help to make clients able to autonomously adapt to use different \ac{rest} \acp{api} defined in each device or to different versions of these \acp{api}. % i.e. evolvable
	  Of course, this would come at a cost of more complex clients. % cacharros parseando HTML, interpretandolo, etc.
	  
	  Our \ac{api} is not hypermedia-driven.
	  This requires third applications to implement a uniform \ac{api} on top of them to allow our middleware consume their data.
	  However, the \ac{api} is resource-oriented and simple to implement. % TODO definir qué es y qué no es resource oriented!
	  
    \item Another obstacle for achieving a real application-level interoperability in our middleware is the use of distinct isolated spaces.
	  In other words, two applications using two different spaces would not share their content.
	  Although this could be avoided by defining a default space, as explained in Section~\ref{sec:tsc_soa} their use is justified in terms of scalability.
  \end{enumerate}
  
  \item \emph{Making our data reusable by third applications}:
	
	As explained, we encourage sharing data in resource constrained devices through an \ac{http} \ac{api}. % TODO hacer la distinción de que lo permitimos a través de ambos: 
	Therefore, any other application able to use \ac{http} and manage semantics can reuse data from our middleware.
\end{itemize}


\subsubsection{Scalability}

The high scalability of the chosen architectural network style is negatively affected by the combination of several features:
\begin{itemize}
  \item Data distribution (i.e. accessing to \emph{self-managed graphs})
  \item Lack of hypermedia on the Semantic Web which cannot guide data discovery.
        In other words, a triple location cannot always be discovered following its \ac{uri}. % relacionado con todo aquello de 303 y demás?
        This means that we need an out-of-band searching mechanism as the one presented in the Chapter~\ref{cha:searching}.
\end{itemize}

% TODO TODO TODO resulta que si desacoplas la URI de un espacio de quien la proporciona, el searching mechanism también puede ayudar con eso!!!


\subsubsection{Manageability}

Manageability helps to cope with so many devices allowing them to have an autonomous behavior. % specially in environments where they cannot be controlled through centrality
This autonomous behavior encompasses self-management, self-configuration, self-healing, self-optimization and self-protection.
In this thesis we propose self-managed architecture which eases the searching of self-managed graphs.


\subsubsection{Mobility}

The ability to face mobility situations can be improved through the autonomy brought by the coordination paradigm. % TODO hacer constantes con space y time uncoupling!
First, mobility situations may impede two nodes coexist at the same time.
Thus, \emph{time uncoupling} enables their communication.

The primitives care about the data and not the provider of that data.
Thus, if more than a node provide a content, we make consumers independent of which provides it.


\subsubsection{Security and privacy}
% En una versión nueva hablan de Proof

This property can be considered traversal to any solution and therefore it is considered out-of-the-scope of the thesis.
However, we are currently working on a lightweight security solution \citep{naranjo_lightweight_2012}. % TODO definir mejor esto
For our future work, we are considering to integrate it with the presented middleware.


\subsubsection{Reliability}
% definición de reliability modificada yendo a la fuente de iot-a.eu

Reliability in \ac{ubicomp} is crucial to handle connectivity losses in various \emph{ad hoc}-like ways.
% ver cómo se relaciona con REST
Particularly, \ac{http} handles it through caching.
% Honestidad: This feature has not been properly addressed in our middleware.


\subsubsection{Limited Computation}

% depende de lo limitado que sea el cacharro, puede ser tan así como XXX
Mobile and embedded devices impose computation restrictions.
This restrictions result in unacceptable response times from the application perspective.
Due to \ac{http} simplicity, computation complexity is found in semantic processing.
For example, some really limited platforms may found difficult parsing long files.
However, we made the following assumptions:
\begin{itemize}
  \item Data providers will be able to process wildcard-based templates and generate appropriate responses from their self-managed graphs.
  \item Data consumers are able to process the responses they receive at least in one semantic format. % por ejemplo :-P
\end{itemize}


Beyond that point, our middleware will not impose anything.
This means that the middleware owner cannot assume:
\begin{itemize}
  \item The inference process on data providers.
        Reasoning over the content to provide unstated knowledge must be desirable.
        Still, most of the mobile and embedded device are not able to process big amount of data. % TODO citar WoT2012
        Besides, reasoners are not available in all platforms.
  \item The inference process on data consumers.
        Another alternative or complement for inferring data on the providers is to do it in the consumers after receiving the responses.
        Again, it is a desirable property, but our middleware does not require it.
  \item General use of more complex query languages.
        Their expressibility benefits the network performance since they can restrict the results returned.
        However, parsing them is not be trivial and parsing libraries are unavailable for some platforms.
        Therefore, a provider cannot assume that providers know how to process something beyond the wildcard based templates.
        As an alternative, providers can decompose these queries in simpler wildcard based templates and locally process the results afterwards (filtering them).
        % This is out of the scope.
	% TODO \subsection{Querying Distributedly over the Space} => optamos por no meternos en jardines y up to the middleware implementer to decide each ;-)
	% condiciones relajadas: consistencia, razonamiento distribuído
	% out of the scope: cómo implementarlo
	% optaremos por lo más sencillo
\end{itemize}


% por eso, el mecanismo de búsqueda se hace en los cacharros más potentes
The searching mechanism explained in the Chapter~\ref{cha:searching} takes into account computing limitations.



\subsubsection{Limited Energy}

% 2 cosas contribuyen: computación y networking => TODO mostrar diagrama de FoxG20
Both computing data and networking operations directly affect energy consumption.
As an example, Figure~\ref{XXX} shows the energy consumption on an embedded platform in an inactivity state and networking and computing data. % más detalles en la sección correspondiente
There are some crucial characteristics which negatively affect energy autonomy of resource constrained devices:
\begin{itemize}
  \item Too complex computation tasks.
	 % They result in excessive networking throughput in some occasions.
        As explained in the previous section, semantic processing is the clearest example of that.
        % como se ha explicado: hemos limitado
        We have limited its impact on resource constrained platforms through the searching mechanism.
  \item Regarding \ac{http}, its stateleness contributes to a bad \textbf{network performance}.
        Besides, the verboseness of most of the semantic formats also contributes negatively to that aspect.
\end{itemize}


On the other hand, we contribute to energy saving minimizing network usage (i.e. \emph{network efficiency}) at different levels:
% efficiency: minimizar el uso de la red
\begin{itemize}
  \item \ac{http}'s network \emph{efficiency} can be improved with \emph{code on demand} and \emph{caching}.
  % por eso, el mecanismo de búsqueda se hace en los cacharros más potentes
  \item The searching mechanism improves \emph{recall} and \emph{precision} (see Chapter~\ref{cha:searching}).
\end{itemize}



\subsection{Summary}
\label{sec:middleware_eval_summary}

In the previous sections we have get to the bottom of the strengths and weaknesses of our proposal.
%  Ventajas/inconvenientes frente a WoT:
%      comunicación indirecta para cacharros
%          local+ref
%              de contenidos relativos a coordinación
%              podríamos cachear contenidos de otros en el espacio central (reliability+energy)
%          ref
%              sacar a colación los problemas aquellos de query sobre muchos recursos
As regard of other semantic \ac{http}-based solutions, the solution presented provides a better handling of mobility situations. % e.g. wot
The uncoupling properties which provide it come at the cost of simplicity to implement the middleware.
%   Para TSC (respecto a otras como Smart-M3):
%      limitamos qué gestiona la entidad central => sólo coordinación, el resto lo coge de los cacharros
%      se reconoce su autonomía permitiendo gestionarse su propio contenido => REST, escalabilidad
%      interop => SW y posibilidad de extender a otros servicios REST
%      multiplataforma => facilidad para portarlo
The main strengths of our \ac{tsc} middleware is that it enriches the coordination properties with semantic content provided by embedded and mobile devices.
Also, it pursues interoperability with other \ac{http}-based applications.

Within its main weaknesses we found performance problems and a coupling between applications due to a no-hypermedia driven \ac{api}.
The performance problems may result of handling semantic data and are particularly important for resource constrained devices.
The coupling between applications is solved through a mandatory \ac{http} \ac{api}, but it complicates reusing third applications' data.


% Ventajas de meclar 2 cosas aparentemente separadas?
% what benefit will derive from the integration of both approaches?
The benefits of bringing together \ac{tsc} and \ac{http} is an unified and simple \ac{api} for the developer. % lo de simple no está demostrado
Due to the primitives' simplicity and to the standard and well accepted technologies uses (i.e. \ac{http} and \ac{semantic} standards), this \ac{api} can be ported to a range of platforms.

%: ----------------------- On how to Search in TSC ------------------------

% this file is called up by thesis.tex
% content in this file will be fed into the main document

%: ----------------------- introduction file header -----------------------
\begin{savequote}[50mm]
The logic of validation allows us to move between the two limits of dogmatism and skepticism. 
\qauthor{Paul Ricoeur}
\end{savequote}


\chapter{Searching in a Distributed Space}
% ó Energy-aware Architecture for Information Search in the Semantic WoT
\label{cha:searching}
\newcommand{\pathchapfour}{4_search}


% the code below specifies where the figures are stored
\ifpdf
    \graphicspath{{\pathchapfour/figures/PNG/}{\pathchapfour/figures/PDF/}{\pathchapfour/figures/}}
\else
    \graphicspath{{\pathchapfour/figures/EPS/}{\pathchapfour/figures/}}
\fi


%------------------------------------------------------------------------- 

% comprobar y aplicar los ultimos cambios y ajustes introducidos en el svn de IJWGS!
%\section{Introduction}

% Originally, the Internet was basically composed of a small number of computers that were physically connected to a wired network.
% Over the years, the popularity of the Internet grew and connecting computers became easier and cheaper.
% Thanks to wireless technologies, devices can connect to the Internet without having to be connected to a network.
% 
% Nowadays, everyday objects like cars or washing machines start to be connected to the Internet to exchange information.
% This is what is currently known as the \emph{Internet of Things (IoT)} \citep{atzori_internet_2010}.
% We can organize the IoT into \emph{Spaces} or islands which cover a particular knowledge as proposed by \citet{abdulrazakCGF10}.
% Using Spaces, we can ease the creation of different \emph{Ambient Intelligence (AmI)} environments.
% For example, we can create a Space for our home where we connect our devices to form an AmI environment.
% Thus, the washing machine could plan to finish washing the clothes two minutes before the car gets home.

Integrating mobile or embedded devices is not trivial as they usually communicate using different protocols.
To solve this problem, the \acl{wot} initiative proposes to use well-established web standards to ease their communication.
However, the format of the data they exchange is also multifarious and application domain dependent.
This implies that data will not be meaningful in other domains unless a specialized system converts and reinterprets them.
A way to solve this problem is annotating the data semantically as proposed by the \ac{www} \citep{ChuaG10,kimKC11}.

Adding semantics to the \ac{iot} works well for devices with high computational capacity but it adds too much overhead for most of the devices composing the \ac{iot}.
To reduce this overhead in such devices, part of this computation is usually delegated to intermediaries \citep{honkola_smart-m3_2010}.
This approach reduces the overhead of semantically annotated data but brings other problems.
(1) When devices rely on others to provide information, it is not guarantee that the information accessed will accurately represent the last information available in the data providers (e.g., the sensors).
(2) Once we rely on those intermediaries, it is required that they are available at all time.
Otherwise, the devices would not be able to talk to each other.

We propose a solution where we use intermediaries to release some workload from the less powerful devices, but at the same time, we promote the direct communication between the devices.
In particular, our system uses intermediaries to search where the data is located in the \Space{} and queries the final devices directly.

In our solution, the devices can become or stop being intermediaries dynamically.
To decide which device will be an intermediary, we evaluate the state of a device (i.e., energy and computation capacity).
Using this dynamic architecture, the absence of a particular intermediary does not collapse the system.
In addition, our system is flexible enough to support a wide range of scenarios.

To enable the search for information in this system, intermediaries need to know the information available in the \Space{}.
To do that, they aggregate summaries sent by each device.
We propose alternatives to summarize and to aggregate this information taking into account the payload of the shared information and the accuracy of the search.

We compare our approach against other common approaches.
In particular, we focus on the energy aspect, demonstrating that our approach helps energy constrained devices to face fewer unnecessary requests.
We also evaluate other important metrics like the number of messages a device has to exchange to perform a request, and the accuracy of the search results provided by the intermediaries.
We perform this evaluation under different scenarios to prove the flexibility of our proposal.

In summary, we make the following contributions:
\begin{itemize}
\item We present a new architecture for managing semantics on the \ac{wot} which reduces the load for resource-constrained devices.
\item We propose and evaluate different approaches to aggregate and summarize the semantic information available in the \Space{} and enable semantic search in this architecture.
\item We demonstrate the advantages of our approach compared to other typical searching approaches under different scenarios.
\end{itemize}

The remainder of this chapter is organized as follows:
Section~\ref{background} discusses the related work.
Section~\ref{proposal} presents in detail our energy-aware architecture.
Section~\ref{sec:clues} describes the information that devices exchange to maintain this architecture.
Section~\ref{environment} presents our experimental environment and Section~\ref{results} evaluates our solution.
Finally, Section~\ref{conclusion} states the conclusions of our work.
\section{Background}
\label{background}

Querying over the semantic content provided by independent sources transparently is a problem which often appears in the \acf{lod}.
\citet{gorlitz_federated_2011} classify the possible \ac{lod} infrastructures according to
\begin{enumerate*}[label=\itshape(\arabic*\upshape)]
  \item how they store data,
  \item whether the index used to search is distributed, and
  \item whether data sources cooperate.
\end{enumerate*}
Table~\ref{tab:infrastructure_lod} summarizes the resulting infrastructure types: central repository, federation and \ac{p2p} data management.


\InsertTab{tab:infrastructure_lod}{Infrastructure paradigms according to their characteristics \citep{gorlitz_federated_2011}}{}{
  \begin{tabular}{l|c|c|c|c|}
    \cline{2-5}
    ~ & \multicolumn{2}{c|}{Central Data Storage} & \multicolumn{2}{c|}{Distributed Data Storage} \\
    \hline
    \multicolumn{1}{|l|}{Independent} & \multirow{2}{*}{n/a} & ~ & \multirow{4}{*}{Federation} & \multirow{2}{*}{n/a} \\
    \multicolumn{1}{|l|}{Data Sources} & ~ & Central &~ & ~ \\
    \cline{1-2} \cline{5-5}
    \multicolumn{1}{|l|}{Cooperative} & \multirow{2}{*}{n/a} & Repository & ~ & P2P Data \\
    \multicolumn{1}{|l|}{Data Sources} & ~ & ~ & ~ & Management \\
    \hline
    ~ & Distr. Index & \multicolumn{2}{c|}{Central Index} & Distr. Index \\
    \cline{2-5}
  \end{tabular}
}{htbp}


% Lo nuestro se parece más a LOD de dispositivos que a distributed triple store:
%     grafos semánticos de fuentes de datos independientes
The solutions which distribute the semantic content are federation and \ac{p2p} data management.
While the latter distributes the index to look for content in multiple machines, federation does not.
%The difference among them resides in whether the index to look for content in the different machines is distributed or not consecutively.
%A federated \ac{lod} infrastructure maintains the meta information about data sources in the machine in charge of searching (called \emph{federator}).
%On the other hand, the P2P approach maintains the meta information distributed.
Our solution proposes a federation infrastructure with some particularities:
\begin{itemize}
  % Esto es, los índices están replicados, pero quien los usa siempre los coge de la misma máquina.
  % No se necesita coordinación entre quienes guardan los índices para balancear o obtener algunos índices que no tienen...
  \item There is a unique node in charge of managing the main version of the metadata needed to create indexes.
	However, it is dynamically chosen among all the participants and can change over the time.
	Consumers hold copies of the manager's metadata.
  \item Each consumer builds its own index from its metadata replica.
	This allows consumers not to critically depend on the availability of any other node to search. % no depender => MUCHO
	%In other words, in our solution there are as many \emph{federators} as data consumers. % Ya no puedo poner esto, porque no he explicado qué es un federator
\end{itemize}


\bigskip


% presentamos 3 soluciones: FedX, ARQ y XXX.
FedX \citep{schwarte_fedx_federation2011,schwarte_fedx_optimization_2011}, DARQ \citep{quilitz_querying_2008} and SemWIQ \citep{langegger_semantic_2008} are the most relevant solutions to perform federated queries over independent semantic data sources.
%Both in these solutions and in ours relevant sources are know beforehand (i.e. they use a \emph{top-down} strategy).
%\emph{Bottom-up} strategies discover data sources by following the links on the content provided by other sources during the query process \citep{schwarte_fedx_optimization_2011}.
%Although our solution also uses a \emph{top-down} strategy, we explicitly avoid human intervention to discover new data sources or providers. % a diferencia de FedX, por lo que ví en la demo % o si, en realidad no nos metemos mucho con eso
%Instead, we rely on automatic discovery mechanisms which additionally must provide some additional information about the data providers. % BTW, esto es un problema para interoperabilidad!
% esto está sacado sobre todo del estado del arte de schwarte_fedx_optimization_2011
DARQ and SemWIQ require locally preprocessed metadata about data sources, while Fedex uses on-demand queries together with a caching mechanism. % SPARQL ASK
In that aspect, our solution resembles DARQ and SemWIQ since it also requires this metadata (called \clues{} in our solution).
In fact, we propose two types of \clues{} which are equivalent to the metadata used by DARQ and SemWIQ.
DARQ maintains predicate indexes to find relevant data sources.
SemWIQ maintains a local catalog with the type information of RDF entities provided by data sources. % el tipo de cada sujeto se debe conocer de alguna forma
In contrast with our solution, DARQ and SemWIQ also use statistics about data sources to optimize their queries.


%    por lo tanto se podría implementar sobre lo nuestro lo suyo
%            (y permitir optimizaciones si introducimos SPARQL en algunos nodos)
% no hay ninguna solución global con una arquitectura que tenga en cuenta le energía
% cómo el problema difiere: más fuentes de datos, menos info en cada una y dinamicas
A key difference with the mentioned systems is the context where we have applied our solution.
It is intended to serve in \ac{ubicomp} environments.
This means that there will be many more data sources and with less information to process in each of them.
Furthermore, their nature also differs from the usual \ac{lod} endpoints' one: the mobile and embedded devices which may provide data are less reliable, more dynamic and more restricted in computation and energy.
Therefore, we designed and evaluated our solution considering this nature and restrictions.
To the best of our knowledge, this is the first approximation to the problem within the explained context.
% además, abarcamos más parte del proceso?? => discovery, la replica, etc.


% pros y contras:
Considering the constrained architecture of the data sources, we acknowledge that some of them will not be able to process \acs{sparql} \citeweb{sparql2008}.
Consequently, we require the lowest common denominator: basic triple patterns.
This causes the main weakness of our solution compared to FedX, DARQ and SemWIQ: a less sophisticated querying mechanism.
%This is because all these solutions assume that all data sources can process SPARQL and offer optimization techniques on top of it. % todo cita a SPARQL
%In any case, SPARQL can always be decomposed into triple patterns.
For our future work, we are considering using \acs{sparql} just with the endpoints able to manage it as a halfway solution.
This will allow us to introduce some of the optimization techniques presented by other solutions.


\bigskip


% mencionar que se ha hablado de soluciones IoT que usen semántica anteriormente y que to the best of our knowledge no existe nada parecido
% como se ha visto en TAL y TAL sección no existen soluciones de busqueda de contenidos semánticos distribuido para WoT. (o al menos no para funcionar en cacharrillos)
In sections~\ref{sec:sw_intermediaries} and \ref{sec:sw_providers} we have analyzed other solutions for \ac{ubicomp} which use semantics.
However, none of them have addressed federated distributed search in resource constrained devices.
% meter alguna referencia a Cabilmonte y demás?
% Es más de cómo consultar datos semánticos en streaming a una base de datos relacional...


% Mencionar SPARQL SERVICES y demás ??
% No es transparente, pero permiten hacer queries a distintas fuentes.


% mencionar fokoue con lo de los resumenes ?
% meter enlace con lo de summaries de ABox en la parte correspondiente?






% La acepción de Name resolution cambia depende del contexto!
% http://en.wikipedia.org/wiki/Name_resolution

% Name resolution vs recognition
% http://stackoverflow.com/questions/8589005/difference-between-named-entity-recognition-and-resolution

% Mirar en el punto 6.3 para ver unos pasos de consulta en la LOD que parece que pueden ser interesantes para esto
% http://www.unifr.ch/webnews/content/20/attach/4765.pdf
\section{Energy-aware super-peer architecture} % pensar titulo más sexy
\label{proposal}

% In this section, we discuss the content the clues should contain to make them suitable to the \ac{wot};
% we explain in detail the roles the nodes of our solution can have;
% we overview the deployment issues related to these roles, such as discovery or role transition;
% and we discuss how our solution can adapt multiple scenarios.

% AG: ¿node, peer y/o device? ¿Aclarar qué es un node?
% IG: personalment peer no lo veo, y entre node y device, pues creo que usarlos indistintamente, pero aclarando eso en algun punto
% AG: Crear comandos para asegurarnos de llamar a las cosas igual siempre (e.g. clues, acronimos, etc.)
% IG: comandos? esa es muy de crack xD

% En algún momento de la sección: CoAP vs HTTP
% AG: Igual habría que comentar que asumimos HTTP por su estado de madurez, pero que podría aplicarse a CoAP en cuanto fuese un estandar y con implementaciones que lo implementasen por completo. Esto tendría ventajas a nivel de uso de Multicast para agrupar WP, consumidores, etc.
%     Podríamos valorar su uso más minuciosamente en cada sección, pero igual en una primera versión no le daría mucho bombo.
% IG: influye en algo de la solucion es algo particular del experimental environment?
% AG: En su mayor parte no, pero habría algunas cosas que se simplificarían (e.g. comunicación con un grupo de nodos)

%   + Vinculo rapido con la intro (1 o 2 frases) % igual también con el related work
In a semantic \ac{wot}, nodes are part of a network where they share semantically described information.
%   + Mencionar el uso de espacios
These nodes gather their information in \Spaces{} which are useful to create \ac{ami} environments.

%   + Introducir ejemplo de hotel
For example, a hotel can create its own semantic \Space{} with the semantically described information of its services and the data provided by devices (e.g., sensors) spread around the hotel.
Thus, clients can use their own personal devices (e.g., smartphones) to interact with the devices of the hotel and search useful information.
Using this approach, a client, for example, can check the current swimming pool occupancy rate using his personal mobile phone.

However, not all devices can afford the costs of fully managing the additional overhead of semantics.
In this section, we present our architecture to manage such environments and deal with the limitations of resource-constrained devices.

\subsection{Basic roles}
As we introduced before, devices can be part of one or multiple \Spaces{}.
A device can provide data (\providers{}), consume data (\consumers{}) or both at the same time.

\begin{description}
\item[Providers:]
This is the simplest role which any device in our system can carry out, even the smallest sensor.
\providers{} must manage their own semantic information.
In particular, they must organize triples in RDF graphs.
\item[Consumers:]
This is the role a device must take to get information from the system.
These devices need to be able to use semantics to find the data they require.
\end{description}

Note that a device could hold just one or both roles at the same time or switch between them.
For example, a \provider{} could become a \consumer{} to get data from the system or a \consumer{} could stop asking for data.
% Cambios de roles:
% To Consumer: try to read something from the network. => evidente
% To not-consumer: "t" without querying anything (it does not worth to update the clues if we won't query).


\subsection{Use of Intermediaries: \acl{wp}}
%   + Intermediaries
%   + Necesidad de tener white pages, mención de roles. => subsección N
Tasks involving the use of semantics can be expensive for some devices.
To reduce the load on such resource-constrained devices, we need to use intermediaries to carry out some of these tasks.
% No hacen tareas por ellos, más bien les evitan hacerlas.
%   + Por que queremos que ellos administren su info?
However, we do not want devices to completely rely on intermediaries for three reasons:
(1) data consumers need to get fresh data and we can only get this by direct access to the actual \provider{},
% We aim to promote the direct communication between devices to ensure that freshness of the information shared.
(2) we must support mobile scenarios where personal devices carry their own data,
%   + Explicar un poco la naturaleza dinámica de los escenarios (sensores que entran y salen del sistema)
% In particular, we consider the dynamicity of WoT environments, where nodes can join and leave a space making their information available or unavailable at any time.
and (3) we need to reduce the maintenance complexity (nodes may join or leave at any time).

For these reasons, we cannot rely on intermediaries to host and provide \emph{all} the semantic information.
However, we use intermediaries as searching enablers in the semantic \Space{}.
This kind of intermediary is what we call \acl{wp}.

\acp{wp} manage metadata or \clues{} about the information shared by others.
The \consumers{} rely on \acp{wp} to enhance the search process and reduce the number of requests generated.
Consequently, the \providers{} process less semantic data and reduce their overall overhead.
Thus, our architecture enables an intermediary-aided energy-aware information search.


\InsertFig{architecture}{fig:solutionArchitecture}{
  Role of White Pages in our proposal
}{}{0.8}{}


Figure~\ref{fig:solutionArchitecture} shows the role of \acp{wp} in our proposal.
In particular, \acp{wp} manage \clues{} which summarize the semantic information shared by the \providers{}.
These \clues{}, which are described in detail in Section~\ref{sec:clues}, are pieces of information useful to determine which nodes can answer to a certain query.
\consumers{} use these \clues{} to directly access the semantic information on the \providers{}.
Summarizing, the main tasks of a \ac{wp} are:
\begin{itemize}
 \item Manage the \clues{} sent by all the nodes in the system.
 \item Aggregate clues sent by \providers{} in a unique response.
 \item Reply to requests from \consumers{} with aggregated clues or a list of nodes (details in Section~\ref{sec:interacting}).
\end{itemize}
 % debería haber un único \emph{WP} o varios? AG: Optaría por 1 para simplificar.
 % Dependiendo de eso, puede ser más complejo mantener una visión común de los clues de todos
 % (con algún algoritmo de gossiping?)
 % Para simplificar el merging de los differentes gossipings agregados se podría hacer: 1 master, N slaves.

\acp{wp} only store \clues{} for a period and they expire afterwards.
Using this approach, we avoid storing information from no longer available nodes.

% Si hubiese que dar detalles de implementación, había pensado en Redis para el \emph{WP} a poner en un server o cluster
% De forma que la gestión distribuida de dicho gossiping no fuese problema mio en esos casos
% Veo 2 casos:
%    + Un cluster de servidores hacen de WP: que usen algo que ya exista y me deje en paz.
%    + Dos/n móviles hacen de \emph{WP} conjuntamente: entonces sí que debería mojarme y decir como mantienen esa visión
%      conjunta.
% El problema es que al haber cambiado la explicación de tipos de cacharros a roles, esta consideración es difícil de
%  hacer.


\subsection{Versioning Clues}
As previously explained, the \ac{wp} aggregates clues sent by \providers{}.
This aggregation is versioned using a \emph{version number} and a \emph{generation number}.
The \emph{version number} increases each time the \ac{wp} receives a new clue.
The \emph{generation number} ($g_{id}$) is a timestamp that represents the moment a new \ac{wp} was chosen.
This requires the clocks to be synchronized in all the nodes to avoid problems.
In any case, the \ac{wp} must guarantee that the $g_{id}$ is higher than the one used by a previous \ac{wp}.

The \ac{wp} maintains two version numbers:
(1) the one used during the setup to do the first load (\emph{setup version}) and
(2) the version of the last aggregated clue.
The first one is shared through the discovery mechanism.
The second one is shared with both \providers{} and \consumers{} in \acs{http} responses.
The versioning is used to improve the \ac{wp} selection process explained in Section~\ref{sec:selection}.


\subsection{Discovering a \acl{wp}}
% He matizado la frase para que no se entienda que los Consumer SIEMPRE dependen de los WP. Sólo para actualizar clues.
% AG: te has empeñado en quitar el matiz de que "sólo" al principio está vendido a la existencia de un WP :-)
In our architecture, when a \consumer{} joins the \Space{} it relies on a \ac{wp} to find information.
Hence, the first thing a node needs to do afterwards is discovering who is the \ac{wp} in that \Space{}.

To run our proposal, we require a discovery system able to
(1) get the \Spaces{} that a particular node belongs to,
(2) identify the \ac{wp} in the system and its \emph{setup version}, and
(3) provide additional information about nodes to decide which one can be the next \ac{wp}.
To implement our architecture, the information about the nodes must include:
% Estos parámetros se vuelven a explicar en la sección de "Selección de WP"!
battery level, available memory, storage space and the approximate time passed since the node joined the \Space{} for the last time (to estimate its reliability).

How to discover the nodes in the \Space{} is transversal to this paper.
We could either extend approaches like \ac{upnp}, \ac{mdns}\footnote{\url{http://tools.ietf.org/html/rfc6762}} or lmDNS \citep{jara_light-weight_2012} or use \acs{http}, \ac{coap}\footnote{\url{https://datatracker.ietf.org/doc/draft-ietf-core-coap/}} and \ac{dns} together as proposed by \citet{ishaq_facilitating_2012}.
In particular, in Section~\ref{sec:mdns} we use \ac{mdns} to evaluate our solution.
% En la primera ronda de revisiones hubo un revisor que dijo:
%  "How does the system receive the information about the battery level, available memory?
%   The authors justify that this is out the scope of this paper, but it is important to know
%   the potential overhead introduced in these processes."
% Es una forma de responderle. Con esa frase quería decir:
%  " Esto no solo funciona en el plano teorico, usando un sistema como mDNS es perfectamente
%    posible implementarlo porque compartir esa pizca de info, pese a que te llegue un poco
%    desactualizada, no es un problema. "

% - Which is the topology of the network?
% It is hard to evaluate the proposed solution in terms of overhead and energy consumption for performing necessary tasks (WP discovery, WP selection, clues updates,...)
% without any knowledge of the network. Are there wired/fixed nodes? Are WP infrastructural nodes? Are the wireless nodes connected through multihop paths?
Note that we assume that any node in the \Space{} can reach the other nodes.
In particular, we assume IP addressability (no matter whether the node is wired or wireless).
As a consequence, to use devices from Wireless Sensor Networks (e.g., Zigbee or 6LoWPAN) as \providers{}, one should rely on gateways.

% IG el CoAP multicast no tiene una utilidad diferente a lo q se habla?
% Se podría usar UPnP, mDNS o  incluso lo propuesto por Jara (lmDNS).
% También se podría usar multicast CoAP (e.g. los que respondan a /\ac{wp} o /consumers)


% Partimos de la base de que el sistema de descubrimiento nos ha descubierto a otros nodos y los siguientes datos acerca de ellos:
%    - Battery (B): valor [1 a N] y unidad [min, hora, dia] -> siendo -1 infinito o cargando.
%         Problema: puede que no sea fácil hacer ese cálculo.
%         Solución: mejor eso a carga relativa o medida de carga, cuando no sabes cuanto consume el cacharro de normal.
%    - Memory (M): Valor [1 a 1023] y unidad [KB, MB, GB]
%    - Persistence (P): Valor [1 a 1023] y unidad [MB, GB, TB]
%    - Joined since (JC): Tiempo desde su unión al espacio (para permitirnos medir su estabilidad): medido en ciclos de media hora
%    - Stores aggregated clue (SAC): [T o F] ¿Tienen una versión del clue agregado?
%
% Hemos tratado de no crear mucho tráfico en el mecanismo de descubrimiento almacenando datos que varíen con poca frecuencia.
% Los campos más variables: la batería y JC, pueden no actualizarse en media hora sin que el algoritmo se vea sustancialmente alterado.
%
% Comprobar que no son datos que costase mucho embeber en cualquier sistema (registro TXT, lo que sea que use UPnP...)
% Jara habla de 80 bytes en 6LoWPAN, probablemente tampoco nos orientemos a cacharros tan peques, pero oye, por comentarlo...


\subsection{Interacting with a \acl{wp}}
\label{sec:interacting}
When a node needs to access a piece of information, it asks the \ac{wp} for the information about the other devices in the \Space{}.
Using this information, the node can find the owner of the required data and ask it for this information.

In our architecture, \providers{} and \consumers{} have additional duties.

\noindent\textbf{Providers.}
They manage their own semantic information and generate \clues{} about the information they host.
Then, they send these \clues{} to the \Space{} \ac{wp}.
As a response, they receive the last version of the aggregated clue they have contributed to.
\providers{} will send their \clues{} to the \ac{wp}
(1) every time a \emph{clue} is updated,
(2) before the lifetime of a \emph{clue} expires, or
(3) whenever there is a new \ac{wp} in the \Space{} with a lower \emph{setup version} than the one in the \provider{}.
Note that \clues{} do not change frequently since they represent the type of information the nodes host rather than the specific data which is constantly generated.

\noindent\textbf{Consumers.}
When a \consumer{} needs to get information from the system, it first needs to find the \ac{wp} in the \Space{} and use this \ac{wp} to obtain an aggregated clue of all the nodes.
Then, it processes this aggregated clue to decide which are the nodes to query.
Finally, it sends the query through \acs{http} requests.

They perform this process synchronously for the first query and periodically in an asynchronous manner for the following ones.
% cada cuanto? arbitrario? que lo decida en base a parámetros de funcionamiento? => Mi propuesta abajo
This period should have an upper limit to ensure a fresh view of the \Space{} and a lower limit to guarantee that the \ac{wp} is not flooded.
To adjust the update frequency within these two limits, we evaluate the frequency of the last 10 requests to that node.
Thus, the view of the \Space{} will be fairly up to date when the \consumer{} processes the next query.
% TODO Añadir mejora: cuando se detecta un nuevo nodo por el mecanismo de descubrimiento, esperas un poco y actualizas clue para ver si ya ha enviado su info al WP.
% Mejora de la mejora: exclusivamente preguntas por la clue de el/los nuevo/s nodo/s detectado/s. Asi envias un JSON más peque.

\noindent\textbf{Optimizations for resource constrained devices.}
Note that \providers{} will update the \ac{wp} with new information and \consumers{} will periodically check the \ac{wp}.
Using this approach, the \ac{wp} reduces the network load and in particular, it decreases the load for resource constrained devices.

In addition, \emph{aggregated clues} can be too long for some really limited devices.
As an optimization for such devices, the \ac{wp} is able to answer specifically the node to address a query.
Thus, these nodes will only maintain a list of the nodes they should ask for several predefined queries.

% Limitar la información enviada también a otros nodos?
% Si tenemos 1024 nodos y cada gossiping es de 100KBs, vamos a tener 100MBs de gossipings :-S
% Aunque esto es justo lo que hay que discutir en el apartado de evaluación de clues...


\subsection{Selecting a \acl{wp}}
\label{sec:selection}
%   + Que queremos solucionar => subsección K (tipos de sistemas)
Depending on the setup, having a dedicated \ac{wp} (e.g., a high-end server) can be too expensive.
For example, in a domestic environment, it is not worth dedicating a full server to be a \ac{wp}.
However, when we have a large setup (e.g., a hotel), it is convenient to dedicate a few servers to decrease the load in small devices.
%   + Necesidad de que esos roles sean algo dinámico y el problema que ello presenta
%      (descubrimiento, pero sobre todo selección) => subsección M

We have to be able to manage the complexity of a system composed by heterogeneous devices.
For example, making a small device a \ac{wp} may be inappropriate in highly populated environments.
On the contrary, having multiple dedicated servers implies a high management overhead which is unnecessary in simple environments.
Our architecture is flexible and any node can be a \provider{}, a \consumer{}, a \ac{wp} or any of them at the same time.

\noindent\textbf{When does the \ac{wp} selection start?}
The selection process can start (1) when no \ac{wp} is available or (2) when the current \ac{wp} gets a worse score than other nodes.
In the first case, the first node to realize there is no \ac{wp} starts the selection process.
In the second case, the current \ac{wp} periodically checks if there is another node with a better score.
If this actually happens, it starts the process.
We limit the frequency for this checking to avoid the overhead associated with this change.

\noindent\textbf{How to select a \acs{wp}?}
The node in charge of selecting a \ac{wp} retrieves information of available nodes using the discovery mechanism.
It ranks them, selects the top one, and informs it that it should become a \ac{wp}.
If a selected node rejects the new role, the selector chooses the next one on the ranking.

The node in charge of the selection sends its aggregated clue to the new \ac{wp}.
If this aggregated clue is fresher than the one in the \ac{wp}, the \ac{wp} will use this one.
Otherwise, the \ac{wp} will use its own.
In case there is no previous aggregated clue, the new \ac{wp} will start from scratch (version to -1).
Once the node becomes the new \ac{wp}, it notifies the version used to initialize the discovery system.
Doing so, \providers{} will realize if they must send their \clue{} again because it is not included in the aggregation provided by the \ac{wp}.

% aclarado que facilitará esa información de cada nodo activo en el espacio
Note that the discovery mechanism must provide the following information about each node in the \Space{} to the selection algorithm:
(1) memory of the device,
(2) storage capacity of the device,
(3) time since it joined the \Space{}, and % esto no debería cambiar a menudo para no generar mucho trafico
(4) its battery charge.
We use the algorithm detailed in Listing~\ref{list:selectionAlgorithm} to rank the nodes.

\begin{listing}
  \begin{algorithmic}[1]
  \State $nodes \gets filter_{threshold}(nodes, ``memory'', threshold_{memory})$
  \State $nodes \gets filter_{threshold}(nodes, ``storage'', total_{nodes} \times storage\_needed_{avg})$
  \If { $anyWith(nodes, battery_{infinite})$ }
    \State $nodes \gets filter(nodes, ``battery'', battery_{infinite})$
    \State $nodes \gets orderBy(nodes, ``battery'')$
  \Else
    \If { $anyGreaterThan(nodes, ``joined\_since'', joined_{threashold})$ }
      \State $nodes \gets filter_{score}(nodes, ``joined\_since'')$
    \EndIf
    \State $nodes \gets filter_{score}(candidates, ``battery'')$
    \State $nodes \gets orderBy(nodes, ``memory'')$
  \EndIf
  \State \Return $nodes$
  \end{algorithmic}
  \caption{\acl{wp} selection algorithm.}
  \label{list:selectionAlgorithm}
\end{listing}

Lines 1 and 2 of the algorithm apply filters based on thresholds.
In particular, the second threshold is variable and depends on the number of nodes (more nodes will generate more \clues{} to store).
After that, we check the power availability of the devices.
We first select the nodes connected to the power grid (represented by $Battery_{infinite}$).
If there are devices connected to the electrical grid, we select devices which are steady enough to apply a filter by $joined\_since$ field.
Filters in lines 8 and 10 choose nodes with z-scores higher than one for the specified fields.
If no node fits that filter, it returns the nodes with values higher or equal to the mean.

% AG: Intentando aclarar si un móvil puede ser WP o no (a un revisor no le había quedado claro)
Summarizing, the algorithm prioritizes energy autonomous nodes and within the nodes with battery limitations, it prefers steady ones.
% IG pues si quieres aclararlo, remata xD
In both cases, it finally selects a node with enough storage and memory (including mobile devices).

% AG: Problema a mencionar: y si no hay ningún \ac{wp} candidato? Broadcasting entre cacharros?

% \noindent\textbf{\ac{wp} conflict resolution}
Due to out-of-dated information, two nodes may become \acp{wp} at the same time.
Our solution would eventually correct this situation thanks to a conflict resolution algorithm.
When a \emph{\acs{wp} A} detects another \emph{\acs{wp} B} in the \Space{}, it has to check which one has a better score according to Listing~\ref{list:selectionAlgorithm}.
If \emph{B} has a better rank, \emph{A} simply resigns as \ac{wp} and notifies it to the discovery system.
Otherwise, it forces \emph{B} to resign through a \acs{http} invocation.
Other nodes will be aware of these changes through the discovery mechanism.




\section{Shared Clues}
\label{sec:clues}
As we introduced before, \acp{wp} host \clues{} about the information in the \Space{}.
Using a \emph{clue}, a \consumer{} can find which node (or nodes) is the \provider{} of a piece of information.
\providers{} generate these \clues{} by digesting the semantic information they store.

Thanks to these \clues{}, resource constrained devices do not have to process unnecessary requests.
We also limit the length of the \clues{} to reduce the bandwidth, the memory, and the storage overhead on these devices.

In this section, we describe in detail these \clues{}, the information they contain and their format.


\subsection{Querying Basics} % Querying assumptions

% Por qué una nueva sección?
% Porque estaba quedando muy abstracto todo y cosas básicas como cómo un consumidor consulta a un proveedor podía no quedar claro.
% O qué es un template y para qué se usa.

A \consumer{} directly queries the information to the selected \providers{}.
This selection is done using the \clues{} described in the rest of the Section.
However, in this section we explain how \consumers{} directly query for information to the selected \providers{}.

We assume that \providers{} offer a HTTP endpoint which accepts queries.
Whether this endpoint redirects the requests to the RDF graph which describes a specific resource,
it creates a response with RDF triples belonging to different resources or
it offers hyperlinks to the relevant graphs is transversal to our solution.
The design of this endpoint is also out of the scope of this paper, but it should be the same in all the providers.

% explicar query templates
To perform the queries, which enable the selection of a subset of the semantic content hold in a given \Space{}, we require a \textbf{template}.
We assume the use of wildcard templates, which are special triples with optional wildcard subject, predicate and/or object (see Figure~\ref{fig:semanticExample}). % TODO referencia al ejemplo o a figura de ejemplo
We could use more sophisticated query languages (e.g. SPARQL) % TODO footnote?
but processing them may be too demanding for devices with constrained capabilities.
On the other hand, simple wildcard templates can be managed by any node able to manage RDF triples.
In any case, complex queries can be decomposed into wildcard templates and our solution would still apply.


\InsertFig{semantic_example}{fig:semanticExample}{
  Sample triples and query templates
}{
  The top of the figure shows some sample triples.
  It highlights the different parts used for each type of \clue{}:
  the information used in class-based \clues{} in orange, the prefix-based \clues{} in blue, and the predicate-based \clues{} in brown.
  The bottom shows query template examples for these triples.
}{0.8}{}



\subsection{Content of a Clue}
To find out which is the most appropriate solution for \ac{ami} scenarios, we have to consider scenarios populated by mobile devices and sensors. % for \ac{iot} scenarios?
Mobile devices usually share data which rarely changes and is described using a few ontologies (e.g., the user profile and his preferences).
On the other hand, sensors are constantly generating new instances of the same ontology (also called individuals).
In both cases, the data shared by each node is described according to one or few vocabularies or taxonomies.

At this point, it is important to define the \emph{TBox} and \emph{ABox} concepts following the definition of \citet{nardi2003introduction}.
\emph{TBox} contains knowledge describing general properties of concepts or terminology and
\emph{ABox} contains knowledge specific to the individuals of the domain of discourse.
An example of \emph{TBox} information is the device type or the elements it is made of
while \emph{ABox} can describe the mobile phone brand or the temperature sensed by a thermometer.

Each mobile device or sensor usually generates \emph{ABox} according to the same \emph{TBox}.
Using this information, we propose avoiding the use of URIs which represent \emph{ABox} information in general terms.
Specifically, we propose and evaluate three possible type of \clues{} to share.
%The first two are inspired by early works on peer-to-peer semantic web $[CITAEDUTELLA,devadithya2007index]$.

\medskip

\noindent\textbf{Prefix-based clue.}
%A coarse-grained step to find relevant nodes for a query is to ignore those which do not have \emph{ABox} for a specific \emph{TBox}. % no sólo aplicable al caso de la ontologias
The individuals (\emph{ABox}) are described according to different ontologies (\emph{TBox}).
Each ontology usually employs a common \emph{namespace}.
This means that the URIs of the concepts in these ontologies share a unique \emph{prefix}.
Besides, each \provider{} also uses few \emph{prefixes} to generate the URIs of its individuals.
In this type of clue, we propose a coarse-grained method to filter nodes which do not have individuals with the prefixes used by the query template.
To create prefix-based clues, the \providers{} have to extract the prefixes used in their graphs. % obvio?

Following the example in Figure~\ref{fig:semanticExample}, a \provider{} \emph{P} may have that graph.
Thanks to the \clues{}, the \consumer{} knows that \emph{P} has URIs starting with \emph{sweet}, \emph{ssn}, \emph{ex} and \emph{weather}.
Consequently, it will send to \emph{P} a query which uses any of the templates shown in Figure~\ref{fig:semanticExample}.
On the contrary, if the template is \emph{?s pref2:sth ?o}, it will avoid sending the template to \emph{P}.


\medskip

\noindent\textbf{Predicate-based clue.}
The predicates relate subjects with other subjects or literals.
These predicates are defined in the \emph{TBox} (e.g., to state they relate concept \emph{A} with concept B) and used in each triple of the \emph{ABox}.
In this approach, we propose extracting the set of predicates used in the graphs stored by each node.
Using this information, they can be simply matched with the predicate defined in the query template.
% AG: Posible pega, se ha explicado antes lo que es un "template"?
% IG: nope, de hecho en la anterior tb me ha surgido la duda de q era

\medskip

\noindent\textbf{Class-based clue.}
In the third approach, we propose sharing the classes of concepts (\textit{rdf:type}) provided by the nodes.
Using this information and the \emph{TBox}, each node can check if the information matching certain query template is susceptible to be stored in other nodes.
For this kind of clues, we assume that each node should have (or be able to obtain) the \emph{TBox} related to the template used for querying.
This is a reasonable assumption since ontologies describing \emph{TBox} are usually accessible on the URL described by its prefix.

To understand how it works, consider that, according to the \emph{TBox}, a predicate \emph{p} relates the concept \emph{A} with another concept.
We assume that the nodes storing instances of the class \emph{A} are more likely to use this predicate in their graphs.
Therefore, if a query template uses that predicate, we will send it to all nodes storing instances of the class \emph{A}.


% IG: esto de referirte a una ontologia del futuro... no me parece muy asines
% AG: Exactamente mismo caso que arriba, pero con ejemplos concretos.
%     La idea era que no quedase muy abstracto. Lo comento porque ahora puede quedar redundante.
% For instance, consider that the following template is queried: \textit{?s~ssn:observes~dbpedia:CO2}.
% According to the SSN ontology, which will be presented later on, only an instance of the class \textit{ssn:Sensor} \textit{observes} something.
% Therefore, the querying node will send the template to the nodes which have instances of the class \textit{ssn:Sensor}.


\subsection{Reasoning to expand clues} % Esto se relaciona principalmente con class-based, pero igual también puede comentarse para Predicate-based.
Through a reasoning process one can know unstated knowledge.
Using this knowledge, we can detect more relevant nodes.
For example, let us consider the previous example where the class \emph{A} defines the predicate \emph{p}.
Providing that \emph{C} is a subclass of \emph{A}, a node which has many instances of the class \emph{C}, may also use the predicate \emph{p}.
Thanks to the reasoning, we can discover that the node has knowledge of the type \emph{A} and therefore, it may be relevant.
% AG: Comentado por lo mismo que arriba
%(in the example, instances of \textit{ssn:Sensor} class' subclasses such as \textit{Accelerometer}).
The drawback, is that reasoning consumes a lot of resources.
This limitation will be further analysed in Section~\ref{sec:clues_eval}.
% TODO Many authors solve this limitation by transforming the queries in the consumer side \cite{}.


\subsection{Use of ABox in clues}
\label{sec:aboxinclues}
% se podrían usar individuals que no cambian mucho
% location, individual en base al cual se generan todos los elementos.
As stated before, in general terms, we want to avoid the use of \emph{ABox} URIs (individuals) in our clues.
Thus, we fulfil two goals:
(1) we generate smaller \clues{} and
(2) the \clues{} will not change too frequently and therefore, we will require less communication to update clues.
However, in some cases, the use of \emph{ABox} content in the \clues{} may be beneficial.
For instance, assume a URI that refers to the specific location \emph{L}.
If we want to search for devices in location \emph{L}, we cannot deduce anything about it using the proposed \clues{}.

% AG: Atencion, a partir de aqui viene una propuesta que me acabo de sacar de la manga que habra que discutir.
%     El problema de esto es que no está en la evaluación y para el 15 de Octubre, lo veo ajustado de incluir.
For this reason, we need to consider sharing the \emph{N} most queried individuals in our clues.
To do that, the \ac{wp} needs to store a list with the statistics about the information collected by each \consumer{}.
\consumers{} can send this information together with the request to update a clue.
\providers{} can obtain a list of the current most popular URIs before sending their updated \clues{} to the \ac{wp}.
Using this list, \providers{} can know if they have these URIs and include them in the \clue{} to be sent to the \ac{wp}.
% TODO el proceso no me acaba de convencer, y cuando hay una nueva URI? Como saben los proveedores que deben actualizarla?
Note that this process would imply an extra request per \provider{} before each update.

This simple approach implies sending not only \emph{TBox} but also \emph{ABox}.
The amount of extra information added to each \clue{} will depend on the size of this list (\emph{N}).
The effectiveness of this method will depend on the number of queries using one of the \emph{N} URIs in their subjects or objects.
% Propuesta: uso de un porcentaje de las queries realizadas
% Problema: no vas a raspar a cada providers con 1000 comprobaciones tampoco!
% Solución: ¿merece la pena evaluar eso? ¿como?


\subsection{Format}
% Tras explicarlas, poner algunos fragmentos con las clues (JSON) para que se vea de que hablamos
We can use many formats to represent the content of a clue.
% pongo esto de EXI porque se que a la gente del mundillo de IoT se le hace el culo chupicola con sus estandares
One option is the ongoing Efficient XML Interchange (EXI) \footnote{\url{http://www.w3.org/XML/EXI}}. % o cite?
EXI is designed to efficiently interchange XML data and therefore, we could obtain better compression rates than with JSON. % aseveracion cierta o solo para XML?
However, we have chosen JSON for its simplicity and its wide adoption in the WWW.

% TODO no lo he acabado de pillar...
The prefix-based \clue{} is the easiest one to represent in JSON since it is formed by a set of URIs.
We can also represent the predicate-based and the class-based \clues{} using a set of URIs.
However, the prefixes of those URIs are usually repeated and for this reason, we do not use plain URI transmission for these clues.
We show an example in Listing \ref{list:oneClue}.
It first defines the prefixes used and gives them a name and then, it specifies the URI endings for each prefix.

\begin{listing}
  \begin{minted}[frame=single, framesep=3mm, linenos=true, xleftmargin=21pt, tabsize=4]{js}
  {
    "s": [[
      "so",
      "http://knoesis.wright.edu/ssw/ont/sensor-observation.owl#"
    ]],
    "p": {
      "so": ["result", "procedure",
		  "observedProperty", "samplingTime"]
  }  }
  \end{minted}
  \caption{
    Representation of a predicate-based \clue{} in JSON.
    The node sending the \clue{} has RDF triples which use the predicates
    \emph{so:result}, \emph{so:procedure}, \emph{so:observedProperty} and \emph{so:samplingTime}.
  }
  \label{list:oneClue}
\end{listing}

\medskip

These are isolated \clues{} sent from a \provider{} to the \ac{wp}.
However, the \ac{wp} gathers all these \clues{} in an aggregated clue which is sent to the \consumer{}.
We show an example of an aggregated clue in Listing \ref{list:aggregatedClue}.
As one can see, the aggregated clue defines the type of \clues{} wrapped through the numeric field \emph{i}.
\emph{G} and \emph{v} form the version of this aggregated clue.
The field \emph{s} defines the prefixes.
Finally, for each node, each prefix is related with the URI endings.

% TODO sacar un ejemplo real? poner esquematicamente?
\begin{listing}
  \begin{minted}[frame=single, framesep=3mm, linenos=true, xleftmargin=21pt, tabsize=4]{js}
  {
    "i": 1,
    "g": 2435467,
    "v": 556,
    "s": [
	  ["dc", "http://purl.org/dc/elements/1.1/"],
	  ["dul", "http://www.loa.istc.cnr.it/ontologies/DUL.owl#"],
	  ["ssn", "http://purl.oclc.org/NET/ssnx/ssn#"] ],
    "p": {
      "node1": {
	"ssn": ["observedBy", "observationResult"],
	"dul": ["isClassifiedBy"] },
      "node0": {
	"ssn": ["observes"],
	"dc": ["description"]
  } } }
  \end{minted}
  \caption{
    Representation of an aggregated clue in JSON.
    Line 2 defines that it embeds predicate \clues{} (i.e. type 1).
    Lines 3 and 4, contain the version of the aggregated clue.
    The remaining lines express the predicates used by two nodes.
    For example, \emph{Node1} has at least a RDF triple which uses the predicate \emph{ssn:observedBy}.
  }
  \label{list:aggregatedClue}
\end{listing}

% Mencionar: Note that we can gossip sleeps periods to know when to query them...
% Comentar que se podrían gossipear los sleeps de los nodos embebidos para saber cuando interrogarles.
% AG: Esto es un aspecto que puede ser guay comentar, de cara a IoTificar la propuesta.
%     Pero pensandolo mejor, debería ir en la sección 3, a la hora d explicar el formato de las
%     clues.
\section{Experimental Environment}
\label{environment}
We have used simulation to study the performance of our solution and compare it against a flooding-based one.
Using a simulation, we can evaluate multiple scenarios and repeat experiments for different approaches under the same conditions.

\subsection{Methodology}
Table~\ref{tab:configurationParameters} shows the main parameters of our simulator.
We vary these parameters to simulate a wide range of scenarios.
% Dado que los revisores parece que no hicieron cado de este aviso, mejor quitarlo.
% We omit topology considerations as we assume devices have IP connectivity.


\begin{table}
  \centering
    \begin{tabular}{l p{7cm}}
      \hline
      Name & Description \\
      \hline
      Network size & The number of nodes in a network. In the simulations conducted, all the nodes are \providers{}. \\
      Number of writes & Amount of writes performed during the simulation period. \\
      Number of queries & Amount of queries performed during the simulation period. \\
      Number of \consumers{} & Amount of nodes querying to other nodes in the \Space{}. \\
      Distribution strategy & Our solution or \ac{nb}. \\
      & In \ac{nb} nodes write locally and spread the queries to the rest of the nodes in the \Space{}. \\
      Drop interval & At the beginning of this interval a node abruptly leaves the network. \\
      & At the end, the node joins the network and a new one is chosen to leave it. \\
      % FUTURE WORK: añadir métrica de accuracy [ http://en.wikipedia.org/wiki/Recall_(information_retrieval) ]
      \hline
    \end{tabular}
    \caption {Configuration parameters.}
  \label{tab:configurationParameters}
\end{table}

% Since we wanted to simulate the Triple Space Computing paradigm over HTTP, the communication between the nodes was point to point and the data exchanged were RDF Triples.
% The node discovery process was ignored since it will show similar additional overhead for each strategy.
% It can be considered transversal to what was being measured.
% IG: version del parrafo anterior
% AG: Ok, pero esta es una de las cosas que en el uptade debería volar o cambiar. Sobre todo, porque no hablamos de TSC
%     en ningún otro lado. Puede ser más fácil aún, ya que como asumimos que queremos semántica en WoT, el HTTP está
%     justificado per se.
As we simulate \ac{http}, we assume a point to point communication between devices which exchange \ac{rdf} Triples.
We discuss how the discovery process affects the solution in Section~\ref{sec:dynamic}.
% TODO ENERO
In the rest of the sections we omit the node discovery process as it represents the same overhead for all strategies.

%To represent the data managed by each node, we first considered using LUBM\footnote{\url{http://swat.cse.lehigh.edu/projects/lubm/}}, a synthetic benchmark.
%Unfortunately, it creates instances from very few classes for each node, which makes all the nodes to have the same TBox.
%In our opinion, this does not faithfully represent Internet of Things scenarios with heterogeneous devices sharing disparate information.

%In our second attempt, we found a dataset which could represent this heterogeneity.
To represent the data managed by each node, we used data from diverse sensor network environments.
These data follow the \textit{Semantic Sensor Network Ontology} (SSN)\footnote{\url{http://www.w3.org/2005/Incubator/ssn/XGR-ssn-20110628/}}.
% created by the W3C Semantic Sensor Network Incubator Group. to represent diverse sensor network environments.
SSN has been used in many projects and scenarios to describe semantically the data provided by heterogeneous sensors.
% TODO IG hasta aqui hemos llegado
Specifically, we used data from the following datasets:
AEMET metereological dataset\footnote{\url{http://datahub.io/dataset/aemet}},
University of Luebeck Wisebed Sensor Readings\footnote{\url{http://thedatahub.org/dataset/university-of-luebeck-wisebed-sensor-readings}},
\emph{Kno.e.sis} Linked Sensor Data\footnote{\url{http://wiki.knoesis.org/index.php/LinkedSensorData}}
and Bizkaisense\footnote{\url{http://helheim.deusto.es/bizkaisense/}}.%http://dev.morelab.deusto.es/bizkaisense
These datasets contain descriptions about the sensing stations and the data sensed by them during certain periods.
The analogy between stations which have different sensors and the \ac{iot} devices is reasonable.
The datasets have been adapted to provide just one measure of each sensor at each moment (to emulate the storage restrictions from embedded devices) and to use as many stations as nodes has the network (depending on the network size).

However, sensors can not only be found in AmI environments.
Personal devices such as mobile phones usually populate them also.
To represent such circumstance, we added semantic data of people to represent their profiles\footnote{\url{http://www.morelab.deusto.es/joseki/articles}}. % TODO TODO TODO actualizar con la nueva plataforma que pongan!

% poner disponible el código
We use SimPy\footnote{\url{http://simpy.sourceforge.net}} to simulate each scenario.
SimPy is a process-based discrete-event simulation language for Python.
To accurately simulate the time needed by each node to provide a response, we considered measures taken from
real embedded web servers running on a \textit{ConnectPort X2} IP gateway\footnote{\url{http://tinyurl.com/connectportx2}} for \textit{Digi}'s \textit{XBee sensors}\footnote{\url{http://tinyurl.com/xbee-sensors}} (\textit{XBee} from now on),
on a FoxG20\footnote{\url{http://www.acmesystems.it}} and on a Samsung Galaxy Tab.
Besides, we also provide measures taken from a regular computer.
Table~\ref{tab:measures_embedded} shows the measures used for the parametrization.

% TODO TODO TODO poner o referenciar las características de los dispositivos cogidas de WoT 2012

\begin{table}
  \begin{center}
	\begin{tabular}{p{2.5cm} r r r r}
	  \hline
	  & \multicolumn{4}{c}{Devices} \\
	  \cline{2-5}
	  Concurrent & \multirow{2}{*}{XBee} & \multirow{2}{*}{FoxG20} & Regular  & Samsung \\
	  requests   &  ~    &   ~     & computer & Galaxy tab \\
	  \hline
	  1  &  77 (1)	&  17 ~(0)  &   5 ~(1)  &  223 (349) \\
	  5  & 392 (8)	&  97 (16) &   8 ~(4)  &  256 ~(76)  \\
	  10 & 775 (8)	& 174 (28) &  13 ~(8)  &  372 (171) \\
	  15 &  -   	& 282 (43) &  18 (13) &  497 (191) \\
% Tampoco creo que haga falta tanto detalle, ¿cuando se van a dar tantas peticiones concurrentes?
%	  20 &  -	    & 375 (30) &  23 (13) &  661 (444) \\
%	  25 &  -	    & 460 (30) &  30 (18) &  748 (288) \\
%	  30 &  -	    & 540 (35) &  38 (22) &  929 (805) \\
%	  35 &  -	    & 632 (29) &  38 (20) & 1029 (672) \\
	  \hline
	\end{tabular}
  \end{center}
  \caption{Mean of the measurements taken in different devices with the standard deviation in parenthesis (milliseconds).}
  \label{tab:measures_embedded}
\end{table}


\subsection{Performance Metrics}
To evaluate the fundamental properties of the strategies, we use the following metrics:
\begin{itemize}
  \item \textit{Precision}: the fraction of nodes which answered relevant results (those responses which were not \textit{not found} responses).
                            It measures the exactness of the results.
  \item \textit{Recall}: the fraction of relevant answers that are returned.
                         It measures the completeness of the results.% overload, failed search, etc.
  \item \textit{Size}: the size of each type of clue.
%  \item \textit{Throughput}: the average rate of successful message delivery over a communication channel.
%  \item \textit{Idle time}: the average time each node is in an idle state and therefore is consuming less energy.
%  \item \textit{Failed requests}: the amount of request which could not be answered due to the physical limitations of
% the nodes (i.e. both due to timeouts or connection rejections when the server is overloaded).
  \item \textit{Total requests}: the number of \acs{http} requests performed during a simulation.
  \item \textit{Response time}: the average time needed to obtain an \acs{http} response.
  \item \textit{Active time}: the total time spent by each node either querying other nodes or handling a query.
\end{itemize}
\section{Results}
\label{results}

\subsection{Types of Clues Shared}
\label{sec:clues_eval}
As presented in Section~\ref{sec:clues}, the type of clue used will affect
(1) the \emph{precision} and \emph{recall} to find the nodes with the appropriate information; % explained in section 4
and (2) the amount of information to transfer over the network (both requests and responses) and nodes have to process.
Increasing \emph{precision} reduces the number of unsuccessful requests to handle and thus, it reduces the energy consumption. 
Similarly, sending more information over the network implies more processing time and more energy consumed.

\medskip

\noindent\textbf{Precision and recall.}
We evaluate the \emph{precision} and the \emph{recall} of the proposed algorithm in a network of 470 nodes issuing the query templates shown in Table~\ref{tab:evaluationTemplates}.
In average, the nodes manage instances belonging to 6.34 different classes ($SD=1.31$) among a total of 17 distinct classes in the space.
%When these classes are expanded with a reasoning process, the nodes store 20 different concepts among a total of 113 classes in the space ($SD=4$). % TODO actualizar
The distinct predicates managed by each node in average are 16.01 ($SD=1.53$) out of 68 different predicates in the space.

\begin{table}[h!tbp]
\centering
\begin{tabular}{| c | c |}
\hline
Name & Template \\
\hline
t1 & ?s~~~rdf:type~~~ssn-weather:RainfallObservation \\
t2 & ?s~~~wsg84:long~~~?p \\
t3 & ?s~~~ssn:observedProperty~~~?o \\
t4 & bizkaisense:ABANTO~~~?p~~~?o \\
t5 & ?s~~~dc:identifier~~~?o \\
%t9 & ?s~~~foaf:family_name~~~?o \\ % TODO futura version del paper con perfiles de usuarios, para meter mas variabilidad en contenidos
\hline
\end{tabular}
\caption{Templates used in the evaluation.}
\label{tab:evaluationTemplates}
\end{table}
% Las consultas se podrían sacar también de bizkaisense o alguna app que hayamos hecho para darle mayor verosimilitud

% Explicar class based
In Figures~\ref{fig:recall_measures}~and~\ref{fig:precision_measures}, class-based clue shows a good \emph{precision} and \emph{recall} for \emph{t1} and \emph{t2}.
\emph{T1} asks exactly for the information this type of clues define (i.e. nodes having instances of a certain class).
\emph{T2} evaluates which nodes have instances in the domain of the \emph{long} predicate (\emph{SpatialThing}).
Note that this works thanks to the RDFS inference because some nodes in the space only write \emph{Point} instances (a subclass of \emph{SpatialThing}).
The domain of \emph{t3} and \emph{t5}'s predicates could not be inferred just using RDFS inference. % (some properties of OWL are used => inverseof).
Even solving this limitation, we would expect a bad \emph{precision} since both predicates relate very general concepts.
In addition, when a class-based clue has no enough information to predict the nodes, it simply floods the query.
This is why the \emph{recall} of \emph{t4} is high.



\InsertFig{clues_recall}{fig:recall_measures}{
  \emph{Recall} for each type of clue used.
}{
}{1}{}

\InsertFig{clues_precision}{fig:precision_measures}{
  \emph{Precision} for each type of clue used.
}{
}{1}{}
% IG: TODO mencionar en el texto de referencia o en la caption algo como: The higher the better.

% Explicar predicate-based
We can see a bad prediction for \emph{t1} and \emph{t4} for predicate-based clues.
\emph{T1} defines a very common predicate and therefore, it cannot discriminate any possible node.
\emph{T4} suffers the same problem explained for the class-based clues.
We proposed a possible solution for this problem in Section~\ref{sec:aboxinclues}.

% Explicar schema-based
Finally, schema-based clue shows a slightly better \emph{precision} for \emph{t4}, since it can discriminate some nodes not using the \emph{bizkaisense} prefix.
On the other hand, it obtains marginally worse \emph{precision} than predicate-based clues for \emph{t3} and \emph{t5}.
This worsening could be greater if few nodes using the schemas \emph{ssn} and \emph{dc} will use the predicates defined in both templates.

\medskip

\noindent\textbf{Verbosity.}
% Una frase para retomar lo anterior, y al grano.
The clues verbosity is also a critical aspect for resource constrained devices.
Figure~\ref{fig:clueSize} shows a higher variance for schema based clues' length and lower verbosity of class based clues.
This is because the nodes virtually have a different number of sensors.
In addition, the links to concepts of other ontologies varies within the two datasets used in the parametrization.
In any case, the diagram shows a similar verbosity for all the clues for the semantic content considered in this evaluation. % TODO analizar si la media varía significativamente
% No es adecuado para cacharros pequeños: poner evaluación de WoT de inferencia


\InsertFig{clues_length}{fig:clueSize}{
  Length of the clues alternatives
}{
}{0.7}{}

% Añadir nuevas barras: Predicate+schema y predicate+schema+MostCommonsIndividuals
% Poner otra barra para saber cuanto contenido semántico guarda un nodo de media en el experimento?
% Comentar que cabe en MTU de ethernet y UDP en caso de querer enviarlo por CoAP


% TODO Añadir nuevo diagrama para el gossiping agregado y ver como crece a más elementos añadidos

\medskip

\noindent\textbf{Summary.}
Class-based clues are useful for templates asking for a specific type of content.
However, they still require inference to obtain a good \emph{precision}.
In \citet{gomez-goiri_restful_2012}, we tested the inference process on the devices and data used in this simulation. % ref a WoT2012 para más detalles
We could not run any reasoner in the ConnectPort X2 Gateway.
Actually, we could only run RDFS reasoners in more powerful embedded and mobile devices such as the FoxG20 and the Samsung Galaxy Tab.
In the FoxG20, it took 48.9 seconds the first load of all the ontologies used and 1.4 seconds to reason over each measurement written.
In a Samsung Galaxy Tab, it took 17.3 seconds and 0.2 the following measurement writings.
Considering these results, we can conclude that there is a clear need for efficient embedded reasoners.
Therefore, the class-based approach is promising but it is impossible to adopt in current embedded and mobile devices.

Between the predicate-based and schema-based clue, we propose to use the predicate-based clue since it subsumes much of the information provided by the schema-based clue.
The rest of the schemas, are referred in the subjects or the objects. % más que schemas son prefixes :-S
They could be easily added to predicate-based clues on the prefixes field.
In addition to the use of predicate-based clues, we could implement the solution for the specific individual search proposed in Section~\ref{sec:aboxinclues}.

% Posible TODO
% GRAFICO 3: \emph{precision} y \emph{recall} comparando con y sin de cada uno

% GRAFICO 4: tamaño de gossiping individual y agregado comparando con y sin cada uno
%            (barra con barra superpuesta encima con cuanto más añade)
%            Según hay más nodos, cómo crece el gossiping a manejar?


% Discussion: ¿elegir un tipo u otro dependiendo del modo de operación?
% (i.e. si hay que perder precisión a costa de no intercambiar MBs...)
% AG: Interesante, pero yo no complicaría una sección ya de por si bastante liosa.


% TODOs importantes que me gustaria hacer para la siguiente version:
%   - Ver cómo crecen los clues agregados.
%   - Proponer una clue híbrida que mezcle a las anteriores.
%   - Evaluar de alguna forma la mejora propuesta para ABox.



\subsection{Network usage} % IG: network role
\label{sec:NetworkUsage}

% Donde ``explicar'' negative broadcasting y/o centralizado? En seccion 4?
% Explicar porque no se pone centralizado en la comparación
%    No es directamente comparable dado que depende directamente de otro factor distinto: frecuencia de escritura.

We have conducted a simulation to evaluate the benefits of our solution against a flooding-based approach (i.e., \emph{negative broadcasting}).
The simulated environment has nodes that join to the same space.
All the nodes are \emph{Providers} and write new information periodically in the space.
During one hour, 1, 10 and 100 \emph{Consumers} perform 1000 queries in total using the templates described in Table~\ref{tab:evaluationTemplates}.

As expected, our solution scales much better than the \emph{negative broadcasting} (see Figure~\ref{fig:requestsByStrategies}).
In particular, the increase on requests per consumer in the space is not appreciable.


\InsertFig{requests_by_strategies}{fig:requestsByStrategies}{
  Required requests for negative broadcasting (\emph{nb}) and our solution with 1, 10 and 100 consumers
}{
}{0.7}{}


In Figure~\ref{fig:requestsByRoles}, we take a closer look to the origin of the traffic of our approach in a space with 100 consumers.
The communication between the \emph{Providers} and the \emph{WP} is much more infrequent than the other communication types.
The reason is that writing to a node only results in a clue update when the structure of the managed information changes.
The first time the metadata about the node (sensor) is written, the second time the first measure and following writings, just add or replace a measure.
Therefore, the clue does not change after the second step.
This matches with the assumption made to share \emph{TBox} information in our clues.

The communication between \emph{Consumers} and \emph{WP} is in between the other two communication patterns.
It is greater than the one from \emph{Providers} to \emph{WP} because \emph{Consumers} need to maintain an updated view of the space.
Recall that the update time depends on the query frequency of each \emph{Consumer}.
The maximum and minimum updating frequency were set to 10 and 1 minute(s) respectively.

The communications between \emph{Consumers} and \emph{Providers} is most of the total communications.
This shows that the overhead added by the use of \emph{WP} on our solution is not significant and it is justified by the reduction of the total number of communications shown in Figure~\ref{fig:requestsByStrategies}.


\InsertFig{requests_by_roles}{fig:requestsByRoles}{
  Requests between roles in our solution in a space with 100 consumers
}{
}{0.7}{}



\subsection{Energy consumption}
\label{sec:energyConsumption}
% Idea: Lo de arriba está muy bien, pero específicamente, cómo afecta a los cacharros?
Our solution tries to save energy by making \emph{Providers} handle fewer requests from \emph{Consumers}.
These savings contrast to the overhead added by the communication with the \emph{WP}.
However, our results have demonstrated that this overhead is small in comparison to the total number of communications.

The energy consumption in mobile and embedded devices increases each time a device needs to process something or communicate with another node (see Figure~\ref{fig:energy_consumption}).
To analyse how communications impact their energy autonomy, we have to consider not only the number of communications but also their time length (see Table~\ref{tab:measures_embedded}).
For example, a mobile phone will consume less energy asking clues to a server than asking them to an embedded device as it has to wait longer for the response.


\InsertFig{energy_consumption}{fig:energy_consumption}{
  Energy consumption for FoxG20 during different activity periods (miliwatts)
}{
}{0.8}{}


The experiment consists of 300 nodes joined to a space running on 1 server, 30 galaxy tabs, 75 FoxG20 and 194 Xbees. % TODO marca del xbee!
We increase the number of devices as their price and capacity decrease.
Using this approach, we mimic a typical space where cheap devices are more common.

As shown in Figure~\ref{fig:activity_measures}, our solution reduces the activity of each device by more than 5 times compared to \emph{negative broadcasting}.
The diagram on the right details the average activity for each type of device.

In our solution, we can check how the load moves from the embedded devices (XBee and FoxG20) to the server (which is indeed chosen as a \emph{WP}).
The exceptional activity registered by the Galaxy Tabs is caused by their extremely high response time.
% This implies that handling a request in the Galaxy Tab takes much longer than the average.
However, future work will cover the change of the HTTP library used in our Android implementation to reduce this response time.


\InsertFig{activity_measures}{fig:activity_measures}{
  Activity time for each strategy
}{
  In the first diagram we appreciate the average active time a node spends in each strategy. In the second one, the active time is classified by the type of device each node is running on.
}{0.8}{}
% TODO añadir a la derecha los valores en NB?

%%%%%%%%%%%%%%%%%%%%%%%%%%%%%%%%%%% QUIZA para 2da VERSION %%%%%%%%%%%%%%%%%%%%%%%%%%%%%%%%%%%%%%%%%
%  + Vamos a medir y comparar distintas situaciones entre sí:
%   Nota: cuando hablo de servidor, movil o dispositivo embebido, ejemplifico capacidades.
%         dispositivo embebido puede almacenar menos cosas y hace consultas más concretas.
%
%      - Situación 1: Negative broadcasting (o todo dispositivos embebidos)
%      - Situación 2: 1 servidor (WP), 20 moviles, 120 dispositivos embebidos (proporción 1:20:120)
%      - Situación 3: 20 moviles, 120 dispositivos embebidos (proporción 1:6, entre móviles 5 con el cargador enchufado)
%      - Situación 4: todo dispositivos embebidos
%  + Eje Y: Tiempo medio en ejecución
%  + Eje X: A parte de las situaciones (ver esquema de \emph{recall}) algún otro aspecto que también afecte al consumo de energía:
%      - número de consultas?
%      - frecuencia de las mismas?
%      - nodos consultores?



\subsection{Performance in Dynamic Environments}
\label{sec:dynamic}
We have evaluated the network usage of our solution in ordinary situations in Sections~\ref{sec:NetworkUsage} and \ref{sec:energyConsumption}.
Nevertheless, we have not conducted simulations where the nodes frequently join and leave the space.
In such situation, the communication needed to manage the clues might be a burden.


\InsertFig{dynamism}{fig:dynamic_situations}{
  Effects of dynamic scenarios in our solution
}{
  Note that the last interval in the x-axis represents a simulation with no drops.
}{0.8}{}


To assess the effect of  dynamic networks on the performance of our solution, we used the scenario presented in Section~\ref{sec:energyConsumption},
Then, we simulate nodes joining and leaving the space at different intervals: 30 seconds, 1 minute, 5 minutes, 10 minutes, 20 minutes, 30 minutes, 45 minutes.
Particularly, for our solution, we tested the most harmful situation: the node leaving the space abruptly is always the \emph{WP}.
We also added an scenario with no drops as a baseline.
Note that we represent this scenario by configuring the drop-interval with a greater value than the simulation time.

In Figure~\ref{fig:dynamic_situations}, we see the results of these simulations.
These results show that even in such dynamic situations, our solution requires fewer communications than negative broadcasting.
In our solution, most of the communications are between \emph{Consumers} and \emph{Providers}.
To evaluate the overhead added by our solution, the right graphic shows the communications involving \emph{WPs}.

We can appreciate that the updates on \emph{Consumers} are independent of the number of times the \emph{WP} changes.
We can also see a minimal change between the scenario where the \emph{WP} is always available and the one with 5 minutes drop-interval.
In this case, when the \emph{WP} drops, many \emph{Consumers} have the latest version of the aggregated clue.
% Dado que se parte del que tenia el WP o del que le ha dado el consumer que le ha elegido. Ambos tienen más opciones de tener una versión actualizada.
This situation increases the chances of getting an updated version for the initialization of the new \emph{WP}.
Thus, it reduces the number of messages from the \emph{Providers} to the new \emph{WP}.


\subsection{Effects on Discovery Mechanisms}
\label{sec:mdns}
We want to prove the feasibility of our solution using a common discovery mechanism.
To that end, we simulate the behavior of the \emph{Multicast DNS} (mDNS)\footnote{http://tools.ietf.org/html/draft-cheshire-dnsext-multicastdns-15}
and \emph{DNS service discovery} (DNS-SD)\footnote{http://tools.ietf.org/html/draft-cheshire-dnsext-dns-sd-11} protocols.
Both protocols are based on the well known and widely accepted Domain Name System protocol (DNS).

On the one hand, DNS-SD proposes the use of the data field of TXT records to share key/values.
We use these records to share the information needed by the selection algorithm among the nodes.

On the other hand, mDNS defines how this and other records are shared through UDP multicast (or unicast in certain situations).
We ignore the cost of browsing the nodes to discover new nodes because is equal for both strategies.
% TODO ENERO IG este frase del however no enlaza demasiado bien con la anterior
However, as explained above, our strategy does differ from \emph{Negative Broadcasting} in the use of TXT records.

The nodes announce records during the start up or whenever they have a resource record with new data.
Therefore, each time a record is updated, we send a multicast message that increases the network traffic.
In our solution the TXT record may change 
(1) when a new \emph{WP} is selected or
(2) when we update the time elapsed since it joined the space and its battery charge level.
The last two parameters need to be updated to select an appropriate \emph{WP} but they do not need to change too frequently.

% Regarding the first case,  % no pillo la introduccion del first case
In the most static scenario the TXT record is only written once.
The more dynamic scenario from the previous section, one on the contrary, updates that record 126 times after writing it for the first time.
This demonstrates that the overhead generated on the discovery system by our solution is minimal even in the worst-case scenario.

\section{Conclusion}
\label{conclusion}

In this chapter, we have presented a dynamic architecture to enhance the search of semantic contents in the \acl{wot}.
%  by any node part of the Web of Things.
In particular, this architecture chooses an intermediary according to its capabilities to support resource constrained devices.
Intermediaries can use different types of clues to summarize the information in the semantic \Space{}.
% managed by the nodes belonging to a space.
Thanks to this support, the devices can directly interrogate others to obtain fresh information while reducing their semantic overhead.

% Esto lo he cambiado porque Ipiña decía que habría que volver a aclarar porqué lo he comparado con flooding en la conclusión (por si alguien no se ha leído el paper ;-))
The main characteristic of our solution is the ability of the devices to share semantic data directly, no matter how simple they are. % o algo como guides our design
Under this assumption, a flooding-based strategy would obtain a high recall. % or a high fraction of relevant answers (por si no se entiende aún que es recall)
However, our evaluation shows that our solution requires less messages between devices than a flooding-based strategy (i.e. is far more precise).
% he puesto more para que no tenga la acepción peyorativa de "many" (no se requieren muchísimos más para que sea mejor)
Even if we use caching strategies to alleviate these effects, our solution performs better for scenarios with more data consumers and query types.
In addition, our approach reduces the workload of mobile and embedded devices which indirectly results into energy savings.

% Para trabajo futuro, se podría considerar: 
% query optimization y citar a alguna solución de distributed triple stores
For our future work, we will consider allowing more expressive query languages such as SPARQL. % use acronym
This would force \consumers{} to discriminate between the \providers{} able to process it and the rest.
%The latters, could implement a trade off solution to send multiple triple pattern templates at once. % optimizando la red
On top of these queries, we could use query optimization techniques in the \consumers{} \citep{schwarte_fedx_optimization_2011}.
Using these techniques, we could transform the queries before sending them to a) obtain more results or b) use the network more efficiently.


% ----------------------------------------------------------------------



%: ----------------------- On how to actuate through TSC ------------------------

% this file is called up by thesis.tex
% content in this file will be fed into the main document

%: ----------------------- introduction file header -----------------------
\begin{savequote}[50mm]
Gonna act like a man or a bitch?!
\qauthor{Det. Alonzo Harris}
\end{savequote}


\chapter{Remotely Actuating}
\label{cha:actuate}
\newcommand{\pathchapfive}{5_actuate}


% the code below specifies where the figures are stored
\ifpdf
    \graphicspath{{\pathchapfive/figures/PNG/}{\pathchapfive/figures/PDF/}{\pathchapfive/figures/}}
\else
    \graphicspath{{\pathchapfive/figures/EPS/}{\pathchapfive/figures/}}
\fi


%------------------------------------------------------------------------- 

In the previous chapter, we presented a energy aware technique to search on the completely distributed semantic space.
%This search mechanism promotes the end-to-end interaction between the objects.
However, as important as observing what happens in an \ac{ubicomp} environment, is to act over it.
Although this dissertation describes less thoroughly this problem, we considered interesting to discuss about it not to leave loose ends. % transmit a broader vision. % the big picture


Therefore, this chapter presents and compares two techniques to change the environment by \emph{writing} into the space. % no se si se entiende que la segunda también emula "escribir en el espacio"
The first one is based on \ac{ts} patterns and uses tasks and a subscription mechanism.
The second one starts with a goal and knowledge about the environment.
Then, it reasons to create a plan to fulfill the goal.
This plan comprises calls to \ac{rest} services.
%These services are semantically annotated and are read from the space.


% ----------------------------------------------------------------------



\section{Actuating through the Space}


% explicar brevemente patrones
% requirements: sistema de suscripciones (out of the scope)
% escenario de ejemplo en Ubicomp


In space-based computing, participants coordinate by reading and writing in a shared Space. % TODO usar define a nivel general para Space!
This encourages an uncoupled communication between applications using the same Space.

% IG: esto lo he homogeneizado porque tenia 3 estilos distintos :)
This section analyzes how to achieve this uncoupled communication following a top-down approach.
First, Section~\ref{sec:ts_patterns} briefly explains \acl{ts}'s most common application patterns. % according to the classification of \citet{freeman_javaspaces_1999}.
Second, Section~\ref{sec:envisaged_scenarios} uses some of these patterns to show how to build an \ac{ubicomp} scenario.
The scenario shows some requirements the middleware must fulfill.
Finally, Section~\ref{sec:notification} analyzes these requisites.



\subsection{\aclp{ts}' patterns}
\label{sec:ts_patterns}

According to \citet{freeman_javaspaces_1999}, there are four main application patterns which can be used with \ac{ts}: % IG: with o in?

\begin{description}
  \item[Replicated-worker pattern.] In this pattern, there is a master process and many worker processes able to compute the same task.
				    %% REDUCIENDO esto: %%
%				    The master takes a problem and divides it in tasks which are solved by any of the workers.
%				    More concretely the pattern is composed by the following steps:
% 				    \begin{enumerate}
% 				      \item The master divides a problem up into smaller tasks.
% 				      \item It writes them in the space.
% 				      \item Any of the many workers takes a task.
% 				      \item That worker computes the task.
% 				      \item It writes the result of that computation.
% 				      \item The master collects the results for the tasks it wrote.
% 				      \item Once the master has all the results, it combines them into a meaningful overall solution.
% 				    \end{enumerate}  
				    The master takes a problem, divides it into smaller tasks, and writes them tasks into the space.
				    Any available worker takes a task, processes it, and writes the result back into the space.
				    When all the workers have written their results, the master takes the results and combines them into a meaningful overall solution.
				    %Note that the workers accept new tasks whenever they are available and able to work.
				    %Therefore, while a worker is busy computing a big task, another one can solve many small tasks.
				    %In other words, this pattern naturally \emph{balances the load} on the space.
				    %Besides, it is \emph{scalable} since we can add more workers running in more machines without rewriting our code.
				    This pattern is scalable and naturally balances the load on the space.
  \item[Command pattern.] This pattern encapsulates the behavior into the tasks shared in the space.
			  Therefore, it requires (1) to share behavior through the space, and (2) any generic worker to be able to compute the behavior.
			  For instance, a graph in \ac{tsc} could include code in an interpreted programming language.
			  % nosotros pasamos de este jaleo
			  % hacer referencia a on-demand-code
			  %In that way, a generic worker application can compute whatever code other processes define and share through the space.
			  %This pattern is only possible in \aclp{ts} where the behavior can be described generically (e.g. in object-oriented \acp{ts}).
			  %Furthermore, the tasks to be performed over \ac{ubicomp} environments are not generically processable by any node.
			  %In contrast, each device is responsible of managing its own actuators on behalf of the rest.
			  %Therefore, this pattern will not be considered.
  \item[Marketplace pattern.] In this pattern, producers (or sellers) and consumers (or buyers) of resources interact to find the best deal. % TODO cita TripCom?
			  %It should be noted that the resources are products or services which can be bought and sold.
			  %Therefore, it is not applicable to the environments considered. % o use cases
			  %esto está relacionado con el paper aquel de TripCom y con el de Simon de hace un año en WoT
  \item[Specialist patterns.] In contrast to the replicated-worker pattern, in this pattern, each worker is specialized.
                              Therefore, each worker performs a particular task.
                              \citeauthor{freeman_javaspaces_1999} enumerate three subtypes:
			      \begin{description}
				  \item[Blackboard pattern.]
					It associates the concept of the space to a \emph{blackboard}.
					Following this analogy, the master is associated with a teacher, tasks with \emph{problems} and the specialized workers with \emph{students}.
					The blackboard pattern starts when the \emph{teacher} writes a \emph{problem} in the \emph{blackboard}.
					The \emph{students} observe the space and write their intention to contribute to solve the problem (i.e. \emph{raise their hands}).
					The \emph{teacher} selects to the expert which will make a modification.
					After the modification, the \emph{teacher} decides if it found the solution or another \emph{student} should contribute.
				  \item [Trellis pattern.]
					In this pattern, the master arranges the problem into low-level, mid-level, and high-level pieces.
					The workers of each level benefit from the refined data provided by the immediate level below.
				  \item [Collaborative patterns.]
					It encompasses all the patterns which allow nodes to collaborate to complete a greater task (i.e. by creating a workflow).
					%They are specific to the domain so they cannot be generalized, but it includes all the patterns where the nodes collaborate to complete a greater task
					%(e.g. by creating a workflow).
			      \end{description}
\end{description}


% Voy preparando el terreno para que se vea que los primeros son más de computación y el otro de negociación
To summarize, the \emph{replicated-worker pattern} is centered in optimizing the computation by parallelizing tasks.
The the \emph{command pattern} can be seen as an abstraction of the latter where the behavior is shipped in the tuples.
The \emph{marketplace pattern} allows negotiation of two entities through the space.
% Coño! pero si el especialista te permite colaborar...
Finally, the \emph{specialist pattern} allows nodes with distinct capacities to cooperate towards a common goal.


% Posible TODO Un esquema dividido en 4 en el que aparecen todos los patrones explicados.
% No es fácil mostrar el 2do y el 3ero.
% Una intentona de estos y otros patrones inventados en Lancaster:
%     https://docs.google.com/a/deusto.es/document/d/1QXGsQ4-wAByc_ByQ-Tts-TSphSQAzfO_c3mcgoZdloc/edit


% IG: ves... esta seccion parece q sale de la nada y creo que es el objetivo principal del capitulo, no? a ver si puedes orientar la intro del capitulo para mostrar esto mas... basicament que se vea que esto es de implementar movidas.
\subsection{Envisaged Scenarios} % IG: he aprendido palabro nuevo...
\label{sec:envisaged_scenarios}
% 1era iteración en reescritura, se puede mejorar mucho

This section devises two stereotypical scenarios for home automation.
Both scenarios emphasize how devices can coordinate in a decoupled mode using the patterns seen in the previous section. % más de uno?
%The scenarios are formed by mobile phones, and embedded platforms with sensors and/or actuators.
Specifically, we believe that the \emph{specialist patterns} fit the best to the needs of \ac{ubicomp} environments:
\begin{itemize}
  \item The devices in \ac{ubicomp} often serve to very specific needs. % IG: comentario general de esto de las comas, que me las pones todo al reves... el sujeto y el predicado sin coma entre medio (a noser que tenga de esas cosas descriptivas, en plan "jose, que es negro, se ha quemado"). Y la coma se pone cuando es algo complementario del tipo de: "en este capitulo, ..."
        For example, let us imagine a mobile phone showing a message or an embedded device turning on the lights of a room. % IG: lo de mobile es el estandard de la tesis? lo de cell no, no?
        %Although there may exist some redundancy in the tasks which different nodes can be accomplish, the task are not nodes are not as interchangeable as in... son tareas simples, no necesitan optimizar ni lanzar muchas
  \item These tasks are usually lightweight. % TODO primera vez que se habla de "task", comentar que para TripCom eran servicios web y que crearon un sistema (WXSW?)
        They do not require large computation resources, but to achieve a concrete and simple goal. % IG: comor? frase rara
\end{itemize}
Therefore, \ac{ubicomp} generally faces a problem of collaboration between nodes with different capacities.
In this problem, computation or negotiation aspects are secondary. % se entenderá a qué me refiero con computation??


%In this pattern, a master writes a task into the space and waits for its result, which is performed by some of the workers specialized in this particular task (e.g. show a message or regulate the temperature).

% IG: hala! y esto de sopeton... lo que decia, preparame el terreno un poco
\subsubsection{Scenario 1: Temperature Regulation}

The first scenario presents a room populated with several kind of sensors such as Oracle's SunSPOTs \citeweb{sunspot},
% TODO TODO TODO referenciar a la parte donde se hayan referenciado bien los dispositivos! (e.g. see ~\ref{environment})
% Para el resto proveer al menos una URL!
Digi's XBee sensors with an IP gateway,
the sensors on a KNX domotic bus \citeweb{knx}, and a fan connected to a FoxG20 \citeweb{foxg20} embedded platform as an actuator (see Figure \ref{fig:devices_scenario}).
In addition, an Android application \citeweb{android} semantically stores the temperature preferences of the user.


\InsertFig{devices_scenario}{fig:devices_scenario}{
  The devices used in the devised scenarios.
}{
}{0.7}{}


An independent node (i.e. the master node) continuously reads from the space % using \emph{read primitive}
(1) the room's temperature, and
(2) the user's desired temperature.
Note that the master node does not care about who exactly provides the temperature information.
It just takes the first available graph from the space.
When the second one is below the first one, it generates a ``\emph{decrease temperature during a certain period}'' task which can be consumed by different independent worker nodes.
In this case, the FoxG20 periodically checks for orders it can fulfill extracting them from the space. % consumes them with a \textit{take primitive}.
% zasca, acabamos de poner en bandeja que nos rechine lo de comprobar periodicamente



\subsubsection{Scenario 2: Sedentary Lifestyle Checker}

The second scenario presents an application which helps the user to avoid a sedentary lifestyle by giving him different warnings.
We consider the recommendation that states that an adult should walk at least 10.000 steps in an ordinary day \citep{tudor2002taking}.
Taking into account the expected steps which should have been completed at each moment of the day, the application generates different priority level messages.

There are several nodes involved in this task.
Each node runs on a device. % which belongs to a user.
First, an Android phone periodically updates the number of steps covered by a user that day \citeweb{pedometer}.
%Besides, it writes her profile, more important to this problem: her age. % simplificado un poquillo
Second, there is an undefined number of devices which know how to warn the user about her unhealthy behavior.
Third, there is a node in charge of generating the warnings.


\medskip


The application follows the \emph{blackboard pattern}.
The node which generates the warnings is the \emph{teacher}.
The \emph{problem} is the warning for the user.
The \emph{students} are the devices which can warn the user. % IG: no es nada intuitivo. Leyendo despues lo entiendo, pero... es raro


When the \emph{teacher} writes a warning for a user, the devices able to warn the user write their characteristics into the space. % IG: alguna cosa mejor que characteristics?
In this example, these characteristics are the device ownership and the intrusiveness of the warning method.
The \emph{teacher} reads these characteristics, chooses the most appropriate \emph{student} to solve the \emph{problem} and updates the \emph{problem} to indicate the selection.
%If no device writes their characteristics, the \emph{teacher} can wait longer.
Finally, the selected device takes any warning for it and warns the user. % no tiene respuesta
If the chosen device does not take messages, the \emph{teacher} can update the \emph{problem} with a new selection.


With this application, if the priority level is low, the user receives the warning in a less intrusive way.
For example, room's light brightness can be increased for low priority notifications, a Chumby \citeweb{chumby} can show an icon for normal priority ones and the user's mobile phone show directly the message when it is a high priority warning. % IG: yo anyadiria una tablita para los dos scenarios mostrando los actuadores y los niveles. Aqui esto me salta un poco de repente.
The remarkable characteristic of this mechanism is that the devices present in the environment at each moment can vary and the application will still fulfill its goal.




\subsection{Notification mechanism}
\label{sec:notification}


In the scenarios described in the previous section, some writings in the space trigger other node's action.
For example, a device must show a message whenever a new warning is written in the space. % e.g. a task or a result
To be aware of these writings, the node can either poll the space or rely on a notification mechanism. % async or blocking
Obviously, the later encourages a more efficient use of the network. % is more optimal and more scalable approach.
% IG: encourage no creo que sea la palabra...

As a consequence, a notification mechanism is highly advisable to fulfill \ac{ts}'s patterns in a distributed environment.
Although the implementation of this notification mechanism is out of the scope of this thesis,
% Hablar de requisitos, describir nuevas primitivas de suscripciones o alternativas de implementación
it should comply with the following aspects:
\begin{itemize}
  %  Debe permitir polling (sync), por ello proponemos nuevas primitivas: read_async y take_async.
  \item Do not substitute the polling mechanism. % porque todavía puede ser interesante en otros casos
        The middleware must provide additional \emph{read} and \emph{take} primitives.
  % Se evalúan sobre el espacio, no sobre el outer-space. 
  \item Since the notification mechanism intends to allow coordination patterns,
        it must consider the knowledge from the coordination space. % TODO comprobar que lo llamé así en el anterior capítulo
        In other words, knowledge from the \emph{outer-space} will not trigger notifications. % TODO usar constante para esto
  % Evaluarse cada X, no siempre que se escriba algo o se vuelve loco.
  \item The evaluation of the subscriptions must not interfere with the writing process (e.g. introducing a delay).
        Therefore, it must run asynchronously.
  % Debe ser sencillo de implementar por cualquier nodo. Para ello proponemos un mínimo contrato: callback URL.
  \item Ease its adoption by any type of computing platform.
	This can be achieved by reducing the requirements on the \emph{clients}.
        For instance, a callback \ac{url} passed during the subscription can represent a minimal contract between the client and the server. % no viceversa
  % Deben caducar para no mantener suscripciones de dispositivos que ya no existen.
  \item Provide additional subscription removal mechanisms.
	\ac{ubicomp} scenarios are composed by unreliable devices which may frequently join and leave the space.
	In this situation, the correct use of unsubscription primitives cannot be guaranteed.
	This may worsen the performance of the system with useless subscriptions from absent devices.
	Therefore, the device managing the subscriptions should adopt more proactive mechanisms.
	For example, it may let the subscriptions expire after a lifetime or remove them when it discovers the unavailability of a callback \ac{url}. % algo así como garbage collection
	% A) The subscriptions must expire.
        % The expiration will allow to delete subscriptions from absent devices. % no longer present devices
        % This also implies that clients are responsible for periodically updating their subscriptions.
        % B) Additional mechanism
        % For instance, a \emph{garbage collection} agent can check the availability of callback \acp{url}.
        % For unavailable callback \acp{url}, the subscriptions can be removed.
\end{itemize}


The main drawback of any subscription mechanism is that it breaks the \ac{rests} property of the \ac{rest} style.
This implies that network performance will improve at the cost of scalability, simplicity and reliability. % TODO poner un see Section~\ref{X} ???


% hasta AQUI!
% Therefore, although it is not the main focus of this thesis, in this section we present a subscription mechanism which can be used on top of the \ac{tsc} middleware proposed.
% The requirements that this mechanism needs to fulfill are the following ones:
% \begin{itemize}
%   \item It should be independent of \ac{tsc}'s writing and reading primitives.
% 	This requirement ensures that frequent writings do not lead to excessive processing in resource constrained nodes.
% 	The main drawback of this requirement is that a node does not now when relevant data is written in the space \emph{per se}.
% 	A developer needs to advertise when it writes relevant data.
%   \item Any node running our solution must be able to implement this mechanism.
% 	In other words, no new dependencies should be added.
%   \item The new primitives should use elements the current developer is used to (i.e. templates or RDF triples).
% \end{itemize}
% 
% explicar cómo funciona
% https://otsopack.readthedocs.org/en/latest/subscriptions.html
% \subsubsection{Subscription primitives}
% 
% The subscription primitives the developer should use are the following ones:
% \begin{itemize}
%   \item \emph{Subscribe}. The node subscribes to the given template returning an URI which identifies the subscription.
%     \begin{minted}{java}
% URI subscribe(space_uri, template, listener)
%     \end{minted}
%   \item \emph{Unsubscribe}. Unsubscribes to a subscription given its subscription URI.
%     \begin{minted}{java}
% void unsubscribe(space_uri, String subscriptionURI)
%     \end{minted}
%   \item \emph{Notify}.
%     \begin{minted}{java}
% void notify(space_uri, template)
%     \end{minted}
% \end{itemize}
% 
% 
% \subsubsection{Deployment}
% 
% The nodes responsible of handling subscriptions and notifications are called \emph{bulletin boards}.
% Other nodes belonging to the same space, discover them using a out of scope method and publish their subscriptions and notifications using their HTTP API.
% 
% Each \emph{bulletin board}:
% \begin{itemize}
%   \item Belongs to a space.
%   \item Exposes a subscription API.
%   \item Shares its subscriptions with other bulletin boards which belong to the same space.
%   \item Propagates the notifications to the relevant nodes using the \emph{callback url} provided by them.
% \end{itemize}
% 
% 
% The Figure~\ref{} describes a simple subscription use case.
% \begin{enumerate}
%   \item N1 subscribes to BB1 with a template t1.
%   \item BB1 propagates the subscription provided by N1 to BB2 and BB3
%   \item N3 notifies to BB3 about t2.
%   \item Since t1 matches t2, BB3 tries to notify to N1 using the callback URI provided during the subscription process.
%   \item Unfortunately, the BB3 cannot notify to N1 due to unexpected network problems.
%   \item BB3 propagates the notification of t2 to BB2.
%   \item BB2 reaches N1, so it notifies it about t2 using the callback URI.
% \end{enumerate}
% 
% TODO poner la foto esquemática de https://otsopack.readthedocs.org/en/latest/subscriptions.html

\section{Direct Consumption of \acs{rest} services}

% por qué? motivación: permitir a otras app integrarse en nuestro espacio
\ac{rest}-like services constitute an emerging mechanism to actuate in the physical environment. % REST-like, no compliant with Fielding
Therefore, their seamless integration would open the door to reuse the capacities of many existing devices. % interoperate
However, this requires applications to directly use the actuator's node \ac{http} \ac{api}.


In contrast, the actuation mechanism presented in the previous section uses the \Space{} (i.e. it is indirect).
The mechanism implies that the actuator must be aware of the content of the space.
For instance, a heater must check the space to find if a new desired temperature was written.
% lo podrá reusar a través de un intermediario
% interop era una de las cosas que queriamos cuidar
% y qué pasa con las soluciones REST que quieren usar lo nuestro?


In this section we propose to close the gap between both approaches.
To do that, we propose to transparently invoke \ac{rest}-like services on behalf of the developer.
From her perspective, she will still use the space-based patterns presented in the previous section.
However, the middleware will be enriched with the actuation capacities of third \ac{http} \acp{api}.


To define that integration, we first explain the mechanism used to describe and select the appropriate \ac{rest} services (Section~\ref{sec:background_restdesc}).
Then we sketch this integration by answering two key elements:
\begin{itemize}
  \item How do the inputs for this mechanism relate to the ones used in the \ac{ts} patterns? (Section~\ref{sec:inputs_proof})
  \item Which node is responsible for triggering this mechanism? (Section~\ref{sec:responsible_proof})
\end{itemize}



\subsection{Background}
\label{sec:background_restdesc}

In \ac{wot} physical changes in the environment are performed by manipulating \ac{http} resources.
According to the \ac{rest} principles, a client should navigate through these resources with no prior knowledge of the \ac{api}.
% Copiar esta explicación mejor de algún lado:
The client should 
\begin{enumerate*}[label=\itshape(\arabic*\upshape)]
  \item interpret the representations provided by the server and then
  \item choose the appropriate state transition from the hypertext according to its intention. % y su conocimiento básico del protocolo: CRUD
\end{enumerate*}


% es HATEOAS?
As explained in Section~\ref{sec:network_properties}, semantic representations do not include a native way to express the hypertext. % TODO realmente se explica?
To solve this, three solutions can be adopted:
% Unos proponen extender con ontologías
\begin{enumerate}
  \item To use an ontology to represent the hypertext \citep{kjernsmo_necessity_2012},
  \item To embed the hypertext independently to the representations on the \ac{http} headers \citep{mark_web_2010}, and
  \item To provide a description of the resources using the \ac{http} OPTIONS verb \citep{verborgh_functional_2012,verborgh_ijcs_2014}.
\end{enumerate}

The latter two enable to discover resources and state transitions without adding metadata to the representations.
This allows not only to describe semantic representations, but any type of formats.

\citeauthor{mark_web_2010}'s \citep{mark_web_2010} approach is extended by \citet{erik_profile_2013} to define how to embed additional semantics to process a resource representation.
These additional semantics are called profiles and are identified by an \acs{uri}.


\bigskip


\citet{verborgh_ijcs_2014} present a more expressive solution which goes beyond simply describing the resource's type.
It also allows to semantically describe the knowledge needed to use a concrete \acs{http} verb on a resource, and the content this call returns. % o precondition
The materialization of this proposal is called \restdesc{} \citep{verborgh_functional_2012}.


% TODO mirar si tiene cabida la mención de otros enfoques para describir recursos que no son muy RESTful
%Different ways exist to describe \ac{rest} services. % mencionar WADL, etc.?
% citar a donde se hable de \restdesc{} y así ya se empieza a explicar la solución de forma discreta.
\restdesc{} describes \acs{http} methods using rules expressed in the \ac{n3}\footnote{\url{http://www.w3.org/TeamSubmission/n3/}} language.
% he evitado explicar que las reglas tienen premisa y conclusión, porque me parece demasiado obvio
A rule's \emph{premise} expresses the requirements to invoke a \ac{rest} service.
A rule's \emph{conclusion} expresses both the \ac{rest} call that needs to be made and the description of that invocation result.


\citet{verborgh_ijcs_2014} propose a proof-based composition mechanism for Web \acp{api} using \restdesc{}.
This mechanism uses as inputs:
\begin{enumerate*}[label=\itshape(\arabic*\upshape)]
  \item an initial state,
  \item a goal state,
  \item Web \ac{api} descriptions using \restdesc{}, and
  \item optional background knowledge.
\end{enumerate*}
Each of these inputs are semantically expressed and therefore, they can be processed by standard \ac{n3} reasoners.
The reasoners generate proofs about how to achieve the goal starting from the initial state using the rest of the inputs.
These proofs can be seen as steps that need to be made to reach a desired state.


Additionally, \citet{verborgh_ijcs_2014} distinguish between pre-proof and post-proof.
The first, are those which assume that the execution of all \acs{api} calls will behave as expected.
The latter, can be seen as a \emph{revision} of the pre-proof.
It executes the Web \acs{api} of the pre-proofs and uses actual execution results to generate a new proof.



\subsection{Inputs for the Proof-based Actuation Mechanism}
\label{sec:inputs_proof}

The reasoning process presented by \citet{verborgh_ijcs_2014} uses as inputs: % presented by en vez de posesivo, porque "X et al.'s" queda como el culo
an initial state,
a goal state,
Web \ac{api} descriptions using \restdesc{}, and
optional background knowledge.
% inputs: clues
% goal: task


\begin{description}
  \item[The descriptions] must be read from the actuation nodes prior to reasoning. % obvio
        They can be read with a \emph{ad hoc} discovery mechanisms.
        However, in order to simplify the design of the middleware, it is sensible to use one of the already existing mechanisms:
        
        \begin{itemize}
	  \item An agent could discover the descriptions and simply write them in the coordination \Space{}. % y ahora habrá que definir que es un agente?
		Then, any node willing to reason will read them from the \Space{}.
	  \item The descriptions could be part of the \clues{} presented in Chapter~\ref{cha:searching}.
		This option ensures that they will be available in any \consumer{}.
		Besides, the static nature of these descriptions does not break the stability assumption of the \clues{}. % explicar más?
        \end{itemize}

  % TODO TODO TODO especificar que este mecanismo adicional de usar clues es optional y no el usado en nuestro escenario simple
  \item[The initial state and background knowledge] can be acquired both from the coordination space and the outer space.
	However, reading all the space would be highly inefficient.
	To reduce readings on the space (i.e. network usage), we propose a procedure composed by two reasoning steps.
	In the first one we only use local incomplete knowledge derived from the \clues{}.
	Then, we read from the \Space{} just the knowledge needed to confirm the pre-proofs obtained in the first reasoning.
	The second reasoning uses this knowledge to get real pre-proofs. % real o actual pre-proofs
	
	% Y ahora explicado en mayor detalle:
	Therefore, a node which wants to actuate over the space will need to obtain the \emph{clues} from the \ac{wp}.
	%These clues, as explained before, tell what kind of content other nodes provide. % un poco redundante
	Let us assume that these \emph{clues} are composed by the predicates used by the nodes which provide content. % es una de las alternativas que se planteaban
	The existence of a predicate used in a premise does not necessarily imply that this rule can be used.
	Nevertheless, its absence does imply that it will not be used (see Figure~\ref{fig:activation_rules}).
	Therefore, we can create temporary \emph{activation rules} from \clues{} which activate those potential rules. % latter rules las que se han mencionado primero

	\InsertFig{activation_rules}{fig:activation_rules}{
	  Sample clues, two rules and the activation rule created from the first rule.
	}{
	  According to the clues shown, the second rule will never be invoked.
	}{0.8}{}

	An \emph{activation rule} for a rule R contains a \emph{true} in the premise.
	The conclusion is made by R's premise substituting the variables with fictitious \acsp{uri} with a common prefix (see Figure~\ref{fig:activation_rules}).
	These fictitious \acsp{uri} are used to distinguish when a triple should be replaced by actual knowledge from the space. % explicar mejor?


  \item[The goal state] should be created from the task to reuse the primitives and the data from the \ac{ts}-based actuation mechanism.
	For instance, from a task of \emph{regulate temperature to 6ºC} we can deduce the goal state of \emph{temperature of 6ºC}.
	This translation is out of the scope of the dissertation.
\end{description}


\subsection{Responsibility for Triggering the Proof-based Actuation Mechanism}
\label{sec:responsible_proof}

Section~\ref{sec:background_restdesc} described two coarse-grained steps:
\begin{itemize}
  \item Reasoning over the descriptions, background knowledge, an initial state and a goal state.
        The result of the reasoning process if a pre-proof, which can be seen as a tentative \emph{execution plan} to achieve the goal.
  \item Check the execution plan by following it.
\end{itemize}


This mechanism can be performed by any node.
% no obligamos a implementar una u otra, pero recomendamos esto: XXX
In this thesis, we briefly analyze three different alternatives and their trade offs.
However, we do not adhere to any of them.
We leave as a future work to implement and quantitatively compare them.

\begin{enumerate}
  \item An agent which resides in the same machine as the \emph{coordination space} can trigger the process. % TODO buscar su nombre bueno
	Doing so, it can locally consume the knowledge available in the \emph{coordination space}.
  \item A more flexible approach would consist of letting any \consumer{} interested on changing the environment to trigger it.
	If these \consumers{} use the search mechanism presented in the Chapter~\ref{cha:searching}, they will have background knowledge about other nodes.
	This reduces the dependency on the node providing the \emph{coordination space}. % TODO buscar su nombre bueno
	However, it requires them to perform tasks such as reasoning and checking the pre-proofs.
	While the latter increases the network usage, the first increases the computation.
	As we already mentioned in previous chapters, these tasks severely affect to the energy consumption.
	Furthermore, some resource platforms will not be even able to reason.
  \item To mitigate that problem, we could delegate this task only on the nodes able to perform such tasks.
	In fact, these nodes could follow the \emph{replicated-worker pattern}.
	They could read from the space goals to trigger the process (i.e. \emph{reasoning tasks}).
	Apart from balancing the load between all the worker nodes, any node can stop being worker at any time by not taking more \emph{tasks} (e.g. if it has low energy).
	These nodes must be \consumers{} to use the \clues{} from the search mechanism as background knowledge.
\end{enumerate}


% Poner ejemplo adaptado del paper de WoT2013?


%Proceso de suggest:

%\begin{itemize}
% \item obtener todo el conocimiento necesario para el proceso?
% \item razonamiento sobre conocimiento
% \item parsear el resultado para invocar servicios HTTP
% \item monitorizar que se ha complido el cambio?
%\end{itemize}

\section{Comparison}

The first actuation mechanism requires a subscription mechanism and provides space and time autonomy.
However, it does that by creating a dependency on the \Space{}.
Two nodes will not be able to communicate with each other without the \Space{}.


The second approach offers this independence of the \Space{}: any node can directly invoke a service to act over the environment.
Remarkably, it allows the integration of third \ac{wot} applications in our middleware. % apps, APIs o services?
To do that, it requires a rule-based reasoning engine which is not available in many computing platforms. % reasoning engine ya queda bien?



% Tabla: XXX
% autonomía: poner o simplemente indirect vs direct?
% requirements in consumers: estar pendiente del espacio vs razonar
% requirements in providers: estar pendiente del espacio vs (describir semánticamente servicio => equivalente a tasks)
\begin{table}[htbp]
  \caption{Characteristics of the discussed actuation mechanisms.}
  \begin{center}
    \footnotesize
    \begin{tabular}{llll} 
      \hline
      \multirow{2}{*}{Actuation} &
      Communication & % o poner autonomias y demás para enfatizar?
      \multirow{2}{*}{Benefits} &
      Required \\
      &
      style &
      ~ &
      features \\
      \hline
      Space-based & Indirect & Decoupled communication & Subscriptions \\[0.2cm]
      REST-based & Direct & Reuse of third \ac{wot} & Rule-based \\ % mencionar interop? % TODO ampliar la noción de WoT a REST-like?
      & & applications & reasoning \\ % no digo "wot actuation capacities porque ya está verbalizado en la explicación"
      \hline
    \end{tabular}
  \end{center}
  \label{tab:actuation_mechanisms}
\end{table}


\bigskip

The impact of both mechanisms in resource constrained devices will be affected by computing and networking activities.
\ac{ts}-based patterns require all the participants to be aware of what is written in the space to react (i.e. be proactive).
Both consumers and providers read and write from the space, subscribe to specific changes and receive notifications.
They are only affected by the contents they are interested in.
Finally, they only perform trivial computing task: interpreting results and pattern matching. % pattern matching sólo en el caso del coordination space!


The proof-based actuation mechanism requires the consumer to have prior knowledge about the environment to reason over it. % frase susceptible de simplificar
Since this knowledge must be acquired from remote nodes somehow, this approach might demand more network usage than the first one.
% de pre-proof a post-proof se generará más networking activity seguro? por qué? por invocar servicios REST? Es equivalente a andar leyendo tareas y escribiendo resultados...
Besides, the proof-based actuation itself also generates more computation activity on the node responsible of reasoning.
Both aspects have a negative impact in the energy consumption.
Consequently, it will most likely not be an adequate mechanism for resource constrained devices.
As a contrast, this mechanism requires few things from the provider perspective: to serve \ac{rest}-like services and provide their descriptions. % aclarar quien es el provider/actuator?




% Tabla: XXX
% autonomía: poner o simplemente indirect vs direct?
% requirements in consumers: estar pendiente del espacio vs razonar
% requirements in providers: estar pendiente del espacio vs (describir semánticamente servicio => equivalente a tasks)
\begin{table}[htbp]
  \caption{Foreseeable networking and computing impact on the nodes involved in the actuation mechanism.}
  \begin{center}
    \footnotesize
    \begin{tabular}{llp{4cm}p{4.4cm}}
      \hline
      \multirow{2}{*}{Actuation} &
      \multirow{2}{*}{Perspective} &
      \multicolumn{2}{c}{Activity} \\
      ~ &
      ~ &
      \multicolumn{1}{c}{Networking} &
      \multicolumn{1}{c}{Computation} \\
      \hline
      \multirow{3}{*}{\ac{ts} patterns} & Provider & Proactive, limited activity & Limited: Results Interpretation \\
				    ~ & Consumer & Proactive, limited activity & Limited: Results Interpretation \\
				    ~ & Space & Reactive, high activity & Varies with the implementation \\[0.2cm]
      \ac{rest} \acp{api} & Provider & Reactive, limited activity & Limited: Handling requests \\
        consumption       & Consumer & Proactive, high activity &  Demanding: Reasoning \\ % Y comprobar planes...
      \hline
    \end{tabular}
  \end{center}
  \label{tab:actuation_networking_computing}
\end{table} % o si queda muy corto, incluirlo directamente en conclusions
\section{Conclusion}

In this chapter we presented two ways to actuate on the physical environment.
The first is the usual way to operate through spaces and provides a higher degree of decoupling.
However, it requires participants to use our middleware's primitives.
In other words, our middleware is not able to reuse third applications \ac{rest} services.



Since the interoperability is one of the guiding principles of our middleware, we presented an additional actuation mechanism.
This mechanism relies in the semantic description of the services, additional knowlege and in a reasoning process. % additional knowledge: background + initial
With that information, it is able to generate executions plans towards a goal.
Following these plans implies different calls to the different services.


We sketched an allignment of our middleware with this actuation mechanism.
The tasks written in the space can be translated into goals.
The knowledge from the space can be used as an input for the reasoning process.
Different nodes are able to trigger the reasoning and check the plan.
Remarkably, we suggest the use of \ac{ts} patterns to promote load balancing between the devices able to perform that process.


% Future work o conclusion?
The second actuation alternative requires further work to check some of the assumptions made.
For example, will always be possible to translate a task into a goal.
Additionally, it opens the door to solve interesting questions: when there are two or more paths to a goal, how to discern which one to follow.


Furthermore, the middleware must deal with the coexistence of both mechanisms.
Will both methods be triggered indistinctly or will the first prevail over the second?
% uso de nuestra actuación por parte de apps WoT
Finally, the question of how to reuse actuation mechanisms of nodes using \ac{ts} patterns from \ac{wot} has not been even addressed.
However, this chapter has exposed, described and compared the key points towards the actuation through a \ac{tsc} middleware for \ac{ubicomp}.

%: ----------------------- conclusion ------------------------
% ----------------------------------------------------------------------

\begin{savequote}[50mm]
Now this is not the end.
It is not even the beginning of the end.
But it is, perhaps, the end of the beginning.
\qauthor{Winston Churchill}
\end{savequote}


\chapter{Conclusions}
\label{cha:conclusions}
\newcommand{\pathchapsix}{6_conclusion}


% the code below specifies where the figures are stored
\ifpdf
    \graphicspath{{\pathchapsix/figures/PNG/}{\pathchapsix/figures/PDF/}{\pathchapsix/figures/}}
\else
    \graphicspath{{\pathchapsix/figures/EPS/}{\pathchapsix/figures/}}
\fi


%------------------------------------------------------------------------- 

% Presentar de qué ha ido la tesis (TODO añadir capítulos en los que se ha ido tratando?)
% 3 ptos con los que quiero que el lector se quede: searching, acting y compatibility with WoT

This dissertation explores the design of a distributed \ac{tsc} middleware suitable for \ac{ubicomp} environments.
The middleware's most remarkable characteristic is its alignment with \ac{rest}-like services and the \ac{wot} world.
This alignment pursues the interoperability of the middleware with other semantic \ac{wot} solutions.
This way, it encourages the reuse of their data and eases the access to its own data.
% Otra forma de vender explicitamente la importancia de integrarlo con WoT:
% Besides, after an analysis we acknowledged the importance and benefits that the upcoming \ac{wot} has brought to the \acl{ubicomp}.
% This has lead us to have make our final solution compatible with the \ac{wot} and other REST services.


To preserve as many beneficial properties as possible from \ac{wot} and \ac{ts} worlds, the middleware is divided in two \Spaces{}: the coordination space and the outer space. % TODO usar constantes
The first one enables a completely indirect communication, time and space decoupled, between the nodes by writing and reading semantic content in it.
The latter acknowledges the distributed nature of \ac{ubicomp} environments by letting the nodes manage their own semantic content.
However, this content is shared with the rest of nodes forming a \emph{virtual} space.
Developers access to the content of the space in a associative manner keeping the space decoupling property.
Both spaces are accessed through a resource-oriented \ac{http} \ac{api} to ease reusing contents by third \ac{rest}-like applications. % TODO se entiende REST-like?? habrá que explicarlo


The distribution of the information avoids participants' dependency on other devices, but arises other challenges.
This dissertation tackles two main challenges: how to search in a decentralized way and how to act on the environment through the space.


To enhance the search process between devices limiting their dependencies in others we propose an architecture.
This architecture dynamically chooses an intermediary according to its capabilities to support the most limited devices.
The final result is that less requests are required in overall, specifically reducing the load on resource constrained devices.
This results in the reduction of their energy consumption.
% es independiente de nuestra solucion: se puede aplicar a otras cosas


To actuate on the physical environment we first presented how to apply \ac{ts} application patterns to the \ac{ubicomp} vision.
%Specifically, tasks and results for these tasks are read and written into the space.
% Necesidad de suscripciones
Although these actuation mechanisms promote indirect communication between the nodes, they require them to use the coordination space.
This impedes to seamlessly integrate actuators from the \ac{wot} world. % without a gateway.
To avoid that, i.e. to promote the reuse of third \ac{wot} applications' actuation capacities, we presented how to integrate a proof-based mechanism with the regular \ac{tsc} usage.
%Finally, the advantages and drawbacks of both approaches are presented.


\section{Contributions}

% De una charla con Buján recuerdo que dijo que para O. Corcho...
%   + Contribuciones científicas: lo que nadie había hecho antes (innovación)
%   + Contribuciones técnicas: implementar o aplicar en otro campo un aspecto innovado por otro


This section summarizes the contributions of this dissertation.
\begin{itemize}
  \item Chapter~\ref{cha:stateoftheart} presented an in-depth state-of-the-art review.
    \begin{itemize}
      %\item It presents the research fields to which the dissertation contributes. % Background
      \item It portrayed this dissertation regarding the relevant works from related research fields.
	    Particularly, it intensified the analysis on the semantic space-based computing middleware. % deepen, intensify
	    This analysis showed that other works do not address the problems tackled in this dissertation.
    \end{itemize}
  
  \item Chapter~\ref{cha:tsc} depicted a two-space model which
	\begin{enumerate*}[label=\itshape(\arabic*\upshape)]
	  \item preserves \ac{tsc}'s decoupling properties and
	  \item respects \ac{ubicomp}'s distributed nature.
	\end{enumerate*}
    % designs a middleware
    \begin{itemize}
      %\item It scrutinizes the \ac{rest} style and its compatibility with \ac{tsc}.
      \item It described a dual space model. % compromise between
            The first space is a common place where all the participants can write and read.
            The second space is a \emph{read-only} virtual \Space{} formed by contents managed by the different participants.
            It is accessed by any node as a whole using the same primitives as the first one.
      \item It designed the middleware by presenting the adopted primitives and how the spaces can be accessed through an \ac{http} \ac{api}.
	    %Doing so, we retain some of \ac{rest}'s properties and enable interoperability with other elements.
      \item It analysed the properties of the designed middleware from different perspectives.
	    Particularly, it scrutinized which beneficial properties from \ac{tsc} or \ac{rest} are retained.
	    Besides, it remarked the middleware's suitability for resource constrained devices.
    \end{itemize}
  
  \item Chapter~\ref{cha:searching} presented a search architecture for semantic \ac{wot} solutions.
        This architecture is compatible with, but it is not only applicable to, the space model proposed.
    \begin{itemize}
      \item It presented a search architecture that:
      \begin{itemize}
	\item Balances the load that the architecture management generates.
	      This balance considers devices' computing and energy capacities.
	\item Structures the nodes in different dynamic roles with different responsibilities. % roles: WP, Provider y Consumer, dynamic: ppalmente WP
	\item Promotes end-to-end \ac{http} requests between constrained devices.
	      They do not need an intermediary to search (i.e. to select the appropriate providers to request). % sé que esto es tricky, pero no lo necesita necesariamente
      \end{itemize}
      \item It assessed different types of information summaries the nodes can use to improve their search.
      \item It evaluated the proposal with a simulated but yet realistic environment.
    \end{itemize}
    
  \item Chapter~\ref{cha:actuate} analysed how to actuate on the physical environment using a space. % uno direct y otro indirect
	Particularly, it showed how to interoperate with existing \ac{rest}-based actuators from space-based computing.
    \begin{itemize}
      \item It explained how to use \ac{tsc} to coordinate actuator nodes and nodes willing to actuate.
      \item It depicted the need of a subscription system and its key requirements.
      \item It compared the space-based actuation technique with a direct actuation one based on the \ac{rest} architectural style. % the requirements and benefits
      \item It contributed with a hybrid approach which seamlessly reuses \ac{rest} \acp{api} from our space model.
    \end{itemize}
  
  % contribuciones técnicas: implementaciones varias
  \item Besides, as a result of the theoretical work made on this dissertation, we have done the following technical contributions:
    \begin{itemize}
      \item \emph{Otsopack} \citeweb{otsopack}: a \ac{tsc} middleware which works over \ac{http} and implements most of the ideas presented in the dissertation.
            This middleware is publicly available for different computing platforms and has been used in several research projects.
      % Esto deberiamos incluirlo aunque sea brevemente en algun lugar...
      % \item Exploration of the use of semantic content in several embedded platforms.
      %  We have tested several platforms, semantic frameworks and reasoners to verify the feasibility of the ideas and assumptions presented in the dissertation.
      \item A fully parametrizable simulation framework to evaluate different communication strategies \citeweb{simulation}.
      \item Three different implementations of the same simple \ac{ubicomp} scenario which illustrate the ideas explained in Chapter~\ref{cha:actuate}.
	    Much of these implementations are fully reusable for further future evaluations \citeweb{actuation}.
    \end{itemize}
\end{itemize}


% el tema de haber trabajado con otros
Finally, parallel to this dissertation, but deeply influenced by it, the author of this dissertation has collaborated with other academic institutions, research centres and companies to make further scientific and technical contributions.
\begin{itemize}
  \item Lightweight user access control for resource constrained devices. % que puede ser adaptada a esto
  \item Apply the \emph{Otsopack} middleware in several scenarios:
    \begin{itemize}
      \item Homes: automating them according to the user's preferences. % y lo del podometro?
                                                                        % en los escenarios de distintos papers
                                                                        % lo que se hizo en TALIS+Engine en Alicante o no se donde
      \item Hospitals and residences: tracking the evolution of patients with cognitive impairments.  % TODO traducir centro de día mejor que hospital
      \item Supermarkets: aiding to buy and easing mobility using robots.
      \item Hotels: intelligently adapting them to the customers' needs.
    \end{itemize}
  % Poner CLIPS4Android???
\end{itemize}
\section{Relevant Publications}

% International JCR Journals
% International conferences
% Workshops % o si queda muy corto, incluirlo directamente en conclusions
\section{Future Work}

% Limitations => Future Work
% ¿Poner lista de la compra o sólo las cosas que se destilen de las limitaciones del trabajo?

Previous chapters of this dissertation have pointed out several limitations and open issues.
These limitations and issues open the door to interesting future work which is detailed below.


\subsection{Closing the Gap with Third \acs{wot} Applications}

One of the middleware's goals is to achieve interoperability with third applications.
In other words, to reuse their data transparently to the user from our middleware and vice-versa.
This goal can be seen as a two-fold strategy: ease the access to the data and make them reusable.

\begin{description}
  \item[Access to the data.] It must be done by means of an \ac{api}.
			    In this thesis we have defended the need of using \ac{http} for that purpose.
			    An adaptation to the upcoming \ac{coap} would ease its adoption by new limited platforms. % "base de usuarios"
			    However, the most challenging aspect is not to replicate it in other suitable protocols.
			    
			    The most challenging aspect is to automatize the consumption of heterogeneous third \acp{api}.
			    According to the \ac{rest} principles, the \ac{resthateoas} constraint can help automated the clients to use any \ac{api} independently of its specific shape. % shape or implementation
			    Instead, this dissertation proposes a common \ac{api} that needs to be implemented in any device willing to share content. % TODO data api vs domain api => mirar el capitulo 3 y explicar esto
			    
			    In any case, an agent would need knowledge about the domain explored to reuse any \ac{api}.
			    And to reuse many \acp{api} without reimplementing the agent, the data types and states of a domain should be standardized. % TODO poner cita del blog de fielding sobre la discussion esa
			    Therefore, any step in that direction would require the agreement of the community around the domain (e.g. the \ac{wot} community).
			    
			    The counter-pad of third \acp{api} automatic consumption would be to ensure that our \ac{tsc} \ac{api} complies with the \ac{resthateoas} constraint.
			    Some steps have been done to make the \ac{api} human discoverable through \ac{html} representations.
			    However, more effort needs to be done to ease its automatic consumption regardless of its specific shape. % shape or implementation? % which would depend on the version
			    % se pueden usar varias estrategias en paralelo
			    
			    Finally, the \ac{resthateoas} constraint may lead to more network usage, which is particularly harmful for devices with energy autonomy limitations.
			    Therefore, this trade-off between the interoperability and the network efficiency should be also assessed.
			    
  \item[Reusing the data.] The semantic knowledge contributes to that goal.
			  However, as \citet{barnaghi_semantics_2012} point out, the use of semantic descriptions alone does not ensure interoperability.
			  Additionally, they suggest the need of:
			  (1) standardizing ontologies and using them (or at least when using own ontologies, referencing an upper ontology); and
			  (2) interpreting the annotations effectively. % efectivamente: con razonadores que funcionen bien
			  While the first one requires the consensus of the community, the second demands for efficient reasoners for mobile and embedded platforms.
			  This efficiency is particularly important since it directly influences their energy consumption.
			  % TODO Mencionar que estamos trabajando en ello?
			  % further assessment of existing ones, creating news for platforms
			  % pero decir que lo nuestro es un primer paso en paralelo a otros que puedan solucionar los otros problemas
			  % semantica y cacharros pequeños, sigue siendo algo difícil. Necesidad de razonadores ligeros.
			  
\end{description}


\subsection{Increasing Search Expressibility and Efficiency}

% Searching: SPARQL y demás
% 1ero: mayor eficiencia
% reescritura puede mejorar interop
We did not consider using query languages such as SPARQL because of the difficulties most constrained platforms may find in parsing them.
With this decision, we design a uniform \ac{api} for all the types of devices.
However, for the future work, one can adopt SPARQL thanks to an optional \ac{api}.
In that case, all the nodes should be aware that some of them are only able to interpret triple pattern templates.

The use of SPARQL in some nodes would benefit the developers and improve the middleware performance. % searching
The users of the middleware would be able to made more expressive queries.
The performance could be enhanced by introducing optimization techniques from the distributed databases field \cite{schwarte_fedx_optimization_2011}. % see references from Section~\ref{background}) or reference them here?


\subsection{Coordination Space's Distribution, Replication or Migration} % no se me ocurre un título más sexy

The current design demands the coordination space to be accessible through a \ac{http} \ac{api}.
For the sake of clarity, at some points of the dissertation we have assumed that this space is centralized in a unique machine.
However, the access through a \ac{http} \ac{api} does not restrict that.
Some alternatives that can be worth to explore are: % mencionar cloud para algo?
\begin{itemize}
  \item Replication of the space to increase the system performance (e.g. load balancing).
  \item Distribution of the content over different machines to enhance its scalability. % TODO decir que ya ha sido tratado por más soluciones? (e.g. TripCom)
  \item Migration from a machine to another to avoid a single point of failure.
	%No node is essential so the system will still work.
        This strategy will allow, for example, to select the space holder dynamically easing the middleware manageability and deployment.
        This would be similar to the \ac{wp} role of our searching architecture. % mencionar Chapter~\ref{cha:searching}?
\end{itemize}

% aquella idea que tenían los de Lancaster de que el middleware cambia de forma dependiendo del entorno
Besides, monitoring the middleware usage, one alternative or another could be dynamically chosen depending on the specific needs of each situation.


\subsection{Integration of the Proof-based Actuation Mechanism in the Space}

Chapter~\ref{cha:actuate} presented an alignment between a proof-based mechanism and actuation through the space.
However, its implementation and evaluation was left as future work.
To enable the usage of the proof-based mechanism, tasks from the latter must be translated into goals.
This translation must be generalizable to automatize it.
However, it is not clear whether it will always be possible this direct translation for complex goals and tasks.

% Idea del fallido WoT2013:
% Let us imagine that at the end of the previous section more than a path is found for the same desired goal.
% For instance, the light of a room can be increased to a certain level by turning on a lamp or by raising the blinds.
% In that case, different behaviors could be applied to select the preferred path.
% For example, one will want to save energy in the long term by raising the blinds.
% On the other hand, another application will decide to reach the light level no matter of the outside conditions.
% 
% Furthermore, these behaviors might be reused among the applications. % además, seguimos con la idea de reutilizar
% For that, we envision to share the behaviors in the space itself. % se infiere que semánticamente anotados
% The behaviors could be defined by the facility administrators or the developers.
% After that, the developers will be able to consume the behaviors through another middleware abstraction.
% Besides, it will monitor the state to check if the state of the space is effectively changed or other actions need to be taken.
Another interesting issue is to discern which path to follow to achieve a goal when two or more paths are available.
We believe that additional high-level rules or policies can be defined (e.g. to select the path which consumed less energy).
Storing them into the space may enable reusing this \emph{behaviors} by third applications' developers.

Finally, apart from reusing third \ac{wot} applications' actuation capacities in \ac{tsc}, the opposite question could also be explored:
How to reuse actuation mechanisms of nodes using \ac{ts} patterns from \ac{wot} solutions?


\subsection{Security for the Middleware}

% si no, es un poco fiesta y cualquiera podría poner datos falsos
Two key aspect for \ac{ubicomp} deployments are
(1) to authenticate the parts involved in the communication; and
(2) to ensure the privacy of the message being transmitted.
% Ya hemos hecho algo para dispositivos limitados y móviles
Although both problems have been deeply discussed in the literature,
the nodes in an \ac{ubicomp} environment face severe computing and energy restrictions.
To solve that, we have already propose a solution appropriate for sensor networks and mobile devices.

The next step is to evaluate it in more demanding scenarios and the final goal is to integrate it in our \ac{tsc} middleware.



%%% A partir de aquí, no son aspectos técnicos concretos


\subsection{Measuring Ease of Usage}

A subtle goal of any middleware is to ease the development of applications by encapsulating complex tasks (i.e. searching on a distributed space).
% Medir de forma objetiva aspectos del ser humano como la sencillez de uso o de implementación
Some objective indicators can be extracted from the code by implementing the same application with and without the middleware (i.e. number of lines).
However, from the human perspective these indicators do not necessarily express how easy, comfortable or pleasant was to use the middleware.
% facilitar modelado?? => el hecho de que sea difícil anotar, puede echar para atrás a desarrolladores: eso es problema de mi app o de quien?
Although the primitives are rather simple, the configuration or just the semantic annotation can complicate the learning process of any developer.
Future work could measure the impact the middleware has on them.
As a result, some improvements could be proposed (e.g. to create a version for a specific domain which masks all the semantic issues).


% Al final, a modo de colofón
\subsection{Further Evaluation on Real Deployments}

% muchas de nuestras decisiones en la implementación del middleware fueron guiadas por la realidad
Rather than designing the best theoretical solution, in this dissertation we have tried to present a realizable solution.
For that, we have implemented most of the ideas in the \emph{Otsopack} middleware and we have promote its use in different research projects.
This has helped us to identify the key problems to solve.

Besides, to efficiently evaluate our solution under very different and extreme circumstances, in some occasions we have simulated the scenarios. % con muchos cacharros
In this simulations we have tried to present always realistic assumptions.
% de simulaciones => a la realidad
However, there is always a gap between what is perceived as realistic and the reality.

To bridge both worlds, the next natural step would be assessing some aspects of our middleware in a real world environment: % real world es académico???
\begin{itemize}
	% Big ecosystem of interconnected objects and devices
  \item The search architecture.
        For that, we should deploy an scenario with a number of heterogeneous devices running a task in a long term.
        % viendo cuales son las necesidades reales de quienes lo usan
	%    para justificar los parametros que se usaron
        Then, we could measure their network usage and energy consumption to either tune up some parameters of the architecture or redesign some aspects.
  % relacionado con el punto de interoperabilidad: reusar distintas applicaciones existentes y ver que realmente se cumple
  %    una forma de asegurarse, es haciendo que distintos usuarios las creen
  \item The interoperability.
	New unconsidered problems could arise trying to reuse data from existing \ac{wot} applications.
	Each manual intervention needed to make our middleware work together with these applications would be a sign of an interoperability problem.
\end{itemize}
\section{Final Remarks}

% Cura de humildad diciendo lo que querrias haber conseguido
% (possitive) leasons learnt:
% 1. middleware fácilmente implementable en distintas plataformas
% 2. la semántica al principio es complicada, pero muestra sus beneficios en el largo plazo facilitando la integración de apps


% ----------------------------------------------------------------------	


% --------------------------------------------------------------
%:                  BACK MATTER: appendices, refs,..
% --------------------------------------------------------------

% the back matter: appendix and references close the thesis
\backmatter


%: ----------------------- appendix ------------------------

%\appendix


%: ----------------------- bibliography ------------------------



% The section below defines how references are listed and formatted
% The default below is 2 columns, small font, complete author names.
% Entries are also linked back to the page number in the text and to external URL if provided in the BibTex file.


% Original version:

% PhDbiblio-url2 = names small caps, title bold & hyperlinked, link to page 
%\begin{multicols}{2} % \begin{multicols}{ # columns}[ header text][ space]
%\begin{tiny} % tiny(5) < scriptsize(7) < footnotesize(8) < small (9)
%
%\bibliographystyle{Latex/Classes/PhDbiblio-url2} % Title is link if provided
%\renewcommand{\bibname}{References} % changes the header; default: Bibliography
%
%\bibliography{9_backmatter/references} % adjust this to fit your BibTex file
%
%\end{tiny}
%\end{multicols}



% Show all bibliography entries
%\nocite*



% If we want bibliography backreference, use unsrt first and the desidered one after

%\bibliographystyle{unsrt} % Defines the bibliography style

%\bibliographystyle{alpha} % Defines the bibliography style

%\bibliographystyle{apa-good} % Defines the bibliography style
%\bibliographystyle{natbib} % Defines the bibliography style

%\bibliographystyle{plainurl}

%\renewcommand{\bibname}{References} % changes the header; default: Bibliography

%To include the references/works cited/bibliography in your Table of Contents, right before the bibliography command, use the command
%\addcontentsline{toc}{section}{References}
%

\bibliography{9_backmatter/ref_rest,9_backmatter/ref_ubicomp,9_backmatter/ref_space,9_backmatter/ref_semantic,9_backmatter/ref_self,9_backmatter/ref_others} % adjust this to fit your BibTex file

% --------------------------------------------------------------
% Various bibliography styles exit. Replace above style as desired.

% in-text refs: (1) (1; 2)
% ref list: alphabetical; author(s) in small caps; initials last name; page(s)
%\bibliographystyle{Latex/Classes/PhDbiblio-case} % title forced lower case
%\bibliographystyle{Latex/Classes/PhDbiblio-bold} % title as in bibtex but bold
%\bibliographystyle{Latex/Classes/PhDbiblio-url} % bold + www link if provided

%\bibliographystyle{Latex/Classes/jmb} % calls style file jmb.bst
% in-text refs: author (year) without brackets
% ref list: alphabetical; author(s) in normal font; last name, initials; page(s)

% Other styles such as alpha don't go well with natbib, and I want to use \citep, \citeauthor, etc.!
\bibliographystyle{plainnat} % calls style file plainnat.bst
% in-text refs: author (year) without brackets
% (this works with package natbib)


% --------------------------------------------------------------


% HOWTO Citar páginas concretas o comentarios (blogs):
% http://tex.stackexchange.com/questions/89143/citation-with-some-text-inside-square-brackets

%: Declaration of originality

% this file is called up by thesis.tex
% content in this file will be fed into the main document

%: ----------------------- introduction file header -----------------------


% the code below specifies where the figures are stored
\ifpdf
    \graphicspath{{9_backmatter/figures/PNG/}{9_backmatter/figures/PDF/}{9_backmatter/figures/}}
\else
    \graphicspath{{9_backmatter/figures/EPS/}{9_backmatter/figures/}}
\fi

% Thesis statement of originality -------------------------------------

% Depending on the regulations of your faculty you may need a declaration like the one below. This specific one is from the medical faculty of the university of Dresden.

\begin{declaration}        %this creates the heading for the declaration page

I herewith declare that I have produced this work without the prohibited assistance of third parties and without making use of aids other than those specified; notions taken over directly or indirectly from other sources have been identified as such. This work has not previously been presented in identical or similar form to any examination board.

The dissertation work was conducted from 2010 to \the\year \ under the supervision of
\Advisor 
~and~
\Advisortwo
~at 
the University of Deusto.

\vspace{10mm}

% City
Bilbao,

% Signature figure

%\begin{flushright}
%\begin{figure}[htbp!]
%\end{figure}
%\includegraphics{signature}%
%\end{flushright}


\end{declaration}



%-------------------------------------------------------------------------

\thispagestyle{empty}

\hfill
\vfill
\medskip

\begin{center}
\noindent
This dissertation was finished writing 
in 
Bilbao
on
% Date
\today 
\end{center}

\vfill
\medskip
\vspace{1cm}
\bigskip


%-------------------------------------------------------------------------

% Blank page
\clearpage
\thispagestyle{empty}

\hfill
\vfill
\medskip

\begin{center}
\textit{This page is intentionally left blank}
\end{center}

\vfill
\medskip
\vspace{1cm}
\bigskip


% ----------------------------------------------------------------------





\end{document}
